\chapter{Audited dependency plan}

The plan lists the intermediate mathematical obligations, existing upstream results, and canonical Lean interfaces. A linked statement may still contain an explicit proof obligation. The full graph includes companion results; the rooted graph isolates the intended dependencies of Preiss.


\section{Definitions and Euclidean geometry}

\begin{definition}[Linear planes]
\label{def:plane}
\lean{Preiss.Euclidean, Preiss.Plane}





An $m$-plane in $\R^n$ is an $m$-dimensional real linear subspace through zero.
Keep its definition independent of tangent measures and of a topology on planes.
\end{definition}

\begin{definition}[Countable Lipschitz coverage]
\label{def:coverage}
\lean{Preiss.CountablyLipschitzCovered}





A set $S$ is countably Lipschitz covered in dimension $m$ if a countable
family of Lipschitz maps $f:\R^m\to X$ covers $S$. The family may be empty.
No measurability or finite-measure hypothesis is part of this predicate.
\end{definition}

\begin{definition}[Lipschitz pieces and countable partial-domain coverage]
\label{def:parametrized-coverage}
\lean{Preiss.CountablyCoveredByLipschitzPieces, Preiss.LipschitzPiece, Preiss.LipschitzPiece.range}





The target is a metric space.
A Lipschitz piece consists of an arbitrary subset $A\subset\R^m$ and a Lipschitz map $f:A\to X$. A set is countably covered by such pieces if it lies in the union of the ranges of a countable family. The domains and the family may be empty. No measurability or finite-measure hypothesis is built into the geometric piece.
\par\emph{Role: Preiss proof.}
\par Lean declarations: \texttt{Preiss.CountablyCoveredByLipschitzPieces, Preiss.LipschitzPiece, Preiss.LipschitzPiece.range}.
\end{definition}

\begin{lemma}[Euclidean extension of Lipschitz parametrizations]
\label{lem:domain-extension}
\lean{Preiss.exists_lipschitz_extension_euclidean, Preiss.exists_lipschitz_extension_finiteDimensional}



The subset \(A\) is arbitrary, possibly empty. The conclusion requires some finite Lipschitz constant, not necessarily the original one.
Every Lipschitz map \(f:A\to\mathbb R^n\), with \(A\subset\mathbb R^m\), admits a globally defined Lipschitz extension \(\widetilde f:\mathbb R^m\to\mathbb R^n\).
\par\emph{Role: Preiss proof.}
\par Lean declarations: \texttt{Preiss.exists\_lipschitz\_extension\_euclidean}.
\end{lemma}
\begin{proof}[Planned proof obligations]
Reuse extend\_pi through finite-dimensional coordinates; no optimal Kirszbraun bound needed.
\par Source: De Lellis Chapters 4--5; see the rectifiability audit.
\end{proof}

\begin{definition}[Unit uniform measures]
\label{def:unit-uniform}
\lean{Preiss.IsUnitUniform}





Write $\nu\in\mathcal U_m$ when $\nu$ is regular, $0\in\operatorname{supp}\nu$, and $\nu(\overline B(x,r))=r^m$ for every support point $x$ and every $r>0$. This definition permits $m=0$. Its adapter to FormalizingGMT's positive-dimensional uniform predicate assumes $m>0$.
\end{definition}

\begin{definition}[Lipschitz graphs]
\label{def:lipschitz-graph}
\lean{Preiss.IsLipschitzGraph}
\uses{def:plane}










A full $m$-dimensional Lipschitz graph is
$\Gamma=\{v+f(v):v\in V\}$, where $V$ is an $m$-plane and
$f:V\to V^\perp$ is Lipschitz. The function need not vanish at zero, so
this includes affine translates. Containment in a graph and coverage modulo
$\Hm$-null sets are separate assertions. See De Lellis Definition 4.1.
\end{definition}

\begin{lemma}[Compact projection model of fixed-dimensional planes]
\label{lem:grassmannian}
\lean{Preiss.isCompact_rankProjectionSet, Preiss.planeProjection, Preiss.planeProjection_injective, Preiss.proj, Preiss.range_planeProjection, Preiss.rankProjectionSet}
\uses{def:plane}








Use self-adjoint idempotent operators of trace $m$ to model $m$-planes.
This is a nonempty compact set when $m\leq n$.
Its distance functions converge uniformly on bounded sets when the projections
converge. No lattice topology on the submodule type is intended.
\par\emph{Role: Preiss proof.}
\par Lean declarations: \texttt{Preiss.isCompact\_rankProjectionSet, Preiss.planeProjection, Preiss.planeProjection\_injective, Preiss.proj, Preiss.range\_planeProjection, Preiss.rankProjectionSet}.
\end{lemma}
\begin{proof}[Planned proof obligations]
Use self-adjoint idempotent operators of trace $m$ to model $m$-planes.
This is a nonempty compact set when $m\leq n$.
Its distance functions converge uniformly on bounded sets when the projections
converge. No lattice topology on the submodule type is intended.
\par Source: De Lellis Lemma 6.20 and Section 6.4, pp.66-69.
\end{proof}

\begin{definition}[Cones about affine planes]
\label{def:cone}
\lean{Preiss.closedTube, Preiss.cone}
\uses{def:plane}








Work in a finite-dimensional real inner-product space; \(V\) is a linear subspace.
For a linear plane \(V\), put \(Q_V=\mathrm{id}-P_V\) and define
\[
 C(x,V,\alpha)=\{y:\|Q_V(y-x)\|\leq\alpha\|P_V(y-x)\|\}.
\]
Geometric applications use \(\alpha\geq0\).
\par\emph{Role: Preiss proof.}
\par Lean declarations: \texttt{Preiss.closedTube, Preiss.cone}.
\end{definition}

\begin{lemma}[Hausdorff measure on planes]
\label{lem:plane-hausdorff}
\lean{Preiss.flatOn, Preiss.map_hausdorffMeasure_subtype_eq_restrict, Preiss.planeMeasure_ball_and_sphere, Preiss.unitFlat, Preiss.unitFlat_closedBall}
\uses{def:plane}








The linear subspace \(V\subset\mathbb R^n\) has dimension \(m\). Intrinsic and ambient Hausdorff measures use the same diameter convention.
For the inclusion \(\iota_V:V\hookrightarrow\mathbb R^n\),
\[
 (\iota_V)_\#\mathcal H_V^m=\mathcal H_{\mathbb R^n}^m\!\restriction V.
\]
This measure is a positive finite scalar multiple of the induced Haar measure. In particular its ball masses scale as \(r^m\); when \(m>0\), spheres have zero mass.
\par\emph{Role: Preiss proof.}
\par Lean declarations: \texttt{Preiss.flatOn, Preiss.map\_hausdorffMeasure\_subtype\_eq\_restrict, Preiss.planeMeasure\_ball\_and\_sphere, Preiss.unitFlat, Preiss.unitFlat\_closedBall}.
\end{lemma}
\begin{proof}[Planned proof obligations]
Split normalization, finiteness, and sphere-null corollaries; do not import tangent or uniform-measure theory.
\par Source: De Lellis Chapters 4--5; see the rectifiability audit.
\end{proof}

\begin{lemma}[A plane avoiding an inner cylinder has a small projected direction]
\label{lem:empty-cylinder-direction}
\lean{Preiss.cylinder, Preiss.exists_unit_small_projection_of_avoids_cylinder, Preiss.exists_unit_small_projection_of_empty_inner_cylinder, Preiss.exists_unit_small_projection_swapped}
\uses{def:plane}








The planes \(V,W\) have the same positive finite dimension, \(p\in V\), and \(r,\rho>0\) with \(r\) sufficiently larger than \(\rho\). No measure hypothesis is used.
If \(\|P_Vz-p\|=\rho\) and
\[
 (z+W)\cap B_r(z)\cap\{q:\|P_Vq-p\|<\rho/4\}=\varnothing,
\]
then some unit \(w\in W\) satisfies \(\|P_Vw\|\leq C\rho/r\).
\par\emph{Role: Preiss proof.}
\par Lean declarations: \texttt{Preiss.cylinder, Preiss.exists\_unit\_small\_projection\_of\_avoids\_cylinder}.
\end{lemma}
\begin{proof}[Planned proof obligations]
\par Source: De Lellis Chapters 4--5; see the rectifiability audit.
\end{proof}

\begin{definition}[Flat measures]
\label{def:flat}
\lean{Preiss.IsFlat}
\uses{def:plane}










A measure is $m$-flat if it equals $c\,\mathcal H^m\!\restriction V$ for an $m$-dimensional linear plane $V$ and a positive finite scalar $c$.
\end{definition}

\begin{definition}[Nonnegative compactly supported flatness energy]
\label{def:flatness}
\lean{Preiss.flatness, Preiss.flatnessEnergy}
\uses{def:plane}








For a compactly supported continuous cutoff $\phi$ and a measure $\mu$, set $E_\phi(\mu,V)=\int\phi(z)\operatorname{dist}(z,V)^2\,d\mu$ and $F_\phi(\mu)=\inf_{\dim V=m}E_\phi(\mu,V)$. The data definition is total. Positivity, attainment and continuity theorems impose $m\leq n$, local finiteness, and the relevant cutoff hypotheses.
\end{definition}
\begin{proof}[Planned proof obligations]
Choose $\phi\in C_c(\mathbb R^n)$ with $0\leq\phi\leq1$,
$\phi=1$ on $\overline B(0,1)$ and support in $B(0,2)$.
Define $F(\mu)=\min_V\int\phi(z)\operatorname{dist}(z,V)^2\,d\mu$.
The minimum is attained. S01 and C01 imply sequential continuity for vague
convergence; S02 implies continuity of $r\mapsto F(D^m_{x,r}\mu)$ for $r>0$.
Flat measures have $F=0$. For a unit uniform measure, $F\leq\varepsilon$
implies that some plane has unit-ball energy at most $\varepsilon$.
\end{proof}

\begin{definition}[Rectifiable sets and measures]
\label{def:rectifiable}
\lean{Preiss.CountablyRectifiable, Preiss.IsRectifiableMeasure}
\uses{def:parametrized-coverage}










Use the Borel measurable structure on the metric target.
A set $S$ is countably $m$-rectifiable if $\mathcal H^m(S\setminus T)=0$ for a countably Lipschitz-piece-covered set $T$. A measure is $m$-rectifiable if $\mu\ll\mathcal H^m$ and is concentrated on a countably $m$-rectifiable set. Local finiteness, measurability of carriers, and finite Hausdorff mass are separate theorem hypotheses. Absolute continuity absorbs the Hausdorff-null remainder.
\par\emph{Role: Preiss proof.}
\par Lean declarations: \texttt{Preiss.CountablyRectifiable, Preiss.IsRectifiableMeasure}.
\end{definition}

\begin{lemma}[Unit ball-law normalization adapter]
\label{lem:unit-uniform-adapter}
\lean{Preiss.isUnitUniform_iff}
\uses{def:uniform-api, def:unit-uniform}








For $m>0$, the pure unit-uniform ball law is equivalent to FormalizingGMT's $m$-uniform predicate together with origin support membership and unit closed-ball mass.
\end{lemma}

\begin{definition}[Conical measures]
\label{def:conical}
\lean{Preiss.IsConical}
\uses{def:dimensional-blowup}










A measure is $m$-conical if $D^m_{0,r}\nu=\nu$ for every $r>0$.
\end{definition}

\begin{lemma}[Gaussian moments of a normalized flat measure]
\label{lem:flat-gaussian-even-moments}
\lean{Preiss.flat_gaussianMoment_diagonal}
\uses{lem:unit-flat-conical, lem:zero-dimensional-flat-moments, lem:flat-moment-projection, lem:conical-even-support-moments, lem:conical-moment-scaling}







For $\sigma_V$ and $j\geq0$, its Gaussian moment form satisfies $b_{2j,s}(x^{2j})=s^j|P_Vx|^{2j}/j!$, and its odd moments are zero. Equivalently the one-dimensional standard Gaussian even moment is $(2j)!/(2^j j!)$.
\par\emph{Role: Preiss proof.}
\par Lean declarations: \texttt{Preiss.flat\_gaussianMoment\_diagonal, Preiss.flat\_gaussianMoment\_odd}.
\end{lemma}
\begin{proof}[Planned proof obligations]

For every $j\geq1$, $b_{2j}^{(j)}(x^{2j})=|Px|^{2j}$ for every $x$.
Use M02, M03 and the Gaussian moment formula of the flat limit, including the
factorial from M08.
\par Source: De Lellis Proposition 9.1, Lemmas 9.2,9.4,9.9,9.10 and Proposition 9.5, pp.95-109.
\end{proof}

\begin{definition}[Pure unrectifiability]
\label{def:purely-unrectifiable}
\lean{Preiss.PurelyUnrectifiable}
\uses{def:rectifiable}








Neither finite measure nor Borel measurability is built into the definition.
A set \(E\) is purely \(m\)-unrectifiable if
\[
 \mathcal H^m(E\cap R)=0
 \quad\text{for every countably \(m\)-rectifiable set }R.
\]
This canonical definition implies the graph-test version used in the source.
\par\emph{Role: Preiss proof.}
\par Lean declarations: \texttt{Preiss.PurelyUnrectifiable}.
\end{definition}

\begin{lemma}[Attainment, positivity and unit-ball comparison]
\label{lem:flatness-minimum}
\lean{Preiss.flatness_exists_minimizer, Preiss.flatness_nonneg, Preiss.flatness_eq_zero_of_isFlat, Preiss.exists_plane_ball_energy_le_of_flatness_le}
\uses{def:flat, def:flatness, def:plane, lem:grassmannian, lem:vague-moving-tests, def:flatness-cutoff, lem:flatness-basic-nonneg, lem:flatness-basic-zero, lem:flatness-minimizer-finite-compacts, lem:flatness-energy-integrable}





For every regular locally finite $\mu$, the minimum defining $F$ is attained and nonnegative. Flat measures have $F=0$. If $F(\mu)\leq\varepsilon$, there is an $m$-plane $V$ with $\int_{B(0,1)}\operatorname{dist}(z,V)^2\,d\mu\leq\varepsilon$.
\par\emph{Role: Preiss proof.}
\par Lean declarations: \texttt{Preiss.exists\_plane\_ball\_energy\_le\_of\_flatness\_le, Preiss.flatness\_eq\_zero\_of\_isFlat, Preiss.flatness\_exists\_minimizer, Preiss.flatness\_nonneg}.
\end{lemma}
\begin{proof}[Planned proof obligations]

Choose $\phi\in C_c(\mathbb R^n)$ with $0\leq\phi\leq1$,
$\phi=1$ on $\overline B(0,1)$ and support in $B(0,2)$.
Define $F(\mu)=\min_V\int\phi(z)\operatorname{dist}(z,V)^2\,d\mu$.
The minimum is attained. S01 and C01 imply sequential continuity for vague
convergence; S02 implies continuity of $r\mapsto F(D^m_{x,r}\mu)$ for $r>0$.
Flat measures have $F=0$. For a unit uniform measure, $F\leq\varepsilon$
implies that some plane has unit-ball energy at most $\varepsilon$.
\par Source: De Lellis Lemma 6.20 and Section 6.4, pp.66-69.
\end{proof}

\begin{lemma}[Continuity of flatness under vague convergence]
\label{lem:flatness-continuous}
\lean{Preiss.flatness_continuous}
\uses{def:flatness, lem:flatness-minimum, lem:grassmannian, lem:vague-moving-tests, def:flatness-cutoff, lem:flatness-vague-general}




For regular locally finite measures $\mu_i\rightharpoonup\mu$, $F(\mu_i)\to F(\mu)$.
\par\emph{Role: Preiss proof.}
\par Lean declarations: \texttt{Preiss.flatness\_continuous}.
\end{lemma}
\begin{proof}[Planned proof obligations]

Choose $\phi\in C_c(\mathbb R^n)$ with $0\leq\phi\leq1$,
$\phi=1$ on $\overline B(0,1)$ and support in $B(0,2)$.
Define $F(\mu)=\min_V\int\phi(z)\operatorname{dist}(z,V)^2\,d\mu$.
The minimum is attained. S01 and C01 imply sequential continuity for vague
convergence; S02 implies continuity of $r\mapsto F(D^m_{x,r}\mu)$ for $r>0$.
Flat measures have $F=0$. For a unit uniform measure, $F\leq\varepsilon$
implies that some plane has unit-ball energy at most $\varepsilon$.
\par Source: De Lellis Lemma 6.20 and Section 6.4, pp.66-69.
\end{proof}

\begin{lemma}[Density-to-tangent adapters]
\label{lem:uniform-adapter}
\lean{Preiss.ae_tangent_uniform}
\uses{def:positive-density, def:tangent-api, def:uniform-api, gmt:doubling-bounds, gmt:uniform-density-adapter, lem:density-bridge, thm:doubling, thm:uniform-upstream}








For $m>0$, positive finite existing density almost everywhere implies that almost everywhere all tangent measures are $m$-uniform and contain the origin in their support.
\end{lemma}

\begin{lemma}[Small Gaussian distance energy forces plane support proximity]
\label{lem:small-energy-no-holes}
\lean{Preiss.dist_support_lt_of_small_gaussian_plane_energy}
\uses{lem:weighted-energy-limit-support, lem:gaussian-plane-energy-integrable, lem:uniform-compactness, lem:plane-projection-subsequence, lem:plane-supported-uniform-flat, lem:unit-flat-support-regularity, lem:vague-support-approximation}




For every $\delta>0$ there exists $\eta>0$ such that if $\nu\in\mathcal U_m$
and $V$ is an $m$-plane with Gaussian distance energy at most $\eta$, every
$x\in V\cap\overline B(0,1)$ has distance less than $\delta$ to the support.
Otherwise rotate planes to one fixed plane, take a convergent subsequence of
measures and alleged holes, and use S04--S05. Vanishing energy puts the limit's
support in the plane; the plane-supported uniform classification identifies it
with $\sigma_V$, contradicting the persistent hole. This lemma does not
require conicality or eigenvalue assumptions.
\par\emph{Role: Preiss proof.}
\par Lean declarations: \texttt{Preiss.dist\_support\_lt\_of\_small\_gaussian\_plane\_energy}.
\end{lemma}
\begin{proof}[Planned proof obligations]

For every $\delta>0$ there exists $\eta>0$ such that if $\nu\in\mathcal U_m$
and $V$ is an $m$-plane with Gaussian distance energy at most $\eta$, every
$x\in V\cap\overline B(0,1)$ has distance less than $\delta$ to the support.
Otherwise rotate planes to one fixed plane, take a convergent subsequence of
measures and alleged holes, and use S04--S05. Vanishing energy puts the limit's
support in the plane; the plane-supported uniform classification identifies it
with $\sigma_V$, contradicting the persistent hole. This lemma does not
require conicality or eigenvalue assumptions.
\par Source: De Lellis Corollary 8.4, Lemmas 8.6-8.7, Proposition 8.5, pp.85-94; author errata for the dimension-two shortcut.
\end{proof}

\begin{lemma}[Homogeneity and parity of conical moments]
\label{lem:conical-moments}
\lean{Preiss.conical_gaussianMoment_homogeneous, Preiss.conical_gaussianMoment_odd}
\uses{lem:conical-moment-scaling, lem:scaled-moment-limits, lem:taylor-two}




Conicality gives $b_{k,s}=s^{k/2}b_{k,1}$. Comparing with M07 makes odd forms
zero and gives $b_{2j,s}=s^j b_{2j}^{(j)}/j!$.
The support identities give $b_{2j}^{(j)}(x^{2j})=|x|^{2j}$ on the support.
\par\emph{Role: Preiss proof.}
\par Lean declarations: \texttt{Preiss.conical\_gaussianMoment\_homogeneous, Preiss.conical\_gaussianMoment\_odd}.
\end{lemma}
\begin{proof}[Planned proof obligations]

Conicality gives $b_{k,s}=s^{k/2}b_{k,1}$. Comparing with M07 makes odd forms
zero and gives $b_{2j,s}=s^j b_{2j}^{(j)}/j!$.
The support identities give $b_{2j}^{(j)}(x^{2j})=|x|^{2j}$ on the support.
\par Source: De Lellis Corollary 8.4, Lemmas 8.6-8.7, Proposition 8.5, pp.85-94; author errata for the dimension-two shortcut.
\end{proof}

\begin{lemma}[All Gaussian polynomial moments of blow-downs have limits]
\label{lem:blowdown-moment-limits}
\lean{Preiss.blowdown_gaussian_polynomial_limit}
\uses{lem:coordinate-product-representation, lem:blowdown-product-limits, lem:gaussian-product-moment-conversion}




For $\nu\in\mathcal U_m$ and every polynomial $P$, the limit of $\int P(z)e^{-|z|^2}\,d(D^m_{0,r}\nu)$ exists as $r\to\infty$. Degree $2j$ moment forms are determined by $b_{2j}^{(j)}/j!$, and odd-degree moments vanish.
\par\emph{Role: Preiss proof.}
\par Lean declarations: \texttt{Preiss.blowdown\_gaussian\_polynomial\_limit}.
\end{lemma}
\begin{proof}[Planned proof obligations]

For odd $k$, $s^{-k/2}b_{k,s}\to0$; for $k=2j$,
$s^{-j}b_{2j,s}\to b_{2j}^{(j)}/j!$.
The factorial is essential; the displayed summary on printed page 74 omits it,
whereas the calculation in Chapter 9 includes it correctly.
Use M03 with $s=r^{-2}$ to deduce the limit of every Gaussian-weighted
polynomial moment of $D^m_{0,r}\nu$ as $r\to\infty$.
\par Source: De Lellis Lemma 7.2, Definition 7.4, Lemma 7.6, Proposition 7.7, pp.70-84.
\end{proof}

\begin{lemma}[Positive covariance form, trace and support identity]
\label{lem:conical-covariance}
\lean{Preiss.conical_covariance_trace_support}
\uses{lem:even-moment-positivity, lem:moment-second-total-trace, lem:conical-support-covariance}





The symmetric form $B=b_{2,1}$ is positive semidefinite,
$\operatorname{tr}B=m$, and $B(x,x)=|x|^2$ on the support.
The trace calculation is the radial Gaussian second moment. This statement
does not require spectral choices.
\par\emph{Role: Preiss proof.}
\par Lean declarations: \texttt{Preiss.conical\_covariance\_trace\_support}.
\end{lemma}
\begin{proof}[Planned proof obligations]

The symmetric form $B=b_{2,1}$ is positive semidefinite,
$\operatorname{tr}B=m$, and $B(x,x)=|x|^2$ on the support.
The trace calculation is the radial Gaussian second moment. This statement
does not require spectral choices.
\par Source: De Lellis Corollary 8.4, Lemmas 8.6-8.7, Proposition 8.5, pp.85-94; author errata for the dimension-two shortcut.
\end{proof}

\section{Rectifiable carriers and cones}

\begin{lemma}[Graph parametrization is Lipschitz and antilipschitz]
\label{lem:graph-parametrization}
\lean{Preiss.graphMap, Preiss.graphMap_lipschitz_antilipschitz}
\uses{def:lipschitz-graph}








The subspace \(V\) is finite-dimensional. A finite Lipschitz bound for \(f\) is given; an optimal bound for \(F\) is unnecessary.
If \(f:V\to V^\perp\) is Lipschitz, then \(F(v)=v+f(v)\) is Lipschitz and satisfies
\[
 \|u-v\|\leq\|F(u)-F(v)\|,\qquad P_V(F(v))=v.
\]
Thus \(F\) is injective and its inverse on its range is \(1\)-Lipschitz.
\par\emph{Role: Preiss proof.}
\par Lean declarations: \texttt{Preiss.graphMap, Preiss.graphMap\_lipschitz\_antilipschitz}.
\end{lemma}
\begin{proof}[Planned proof obligations]
\par Source: De Lellis Chapters 4--5; see the rectifiability audit.
\end{proof}

\begin{lemma}[Pairwise cone condition implies graph containment]
\label{lem:cone-graph}
\lean{Preiss.exists_lipschitzGraph_of_pairwise_mem_cone}
\uses{def:cone, lem:domain-extension}








The set \(F\subset\mathbb R^n\) is arbitrary, \(V\) is an \(m\)-plane, and \(0\leq\alpha<\infty\). No measure hypothesis is needed.
If \(y\in C(x,V,\alpha)\) for every \(x,y\in F\), then there is a Lipschitz \(f:V\to V^\perp\) such that
\[
 F\subset\{v+f(v):v\in V\}.
\]
\par\emph{Role: Preiss proof.}
\par Lean declarations: \texttt{Preiss.exists\_lipschitzGraph\_of\_pairwise\_mem\_cone}.
\end{lemma}
\begin{proof}[Planned proof obligations]

\par Source: De Lellis Chapters 4--5; see the rectifiability audit.
\end{proof}

\begin{lemma}[Rectifiability closure operations]
\label{lem:carrier-closure}
\lean{Preiss.countablyRectifiable_iUnion, Preiss.countablyRectifiable_image, Preiss.countablyRectifiable_mono, Preiss.countablyRectifiable_of_sdiff_null, Preiss.countablyRectifiable_union}
\uses{def:coverage, def:rectifiable}








Use subset-domain parametrizations in metric targets. For images, source and target Hausdorff measures have the same dimension \(m\).
Countably \(m\)-rectifiable sets are closed under taking subsets, countable unions, Lipschitz images, and modifications by \(\mathcal H^m\)-null sets.
\par\emph{Role: Preiss proof.}
\par Lean declarations: \texttt{Preiss.countablyRectifiable\_iUnion, Preiss.countablyRectifiable\_image, Preiss.countablyRectifiable\_mono, Preiss.countablyRectifiable\_of\_sdiff\_null}.
\end{lemma}
\begin{proof}[Planned proof obligations]
Several elementary declarations in one API, not a single conjunction theorem.
\par Source: De Lellis Chapters 4--5; see the rectifiability audit.
\end{proof}

\begin{lemma}[Rectifiability of the zero measure]
\label{lem:zero-measure}
\lean{Preiss.zero_isRectifiableMeasure}
\uses{def:rectifiable}








The zero measure is rectifiable in every natural dimension, on every metric target with its Borel measurable structure.
\end{lemma}

\begin{lemma}[Graphs and their subsets are rectifiable]
\label{lem:graph-rectifiable}
\lean{Preiss.countablyRectifiable_subset_graph}
\uses{def:rectifiable, lem:graph-parametrization}








Use the canonical definition by Lipschitz parametrizations on subsets of \(\mathbb R^m\); the target is Euclidean.
Every full Lipschitz \(m\)-graph is countably \(m\)-rectifiable. Consequently every subset of such a graph is countably \(m\)-rectifiable.
\par\emph{Role: Preiss proof.}
\par Lean declarations: \texttt{Preiss.countablyRectifiable\_subset\_graph}.
\end{lemma}
\begin{proof}[Planned proof obligations]

\par Source: De Lellis Chapters 4--5; see the rectifiability audit.
\end{proof}

\begin{lemma}[Borel sigma-finite carrier]
\label{lem:borel-carrier}
\lean{Preiss.exists_borel_sigmaFinite_carrier}
\uses{lem:carrier-closure, lem:domain-extension, lem:plane-hausdorff}







The target is Euclidean; \(m\) is a natural dimension. No measurability assumption is part of the rectifiability predicate for \(S\).
If \(S\subset\mathbb R^n\) is countably \(m\)-rectifiable, there exists a Borel countable union \(T\) of compact Lipschitz images such that
\[
 \mathcal H^m(S\setminus T)=0,\qquad
 \mathcal H^m\!\restriction T\text{ is sigma-finite}.
\]
\par\emph{Role: Preiss proof.}
\par Lean declarations: \texttt{Preiss.exists\_borel\_sigmaFinite\_carrier}.
\end{lemma}
\begin{proof}[Planned proof obligations]

\par Source: De Lellis Chapters 4--5; see the rectifiability audit.
\end{proof}

\begin{lemma}[Approximate cone criterion]
\label{lem:cone-criterion}
\lean{Preiss.countablyRectifiable_of_ae_cone_approximation}
\uses{def:rectifiable, gmt:hausdorff-density-outside, gmt:hausdorff-density-upper, gmt:vitali, lem:cone-graph}








For a Borel set of finite $\Hm$ measure, positive lower density and
vanishing $r^{-m}$-normalized mass outside a fixed cone at almost every point
imply countable rectifiability (De Lellis Proposition 4.6).
\end{lemma}
\begin{proof}[Planned proof]

Stratify cone directions, apertures, lower-density bounds, and error scales
into countably many sets satisfying the pairwise cone condition.
\end{proof}

\begin{lemma}[Pure unrectifiability calculus]
\label{lem:pure-basic}
\lean{Preiss.purelyUnrectifiable_measure_inter_eq_zero, Preiss.purelyUnrectifiable_mono}
\uses{def:purely-unrectifiable, lem:graph-rectifiable}








Use the canonical pure predicate defined by testing all countably rectifiable sets.
Every subset of a purely \(m\)-unrectifiable set is purely \(m\)-unrectifiable. Its intersection with any countably \(m\)-rectifiable set, in particular any Lipschitz \(m\)-graph, is \(\mathcal H^m\)-null.
\par\emph{Role: Preiss proof.}
\par Lean declarations: \texttt{Preiss.purelyUnrectifiable\_measure\_inter\_eq\_zero, Preiss.purelyUnrectifiable\_mono}.
\end{lemma}
\begin{proof}[Planned proof obligations]

\par Source: De Lellis Chapters 4--5; see the rectifiability audit.
\end{proof}

\begin{lemma}[Carrier API aggregation]
\label{lem:carrier-calculus}
\lean{Preiss.isRectifiableMeasure_of_countable_restrict}
\uses{lem:carrier-closure}







The measure is absolutely continuous with respect to \(\mathcal H^m\). Restriction carriers are measurable and form a countable family covering the measure up to a null set.
Hausdorff-null changes of carriers preserve rectifiability of measures absolutely continuous with respect to \(\mathcal H^m\). Countably many rectifiable restrictions whose carriers cover almost all the measure assemble to a rectifiable measure.
\par\emph{Role: Preiss proof.}
\par Lean declarations: \texttt{Preiss.isRectifiableMeasure\_of\_countable\_restrict}.
\end{lemma}
\begin{proof}[Planned proof obligations]

Refines existing bundled node; declaration names for the individual elementary facts should be selected during implementation.
\par Source: De Lellis Chapters 4--5; see the rectifiability audit.
\end{proof}

\begin{lemma}[Rectifiable/pure decomposition]
\label{lem:decomposition}
\lean{Preiss.exists_rectifiable_purelyUnrectifiable_partition, Preiss.rectifiable_pure_partition_unique}
\uses{lem:carrier-closure, lem:pure-basic}







The set \(E\subset\mathbb R^n\) is Borel and \(\mathcal H^m(E)<\infty\).
There exist disjoint Borel sets \(R,U\subset E\) with
\[
 E=R\cup U,\qquad R\text{ countably \(m\)-rectifiable},\qquad
 U\text{ purely \(m\)-unrectifiable}.
\]
Any two such decompositions agree up to \(\mathcal H^m\)-null sets.
\par\emph{Role: Preiss proof.}
\par Lean declarations: \texttt{Preiss.exists\_rectifiable\_purelyUnrectifiable\_partition, Preiss.rectifiable\_pure\_partition\_unique}.
\end{lemma}
\begin{proof}[Planned proof obligations]

Maximize mass among Borel canonical rectifiable subsets. Graph-only source definition is not required.
\par Source: De Lellis Chapters 4--5; see the rectifiability audit.
\end{proof}

\section{Density and dimension}

\begin{definition}[Positive finite density]
\label{def:positive-density}
\lean{Preiss.HasPositiveFiniteDensitiesAt, Preiss.HasPositiveFiniteDensityAt, Preiss.HasPositiveFiniteDensityAt.densities, Preiss.HasPositiveFiniteDensityAt.hasDensity}
\uses{def:density-api}










The predicate \texttt{HasPositiveFiniteDensityAt} means
$0<\Theta_*^m=\Theta^{*m}<\infty$ in the upstream convention.
The plural \texttt{HasPositiveFiniteDensitiesAt} means only
$0<\Theta_*^m\leq\Theta^{*m}<\infty$.
\end{definition}

\begin{lemma}[Measure comparisons on density strata]
\label{lem:density-measure-comparison}
\lean{Preiss.UpperPowerBoundAt, Preiss.hausdorff_restrict_le_measure_of_ball_bounds, Preiss.measure_restrict_le_hausdorff_of_ball_bounds}
\uses{mathlib:vitali-enlargement, lem:ball-upper-hausdorff-bound, lem:lower-ball-restrict-comparison}








At every centre of \(S\), the appropriate power bound holds for all radii below one common positive cutoff. The measure is locally finite; localize to bounded strata when claiming finiteness.
Uniform small-ball upper bounds on a Borel stratum \(S\) imply
\[
 \mu\!\restriction S\leq C\,\mathcal H^m\!\restriction S.
\]
Uniform positive lower bounds give the reverse comparison, with a finite geometric constant, and finite \(\mathcal H^m\)-measure on bounded strata.
\par\emph{Role: Preiss proof.}
\par Lean declarations: \texttt{Preiss.UpperPowerBoundAt, Preiss.hausdorff\_restrict\_le\_measure\_of\_ball\_bounds, Preiss.measure\_restrict\_le\_hausdorff\_of\_ball\_bounds}.
\end{lemma}
\begin{proof}[Planned proof obligations]

The upper comparison follows from the generic metric-space diameter bound.
For the lower comparison, fix an open superset of the set to be estimated.
Choose arbitrarily small balls centred on that set, contained in the open
superset, and satisfying the lower mass bound. The disjoint-family Vitali
theorem supplies a countable cover by fivefold enlarged balls. Their diameter
cost is bounded by $10^m/a$ times the total mass of the disjoint original balls.
Let the covering diameters tend to zero and then use outer regularity to
shrink the open superset. This argument only needs a separable metric space
and an outer regular measure; the Euclidean statement is a specialization.
\par Source: De Lellis Chapters 4--5; see the rectifiability audit.
\end{proof}

\begin{lemma}[Dimension zero]
\label{lem:zero-dimension}
\lean{Preiss.preiss_zero}
\uses{def:positive-density, def:rectifiable, lem:countable-carrier, lem:zero-hausdorff-ac, lem:zero-density-atoms}










In dimension zero the density hypothesis implies rectifiability.
\end{lemma}
\begin{proof}[Planned proof]
Continuity from above of the finite measure on a bounded ball identifies
the small-ball limit with $\mu\{x\}$. The measure is carried by its atoms;
local finiteness and a countable compact exhaustion make these countable.
Zero-dimensional Hausdorff measure has no nonempty null sets.
\end{proof}

\begin{lemma}[Countable carriers]
\label{lem:countable-carrier}
\lean{Preiss.countablyCoveredByLipschitzPieces_of_countable, Preiss.countablyRectifiable_of_countable}
\uses{def:parametrized-coverage, def:rectifiable}
Every countable subset of a metric space is covered by countably many constant
Lipschitz parametrizations, in every natural dimension. It is therefore countably
rectifiable. The empty set is included.
\end{lemma}
\begin{proof}
Associate to each point the constant map on the whole Euclidean domain.
\end{proof}

\begin{lemma}[Absolute continuity with respect to zero-dimensional Hausdorff measure]
\label{lem:zero-hausdorff-ac}
\lean{Preiss.absolutelyContinuous_hausdorffMeasure_zero}
Every Borel measure on a metric space is absolutely continuous with respect to
$\mathcal H^0$, since a nonempty set has $\mathcal H^0$-measure at least one.
\end{lemma}

\begin{lemma}[Zero-dimensional density equals atomic mass]
\label{lem:zero-density-atoms}
\lean{Preiss.tendsto_densityRatio_zero, Preiss.lowerDensity_zero_eq_singleton}
\uses{def:density-api}
For a locally finite measure on a proper metric space, the zero-dimensional
density ratio at $x$ converges to $\mu\{x\}$ as the radius tends to zero from above.
In particular the lower zero-dimensional density equals this atomic mass.
\end{lemma}
\begin{proof}
Closed balls decrease to the singleton, and one such ball has finite measure.
Apply continuity from above.
\end{proof}

\begin{lemma}[Normalization bridges]
\label{lem:density-bridge}
\uses{lem:density-limit-values}
\lean{Preiss.density_iff_upstream, Preiss.density_iff_openBall_limit}
\uses{def:positive-density, gmt:density-bounds, gmt:density-exists-set, gmt:density-limit}










Positive finite equal densities are equivalent to membership in the existing
positive-finite-density-exists set, and to a positive finite limit of
$\mu(B(x,r))/r^m$.
\end{lemma}
\begin{proof}[Planned proof]

Use the complete order topology of ENNReal. Squeeze open and closed balls
using radii $(1\pm\varepsilon)r$, then let $\varepsilon\downarrow0$.
The radius-normalized limit is $2^m$ times the diameter-normalized one.
De Lellis further divides by $\omega_m$; the resulting constant is positive
and finite. Equality, positivity, and finiteness are invariant under this
change, but exact constants in uniform-measure lemmas must be transported.
\end{proof}

\begin{lemma}[Finite upper density implies Hausdorff absolute continuity]
\label{lem:absolute-continuity}
\lean{Preiss.absolutelyContinuous_hausdorff_of_upper_density_finite}
\uses{lem:upper-density-localization, lem:ball-upper-hausdorff-bound, gmt:goodballs}







\(\mu\) is a Borel measure on \(\mathbb R^n\); \(m\) is a natural dimension.
If \(\Theta^{*m}(\mu,x)<\infty\) for \(\mu\)-almost every \(x\), then \(\mu\ll\mathcal H^m\).
\par\emph{Role: Preiss proof.}
\par Lean declarations: \texttt{Preiss.absolutelyContinuous\_hausdorff\_of\_upper\_density\_finite}.
\end{lemma}
\begin{proof}[Planned proof obligations]

\par Source: De Lellis Chapters 4--5; see the rectifiability audit.
\end{proof}

\begin{lemma}[Finite Hausdorff pieces with comparable measures]
\label{lem:finite-pieces}
\lean{Preiss.exists_finite_hausdorff_pieces_of_density_bounds}
\uses{gmt:density-bounds, gmt:goodsets, lem:density-measure-comparison, lem:density-bounds-good-set}







\(\mu\) is locally finite and \(0<\Theta_*^m(\mu,x)\leq\Theta^{*m}(\mu,x)<\infty\) for \(\mu\)-almost every \(x\).
There are bounded Borel sets \(S_j\) carrying \(\mu\), each of finite \(\mathcal H^m\)-measure, such that \(\mu\!\restriction S_j\) and \(\mathcal H^m\!\restriction S_j\) have positive finite two-sided comparison constants. In particular,
\[
 \mu\bigl((\bigcup_j S_j)^c\bigr)=0.
\]
\par\emph{Role: Preiss proof.}
\par Lean declarations: \texttt{Preiss.exists\_finite\_hausdorff\_pieces\_of\_density\_bounds}.
\end{lemma}
\begin{proof}[Planned proof obligations]

\par Source: De Lellis Chapters 4--5; see the rectifiability audit.
\end{proof}

\begin{lemma}[Dimensions above the ambient dimension]
\label{lem:excess-dimension}
\lean{Preiss.measure_zero_of_dimension_gt}
\uses{def:positive-density, lem:uniform-adapter, lem:tangent-existence, gmt:uniform-dimension-bound}
If $m>n$, positive finite existing $m$-density almost everywhere forces $\mu=0$.
\end{lemma}
\begin{proof}[Proof route]
A positive-density point has a tangent measure. Almost every such point has only
$m$-uniform tangents, contradicting the upstream ambient-dimension bound.
Thus the density hypothesis is concentrated on a null set and $\mu=0$.
\end{proof}

\section{Vague limits and tangent measures}

\begin{definition}[Dimensionally normalized blow-ups]
\label{def:dimensional-blowup}
\lean{Preiss.dimensionalBlowup}

Set $D^m_{x,r}\mu=r^{-m}(T_{x,r})_\#\mu$, where $T_{x,r}(y)=(y-x)/r$. The definition is total in Lean; geometric theorems assume $r>0$.
\end{definition}

\begin{definition}[Dimensional tangent measures]
\label{def:dimensional-tangent}
\lean{Preiss.IsDimensionalTangent}
\uses{def:dimensional-blowup, gmt:vague}

A nonzero regular measure $\nu$ is a dimensional tangent of $\mu$ at $x$ if $D^m_{x,r_i}\mu\rightharpoonup\nu$ for some positive sequence $r_i\to0$. This fixes the scalar normalization.
\end{definition}

\begin{definition}[Tangents at infinity]
\label{def:tangent-at-infinity}
\lean{Preiss.IsTangentAtInfinity}
\uses{def:dimensional-blowup, gmt:vague}

A nonzero regular measure $\nu$ is a tangent at infinity of $\mu$ if $D^m_{0,r_i}\mu\rightharpoonup\nu$ for some positive sequence $r_i\to\infty$.
\end{definition}

\begin{lemma}[Flat limit excludes persistent mass away from its plane]
\label{lem:flat-limit-tube}
\lean{Preiss.tendsto_mass_outside_tube_of_flat_limit}
\uses{gmt:portmanteau, lem:plane-hausdorff, lem:flat-measure-regularity, lem:vague-compact-null-mass}
All measures are regular, \(V\) is an \(m\)-plane, \(0<c<\infty\), and \(\eta>0\). Use the compact closure of the region.
If \(\nu_j\rightharpoonup c\mathcal H^m\!\restriction V\), then the mass of any fixed bounded set separated by a positive distance from \(V\) tends to zero. In particular this applies to the region outside an \(\eta\)-tube inside a fixed ball.
\end{lemma}


\begin{lemma}[Flat limit gives lower mass in moving plane balls]
\label{lem:flat-limit-moving-balls}
\lean{Preiss.eventually_ball_mass_lower_bound_of_flat_limit}
\uses{gmt:portmanteau, lem:plane-hausdorff, lem:flat-measure-regularity, lem:vague-uniform-ball-lower}







All measures are regular, \(0<c<\infty\), and \(V\) is an \(m\)-plane. Extract converging centres and use a fixed smaller open ball.
If \(\nu_j\rightharpoonup c\mathcal H^m\!\restriction V\), then for each \(R,\eta>0\) there are \(b>0\) and \(j_0\) such that
\[
 \nu_j(B_\eta(z))\geq b
 \quad(j\geq j_0,\ z\in V\cap\overline B_R(0)).
\]
\par\emph{Role: Preiss proof.}
\par Lean declarations: \texttt{Preiss.eventually\_ball\_mass\_lower\_bound\_of\_flat\_limit}.
\end{lemma}
\begin{proof}[Planned proof obligations]

\par Source: De Lellis Chapters 4--5; see the rectifiability audit.
\end{proof}

\begin{lemma}[Fixed-normalization compactness at density-bound points]
\label{lem:dimensional-compactness}
\lean{Preiss.exists_dimensional_tangent_subseq_of_density_bounds}
\uses{lem:dimensional-extraction-ball-bounds, gmt:density-bounds}







\(\mu\) is locally finite and \(0<\Theta_*^m(\mu,x)\leq\Theta^{*m}(\mu,x)<\infty\) at the fixed point \(x\).
For every positive sequence \(r_j\to0\), there are a strictly increasing subsequence map \(\phi\) and a nonzero regular measure \(\nu\) such that
\[
 D^m_{x,r_{\phi(j)}}\mu\rightharpoonup\nu.
\]
The limit is a tangent measure in GMT's arbitrary-scalar convention.
\par\emph{Role: Preiss proof.}
\par Lean declarations: \texttt{Preiss.exists\_dimensional\_tangent\_subseq\_of\_density\_bounds}.
\end{lemma}
\begin{proof}[Planned proof obligations]

\par Source: De Lellis Chapters 4--5; see the rectifiability audit.
\end{proof}

\begin{lemma}[Composition, scalar and isometric covariance of blow-ups]
\label{lem:blowup-composition}
\lean{Preiss.dimensionalBlowup_comp, Preiss.dimensionalBlowup_smul, Preiss.dimensionalBlowup_map_linearIsometry}
\uses{def:dimensional-blowup, gmt:blowup, gmt:blowup-regular}






Prove
$D^m_{b,s}(D^m_{a,r}\mu)=D^m_{a+rb,rs}\mu$;
positive scalar multiplication commutes with $D$; and pushforward by a linear
isometric embedding intertwines $D$, at centres in its image. Include inverse
dilation and preservation of flatness, regularity and uniform ball constants.
\par\emph{Role: Preiss proof.}
\par Lean declarations: \texttt{Preiss.dimensionalBlowup\_comp, Preiss.dimensionalBlowup\_map\_linearIsometry, Preiss.dimensionalBlowup\_smul}.
\end{lemma}
\begin{proof}[Planned proof obligations]

Prove
$D^m_{b,s}(D^m_{a,r}\mu)=D^m_{a+rb,rs}\mu$;
positive scalar multiplication commutes with $D$; and pushforward by a linear
isometric embedding intertwines $D$, at centres in its image. Include inverse
dilation and preservation of flatness, regularity and uniform ball constants.
\par Source: De Lellis Chapters 3 and 6; general supporting analysis.
\end{proof}

\begin{lemma}[Support approximation from arbitrary positive outer measure sets]
\label{lem:hull-support-approximation}
\lean{Preiss.IsMeasureDensityOnePoint, Preiss.dist_badSet_rescaled_tendsto_zero, Preiss.exists_badSet_points_rescaled_tendsto, Preiss.dist_set_rescaled_tendsto_zero_of_densityOne}
\uses{def:dimensional-tangent, gmt:density-points, gmt:portmanteau, lem:measurable-hull-holes, lem:rescaled-distance-selection}






Assume $a\in D$, $a$ is a density-one point of its measurable hull $H$,
$\limsup_{r\downarrow0}\mu(B(a,r))/r^m<\infty$, and
$D^m_{a,r_i}\mu\rightharpoonup\nu$ with $r_i\downarrow0$.
Then for every $y\in\operatorname{supp}\nu$,
$\operatorname{dist}(a+r_i y,D)/r_i\to0$.
Indeed, a subsequence of holes $B(a+r_i y,\eta r_i)$ disjoint from $D$ has
zero $H$-mass by A07. It lies in $B(a,Cr_i)$ for a fixed $C$; its remaining
mass is $o(r_i^m)$ by density-one and the upper density bound. Portmanteau
would then make $\nu(B(y,\eta))=0$, a contradiction. Choose approximate
nearest points in $D$ with error $r_i/i$; no closedness of $D$ is needed.
\par\emph{Role: Preiss proof.}
\par Lean declarations: \texttt{Preiss.IsMeasureDensityOnePoint, Preiss.dist\_badSet\_rescaled\_tendsto\_zero, Preiss.exists\_badSet\_points\_rescaled\_tendsto}.
\end{lemma}
\begin{proof}[Planned proof obligations]

Assume $a\in D$, $a$ is a density-one point of its measurable hull $H$,
$\limsup_{r\downarrow0}\mu(B(a,r))/r^m<\infty$, and
$D^m_{a,r_i}\mu\rightharpoonup\nu$ with $r_i\downarrow0$.
Then for every $y\in\operatorname{supp}\nu$,
$\operatorname{dist}(a+r_i y,D)/r_i\to0$.
Indeed, a subsequence of holes $B(a+r_i y,\eta r_i)$ disjoint from $D$ has
zero $H$-mass by A07. It lies in $B(a,Cr_i)$ for a fixed $C$; its remaining
mass is $o(r_i^m)$ by density-one and the upper density bound. Portmanteau
would then make $\nu(B(y,\eta))=0$, a contradiction. Choose approximate
nearest points in $D$ with error $r_i/i$; no closedness of $D$ is needed.
\par Corrections and boundary conditions: New measurable-hull repair; require density one of the hull and finite upper dimensional density.
\par Source: De Lellis Lemma 6.6, Theorem 6.8, pp.57-58,63-65; Corollary 3.9 and corrected Lemma 3.8, pp.15,25-26; author errata.
\end{proof}

\begin{lemma}[Joint continuity of translated dilations]
\label{lem:vague-moving-blowups}
\lean{Preiss.vague_smul_map_sub_div, Preiss.vague_dimensionalBlowup_moving}
\uses{def:dimensional-blowup, gmt:blowup-regular, lem:vague-moving-tests}







If $\mu_i\rightharpoonup\mu$, $a_i\to a$, $r_i\to r>0$, and
$c_i\to c\in[0,\infty)$, then
$c_i(T_{a_i,r_i})_\#\mu_i\rightharpoonup c(T_{a,r})_\#\mu$.
Assume regularity of the measures and finite real scalars; no bound at infinity
on their total masses is needed. The preimages of a fixed compact test support
lie in one compact set because the parameters converge and $r>0$. The pulled
back tests converge uniformly there. S01 proves the statement. Its special case
$D^m_{a_i,r_i}\mu_i\to D^m_{a,r}\mu$ is used repeatedly.
\par\emph{Role: Preiss proof.}
\par Lean declarations: \texttt{Preiss.vague\_dimensionalBlowup\_moving, Preiss.vague\_smul\_map\_sub\_div}.
\end{lemma}
\begin{proof}[Planned proof obligations]

If $\mu_i\rightharpoonup\mu$, $a_i\to a$, $r_i\to r>0$, and
$c_i\to c\in[0,\infty)$, then
$c_i(T_{a_i,r_i})_\#\mu_i\rightharpoonup c(T_{a,r})_\#\mu$.
Assume regularity of the measures and finite real scalars; no bound at infinity
on their total masses is needed. The preimages of a fixed compact test support
lie in one compact set because the parameters converge and $r>0$. The pulled
back tests converge uniformly there. S01 proves the statement. Its special case
$D^m_{a_i,r_i}\mu_i\to D^m_{a,r}\mu$ is used repeatedly.
\par Corrections and boundary conditions: Require limiting radius strictly positive; scalars finite and convergent; obtain one common compact set of pulled-back supports.
\par Source: De Lellis Chapters 3 and 6; general supporting analysis.
\end{proof}

\begin{lemma}[Closedness of arbitrary-scalar tangents]
\label{lem:tangent-closed}
\lean{Preiss.tangent_closed}
\uses{def:tangent-api, lem:vague-diagonal, gmt:blowup-regular}




The source measure $\mu$ is locally finite.




If every $\nu_i\in\operatorname{Tan}(\mu,x)$, $\nu_i\rightharpoonup\nu$, and $\nu$ is regular and nonzero, then $\nu\in\operatorname{Tan}(\mu,x)$.
\par\emph{Role: Supporting API.}
\par Lean declarations: \texttt{Preiss.tangent\_closed}.
\end{lemma}
\begin{proof}[Planned proof obligations]

For a fixed measure and centre, the limit of a vaguely convergent sequence of
dimensional tangents is a dimensional tangent whenever that limit is regular
and nonzero. Prove the analogous arbitrary-scalar statement as a separate
adapter. Constant positive dilations of a dimensional tangent remain dimensional
tangents. At a positive finite density point the dimensional tangent set is
nonempty and compact, because every sufficiently small blow-up has compact
mass bounds and nonzero subsequential limits. A countable-test diagonal lemma
implements closure without assuming a topology on all raw measures.
\par Source: De Lellis Chapters 3 and 6; general supporting analysis.
\end{proof}

\begin{lemma}[Closedness of dimensional tangents]
\label{lem:dimensional-tangent-closed}
\lean{Preiss.dimensionalTangent_closed}
\uses{def:dimensional-tangent, lem:vague-diagonal, lem:dimensional-blowup-support}




The source measure $\mu$ is locally finite.





If every $\nu_i\in\operatorname{Tan}_m(\mu,x)$, $\nu_i\rightharpoonup\nu$, and $\nu$ is regular and nonzero, then $\nu\in\operatorname{Tan}_m(\mu,x)$.
\par\emph{Role: Preiss proof.}
\par Lean declarations: \texttt{Preiss.dimensionalTangent\_closed}.
\end{lemma}
\begin{proof}[Planned proof obligations]

For a fixed measure and centre, the limit of a vaguely convergent sequence of
dimensional tangents is a dimensional tangent whenever that limit is regular
and nonzero. Prove the analogous arbitrary-scalar statement as a separate
adapter. Constant positive dilations of a dimensional tangent remain dimensional
tangents. At a positive finite density point the dimensional tangent set is
nonempty and compact, because every sufficiently small blow-up has compact
mass bounds and nonzero subsequential limits. A countable-test diagonal lemma
implements closure without assuming a topology on all raw measures.
\par Source: De Lellis Chapters 3 and 6; general supporting analysis.
\end{proof}

\begin{lemma}[All flat tangents imply weak linear approximation]
\label{lem:flat-to-approximation}
\lean{Preiss.weaklyLinearlyApproximableAt_of_all_tangents_flat}
\uses{def:weak-approximation, lem:dimensional-compactness, lem:flat-limit-moving-balls, lem:flat-limit-tube, lem:sequential-uniform-scale, lem:weak-approximation-scale-transfer}








The set \(F\) is Borel with \(\mathcal H^m(F)<\infty\), and \(\mathcal H^m\!\restriction F\) has positive lower and finite upper \(m\)-density at \(x\).
If every tangent of \(\mathcal H^m\!\restriction F\) at \(x\) is \(m\)-flat, then \(F\) is weakly linearly approximable at \(x\):
for every \(\eta>0\), some \(\lambda,r_0>0\) work at every \(0<r<r_0\) for both the tube and plane-centred-ball clauses.
\par\emph{Role: Preiss proof.}
\par Lean declarations: \texttt{Preiss.weaklyLinearlyApproximableAt\_of\_all\_tangents\_flat}.
\end{lemma}
\begin{proof}[Planned proof obligations]

\par Source: De Lellis Chapters 4--5; see the rectifiability audit.
\end{proof}

\begin{lemma}[Plane and scaling API]
\label{lem:plane-scaling}
\lean{Preiss.isFlat_dimensionalBlowup, Preiss.isFlat_of_dimensionalBlowup_zero, Preiss.isFlat_map_linearIsometry}
\uses{def:dimensional-blowup, def:flat, lem:blowup-composition, lem:plane-hausdorff, lem:measure-support-transport, lem:flat-plane-blowup}





Flat measures are regular and nonzero. Dimensional dilation and recentering at a point of their supporting plane preserve flatness; pushforward by a linear isometry gives the corresponding flat measure on the image plane. Unit-uniform scaling and blow-up composition have separate interfaces.
\end{lemma}

\begin{lemma}[Tangent existence under two-sided density bounds]
\label{lem:tangent-existence}
\lean{Preiss.tangent_exists}
\uses{def:positive-density, def:tangent-api, gmt:doubling-bounds, thm:doubling}








At a point with positive lower and finite upper $m$-density, a locally finite measure has a nonzero tangent measure.
\end{lemma}

\begin{lemma}[Positive dilations preserve dimensional tangents]
\label{lem:dimensional-tangent-dilation}
\lean{Preiss.dimensionalTangent_dilate}
\uses{def:dimensional-blowup, def:dimensional-tangent, lem:blowup-composition, lem:vague-moving-blowups, lem:dimensional-tangent-operations}




The source measure $\mu$ is locally finite.



For $r>0$ and $\nu\in\operatorname{Tan}_m(\mu,x)$, $D^m_{0,r}\nu\in\operatorname{Tan}_m(\mu,x)$.
\par\emph{Role: Preiss proof.}
\par Lean declarations: \texttt{Preiss.dimensionalTangent\_dilate}.
\end{lemma}
\begin{proof}[Planned proof obligations]

For a fixed measure and centre, the limit of a vaguely convergent sequence of
dimensional tangents is a dimensional tangent whenever that limit is regular
and nonzero. Prove the analogous arbitrary-scalar statement as a separate
adapter. Constant positive dilations of a dimensional tangent remain dimensional
tangents. At a positive finite density point the dimensional tangent set is
nonempty and compact, because every sufficiently small blow-up has compact
mass bounds and nonzero subsequential limits. A countable-test diagonal lemma
implements closure without assuming a topology on all raw measures.
\par Source: De Lellis Chapters 3 and 6; general supporting analysis.
\end{proof}

\begin{lemma}[Transitivity of the support-centered scaling orbit closure]
\label{lem:pseudotangent-transitive}
\lean{Preiss.IsUniformPseudotangent, Preiss.uniformPseudotangent_isUnitUniform, Preiss.uniformPseudotangent_trans, Preiss.uniformPseudotangent_map_linearIsometry, Preiss.uniformPseudotangent_trans_of_regular}
\uses{def:dimensional-blowup, lem:blowup-composition, lem:vague-diagonal, lem:vague-moving-blowups, lem:vague-support-approximation, lem:uniform-vague-closed, lem:uniform-scaling, lem:closed-map-support, lem:proper-map-regularity, lem:vague-cocompact-pushforward, lem:dimensional-blowup-support}






Let $\mathcal O_m(\nu)$ be the closure of
$\{D^m_{y,r}\nu:y\in\operatorname{supp}\nu,\ 0<r\leq1\}$ within regular
locally finite measures. For uniform measures all these limits retain their
ball constant. Prove
$\lambda\in\mathcal O_m(\nu)\Rightarrow
 \mathcal O_m(\lambda)\subseteq\mathcal O_m(\nu)$.
To compose a rescaling centred at a support point $z$ of $\lambda$, approximate
$z$ by support points of the convergent approximating measures (S04), then use
S02 and S03. Products of radii remain at most one. A diagonal selection turns
closure membership into the sequence form in the current Lean statement.
\par\emph{Role: Preiss proof.}
\par Lean declarations: \texttt{Preiss.IsUniformPseudotangent, Preiss.uniformPseudotangent\_isUnitUniform, Preiss.uniformPseudotangent\_map\_linearIsometry, Preiss.uniformPseudotangent\_trans}.
\end{lemma}
\begin{proof}[Planned proof obligations]

Let $\mathcal O_m(\nu)$ be the closure of
$\{D^m_{y,r}\nu:y\in\operatorname{supp}\nu,\ 0<r\leq1\}$ within regular
locally finite measures. For uniform measures all these limits retain their
ball constant. Prove
$\lambda\in\mathcal O_m(\nu)\Rightarrow
 \mathcal O_m(\lambda)\subseteq\mathcal O_m(\nu)$.
To compose a rescaling centred at a support point $z$ of $\lambda$, approximate
$z$ by support points of the convergent approximating measures (S04), then use
S02 and S03. Products of radii remain at most one. A diagonal selection turns
closure membership into the sequence form in the current Lean statement.
\par Source: De Lellis Lemma 6.6, Theorem 6.8, pp.57-58,63-65; Corollary 3.9 and corrected Lemma 3.8, pp.15,25-26; author errata.
\end{proof}

\begin{lemma}[Nonempty compact tangent family at a density point]
\label{lem:dimensional-tangent-compactness}
\lean{Preiss.dimensionalTangent_nonempty_compact}
\uses{lem:dimensional-family-compact-ball-bounds, gmt:density-bounds}







If $0<\lim_{r\downarrow0}\mu(B(x,r))/r^m<\infty$, the dimensional tangent family is nonempty and sequentially compact; every sequence of sufficiently small dimensional blow-ups has a subsequence with a nonzero regular limit.
\par\emph{Role: Preiss proof.}
\par Lean declarations: \texttt{Preiss.dimensionalTangent\_nonempty\_compact}.
\end{lemma}
\begin{proof}[Planned proof obligations]

For a fixed measure and centre, the limit of a vaguely convergent sequence of
dimensional tangents is a dimensional tangent whenever that limit is regular
and nonzero. Prove the analogous arbitrary-scalar statement as a separate
adapter. Constant positive dilations of a dimensional tangent remain dimensional
tangents. At a positive finite density point the dimensional tangent set is
nonempty and compact, because every sufficiently small blow-up has compact
mass bounds and nonzero subsequential limits. A countable-test diagonal lemma
implements closure without assuming a topology on all raw measures.
\par Source: De Lellis Chapters 3 and 6; general supporting analysis.
\end{proof}

\begin{lemma}[Comparison of tangent conventions]
\label{lem:tangent-ray}
\lean{Preiss.tangent_iff_dimensional_ray}
\uses{lem:tangent-ray-density-bounds, def:positive-density}











At a positive finite density point of a locally finite measure, an arbitrary
scalar tangent is exactly a positive finite scalar multiple of a dimensional
tangent. These two sets of tangents are not equal.
\end{lemma}
\begin{proof}[Planned proof]

Extract a nonzero dimensional limit using two-sided local ball bounds.
A compactly supported test with nonzero integral detects the relative scalar
limit; equality against all compact tests identifies the two limits up to
a positive finite scalar. This argument only needs density bounds.
\end{proof}

\begin{lemma}[Blow-downs of tangents remain tangents]
\label{lem:tangent-blowdown}
\lean{Preiss.dimensionalTangent_of_tangentAtInfinity}
\uses{def:dimensional-tangent, def:tangent-at-infinity, lem:dimensional-tangent-closed, lem:dimensional-tangent-dilation}




The source measure $\mu$ is locally finite.



If a dimensional tangent $\zeta$ has a nonzero regular limit at infinity $\chi$,
then $\chi$ is a dimensional tangent of the original measure at the same point.
Use S07's positive dilation invariance and closure. This statement is pointwise;
it needs no almost-everywhere recentering theorem.
\par\emph{Role: Preiss proof.}
\par Lean declarations: \texttt{Preiss.dimensionalTangent\_of\_tangentAtInfinity}.
\end{lemma}
\begin{proof}[Planned proof obligations]

If a dimensional tangent $\zeta$ has a nonzero regular limit at infinity $\chi$,
then $\chi$ is a dimensional tangent of the original measure at the same point.
Use S07's positive dilation invariance and closure. This statement is pointwise;
it needs no almost-everywhere recentering theorem.
\par Source: De Lellis Lemma 6.20 and Section 6.4, pp.66-69.
\end{proof}

\begin{lemma}[Continuity of the scale function]
\label{lem:flatness-radius-continuous}
\lean{Preiss.continuousOn_flatness_dimensionalBlowup}
\uses{def:dimensional-blowup, def:flatness, lem:flatness-continuous, lem:vague-moving-blowups, def:flatness-cutoff, lem:dimensional-blowup-support, lem:flatness-vague-general}




For fixed regular locally finite $\mu$ and center $x$, $r\mapsto F(D^m_{x,r}\mu)$ is continuous on $(0,\infty)$.
\par\emph{Role: Preiss proof.}
\par Lean declarations: \texttt{Preiss.continuousOn\_flatness\_dimensionalBlowup}.
\end{lemma}
\begin{proof}[Planned proof obligations]

Choose $\phi\in C_c(\mathbb R^n)$ with $0\leq\phi\leq1$,
$\phi=1$ on $\overline B(0,1)$ and support in $B(0,2)$.
Define $F(\mu)=\min_V\int\phi(z)\operatorname{dist}(z,V)^2\,d\mu$.
The minimum is attained. S01 and C01 imply sequential continuity for vague
convergence; S02 implies continuity of $r\mapsto F(D^m_{x,r}\mu)$ for $r>0$.
Flat measures have $F=0$. For a unit uniform measure, $F\leq\varepsilon$
implies that some plane has unit-ball energy at most $\varepsilon$.
\par Source: De Lellis Lemma 6.20 and Section 6.4, pp.66-69.
\end{proof}

\begin{lemma}[One null exceptional set for all dimensional recenterings]
\label{lem:recenter}
\lean{Preiss.ae_recenter_dimensionalTangent, Preiss.ae_recenter_dimensionalTangent_of_locallyFinite}
\uses{lem:vague-second-countable, lem:dimensional-tangent-metric-separation, lem:recenter-separated-set-null, lem:recenter-nonzero, lem:dimensional-tangent-operations}






Under positive finite existing $m$-density almost everywhere, prove
\[
 \text{for }\mu\text{-a.e. }a,\quad
 \forall\nu\in\operatorname{Tan}_m(\mu,a),\quad
 \forall y\in\operatorname{supp}\nu,\quad\forall r>0,
 \quad D^m_{y,r}\nu\in\operatorname{Tan}_m(\mu,a).
\]
Here $\operatorname{Tan}_m$ uses fixed dimensional normalization.
It is enough to prove recentering at radius one; S03 and S07 give all radii.
For fixed rational tolerances and radius cutoffs, take the arbitrary set of
bad centres, stratify its density constants and intersect with a bounded ball.
If it has positive outer measure, use a finite metric cover of the compact
family of translated uniform tangents to select one positive-outer-measure
class $D$ whose chosen translated tangents are pairwise close.
Take a density point of its measurable hull lying in $D$: the hull's bad
density points are null and cannot cover a set of positive outer measure.
Apply A08 to select nearby bad centres in $D$ and S02 to compare their blow-ups
with the chosen translated tangent. The triangle inequality contradicts the
defining fixed separation of the bad centres. Countably many tolerances give
one null exceptional set. This repairs the book's unjustified assertion that
its existentially defined set $D$ is measurable.
\par\emph{Role: Preiss proof.}
\par Lean declarations: \texttt{Preiss.ae\_recenter\_dimensionalTangent}.
\end{lemma}
\begin{proof}[Planned proof obligations]

Under positive finite existing $m$-density almost everywhere, prove
\[
 \text{for }\mu\text{-a.e. }a,\quad
 \forall\nu\in\operatorname{Tan}_m(\mu,a),\quad
 \forall y\in\operatorname{supp}\nu,\quad\forall r>0,
 \quad D^m_{y,r}\nu\in\operatorname{Tan}_m(\mu,a).
\]
Here $\operatorname{Tan}_m$ uses fixed dimensional normalization.
It is enough to prove recentering at radius one; S03 and S07 give all radii.
For fixed rational tolerances and radius cutoffs, take the arbitrary set of
bad centres, stratify its density constants and intersect with a bounded ball.
If it has positive outer measure, use a finite metric cover of the compact
family of translated uniform tangents to select one positive-outer-measure
class $D$ whose chosen translated tangents are pairwise close.
Take a density point of its measurable hull lying in $D$: the hull's bad
density points are null and cannot cover a set of positive outer measure.
Apply A08 to select nearby bad centres in $D$ and S02 to compare their blow-ups
with the chosen translated tangent. The triangle inequality contradicts the
defining fixed separation of the bad centres. Countably many tolerances give
one null exceptional set. This repairs the book's unjustified assertion that
its existentially defined set $D$ is measurable.
\par Corrections and boundary conditions: New measurable-hull proof replacing unjustified measurability in De Lellis Lemma 6.6; quantify all tangents, support centers and positive radii outside one null set.
\par Source: De Lellis Lemma 6.6, Theorem 6.8, pp.57-58,63-65; Corollary 3.9 and corrected Lemma 3.8, pp.15,25-26; author errata.
\end{proof}

\section{Weak approximation and the sufficient criterion}

\begin{definition}[Relaxed uniform approximation good sets]
\label{def:weak-approximation-good-set}
\lean{Preiss.weakApproximationGoodSet}
\uses{def:plane}








Parameters \(\eta,\lambda,R\) are positive. Later closedness results assume \(\nu\) finite and regular.
Let \(G(\nu,m,\eta,\lambda,R)\) consist of points \(x\) for which every \(0<r<R\) admits an \(m\)-plane \(V\) satisfying
\[
 \nu\{y:\|y-x\|<r,\ \operatorname{dist}(y-x,V)>\eta r\}\leq\eta r^m,
\]
and
\[
 \nu(\overline B_{\eta r}(z))\geq\lambda r^m
 \quad\text{for all }z\in(x+V)\cap B_r(x).
\]
\par\emph{Role: Preiss proof.}
\par Lean declarations: \texttt{Preiss.weakApproximationGoodSet}.
\end{definition}

\begin{lemma}[Small tube-error mass implies support containment]
\label{lem:tube-support-upgrade}
\lean{Preiss.subset_tube_of_lower_mass_and_small_outside_mass}
\uses{def:plane}

The measure $\mu$ and set $C$ are arbitrary. The centre $a\in\mathbb R^n$ need not lie in $C$. No finiteness or measurability hypothesis is required.
Let $\delta>0$, $0<\eta\leq1$, and $r>0$. If $\mu(B(z,\eta r/2))\geq\delta(\eta r/2)^m$ for every $z\in C$, and the mass in $B(a,2r)$ outside the closed $\eta r/2$ tube about $a+V$ is less than $\delta(\eta r/2)^m$, then $C\cap B(a,r)$ is contained in the closed $\eta r$ tube about $a+V$.
\par\emph{Role: Preiss proof.}
\par Lean declarations: \texttt{Preiss.subset\_tube\_of\_lower\_mass\_and\_small\_outside\_mass}.
\end{lemma}
\begin{proof}[Planned proof obligations]
\par Source: De Lellis Chapters 4--5; see the rectifiability audit.
\end{proof}

\begin{definition}[Weak linear approximation]
\label{def:weak-approximation}
\lean{Preiss.MeasureWeaklyLinearlyApproximableAt, Preiss.WeaklyLinearlyApproximableAt}
\uses{def:plane}










A measure $\mu$ is weakly linearly approximable at $x$ if for every $\eta>0$ there are $a,r_0>0$ such that, for every $0<r<r_0$, some $m$-plane $V$ satisfies both $\mu(B(x,r)\setminus\{z:\operatorname{dist}(z-x,V)\leq\eta r\})<\eta r^m$ and $\mu(B(z,\eta r))\geq ar^m$ for every $z\in(x+V)\cap B(x,r)$. The set predicate specializes this to $\mathcal H^m\!\restriction E$.
\end{definition}

\begin{lemma}[Small projection direction bounds projected tube measure]
\label{lem:projected-tube-slab-bound}
\lean{Preiss.hausdorffMeasure_projected_tube_le}
\uses{lem:plane-hausdorff, lem:plane-slab-measure, lem:projected-tube-contained-slab, lem:empty-cylinder-direction}







The subspaces \(V,W\) have dimension \(m>0\); \(\varepsilon,\eta\geq0\) and \(R>0\). The constant depends only on dimension and normalization.
There is \(C_m>0\) such that, if \(w\in W\) is a unit vector with \(\|P_Vw\|\leq\varepsilon\), then
\[
 \mathcal H^m\!\left(P_V\bigl(\{z:\operatorname{dist}(z,a+W)\leq\eta R\}\cap\overline B_R(a)\bigr)\right)
 \leq C_m(\varepsilon+\eta)R^m.
\]
\par\emph{Role: Preiss proof.}
\par Lean declarations: \texttt{Preiss.hausdorffMeasure\_projected\_tube\_le}.
\end{lemma}
\begin{proof}[Planned proof obligations]

\par Source: De Lellis Chapters 4--5; see the rectifiability audit.
\end{proof}

\begin{lemma}[Relaxed approximation good sets are closed]
\label{lem:closed-weak-approximation-good-set}
\lean{Preiss.isClosed_weakApproximationGoodSet}
\uses{def:weak-approximation-good-set, lem:grassmannian, lem:moving-set-measure-bounds, lem:projection-distance, lem:plane-projection-subsequence}








\(\nu\) is finite and regular; use the compact orthogonal-projection topology on \(m\)-planes and fixed positive \(\eta,\lambda,R\).
The relaxed good set \(G(\nu,m,\eta,\lambda,R)\) is closed.
For each fixed radius, the relation in \(x,V\) is closed; projection along the compact plane space and intersection over radii preserve closedness.
\par\emph{Role: Preiss proof.}
\par Lean declarations: \texttt{Preiss.isClosed\_weakApproximationGoodSet}.
\end{lemma}
\begin{proof}[Planned proof obligations]

Prove joint LSC moving-open-set mass and USC closed-ball mass, approximate centres via projections; project closed relation along compact plane witness.
\par Source: De Lellis Chapters 4--5; see the rectifiability audit.
\end{proof}

\begin{lemma}[Weak approximation gives positive lower density]
\label{lem:weak-approximation-lower-density}
\lean{Preiss.lower_density_pos_of_weakApproximation, Preiss.lower_power_bound_of_weakApproximation}
\uses{def:weak-approximation, lem:lower-density-from-power}








No finite-total-measure assumption is required for this implication. Its pure-measure version concludes a positive local lower power bound.
Weak linear approximation of \(F\) at \(x\) implies
\[
 \Theta_*^m(\mathcal H^m\!\restriction F,x)>0.
\]
Use the lower-ball clause at the plane centre \(z=x\) and one fixed aperture.
\par\emph{Role: Preiss proof.}
\par Lean declarations: \texttt{Preiss.lower\_density\_pos\_of\_weakApproximation, Preiss.lower\_power\_bound\_of\_weakApproximation}.
\end{lemma}
\begin{proof}[Planned proof obligations]
\par Source: De Lellis Chapters 4--5; see the rectifiability audit.
\end{proof}

\begin{lemma}[Weak approximation survives negligible deletion]
\label{lem:weak-approximation-deletion}
\lean{Preiss.negligibleDeletionAt, Preiss.weakApproximation_restrict_of_negligible_sdiff}
\uses{def:weak-approximation}








The pointwise theorem assumes the displayed limit. Shrink the source aperture and positive lower-mass constant.
If \(F\subset E\) and
\[
 r^{-m}\mathcal H^m((E\setminus F)\cap B_{2r}(x))\longrightarrow0,
\]
then weak linear approximation of \(E\) at \(x\) implies that of \(F\) at \(x\).
\par\emph{Role: Preiss proof.}
\par Lean declarations: \texttt{Preiss.negligibleDeletionAt, Preiss.weakApproximation\_restrict\_of\_negligible\_sdiff}.
\end{lemma}
\begin{proof}[Planned proof obligations]
\par Source: De Lellis Chapters 4--5; see the rectifiability audit.
\end{proof}

\begin{lemma}[Aperture slack gives countable good-set exhaustion]
\label{lem:weak-approximation-good-exhaustion}
\lean{Preiss.weaklyLinearlyApproximableAt_iff_goodSet_exhaustion, Preiss.weakApproximationGoodSet_mono}
\uses{def:weak-approximation, def:weak-approximation-good-set}

No finiteness or regularity of $\mu$ is assumed.
For any measure $\mu$ on $\mathbb R^n$, weak linear approximation at $x$ is equivalent to: for every $\eta>0$ there exists $N\in\mathbb N$ with $x\in G(\mu,m,\eta,1/(N+1),1/(N+1))$. The reverse implication uses the good-set condition at aperture $\eta/2$ to obtain strict upper error and open lower balls at aperture $\eta$.
\par\emph{Role: Preiss proof.}
\par Lean declarations: \texttt{Preiss.weaklyLinearlyApproximableAt\_iff\_goodSet\_exhaustion}.
\end{lemma}
\begin{proof}[Planned proof obligations]
\par Source: De Lellis Chapters 4--5; see the rectifiability audit.
\end{proof}

\begin{lemma}[Compact subsets with uniform approximation and density bounds]
\label{lem:uniform-compact-selection}
\lean{Preiss.exists_compact_uniform_weakApproximation_subset}
\uses{gmt:goodsets, gmt:hausdorff-regular, lem:closed-weak-approximation-good-set, lem:weak-approximation-good-exhaustion, lem:weak-approximation-lower-density, lem:compact-measurable-exhaustion}







The Borel set \(E\) has finite \(\mathcal H^m\)-measure and is weakly linearly approximable almost everywhere. Intersect Borel \(E\) with closed good sets before applying inner regularity.
One can choose compact \(C\subset E\) losing arbitrarily little \(\mathcal H^m\)-measure, on which the selected lower mass bounds and finitely many relaxed weak-approximation conditions have uniform positive parameters.
\par\emph{Role: Preiss proof.}
\par Lean declarations: \texttt{Preiss.exists\_compact\_uniform\_weakApproximation\_subset}.
\end{lemma}
\begin{proof}[Planned proof obligations]

\par Source: De Lellis Chapters 4--5; see the rectifiability audit.
\end{proof}

\begin{lemma}[Weak approximation restricts to Borel subsets almost everywhere]
\label{lem:weak-approximation-restrict-ae}
\lean{Preiss.ae_weaklyLinearlyApproximableAt_subset}
\uses{gmt:density-points, gmt:hausdorff-density-upper, lem:weak-approximation-deletion, mathlib:besicovitch-differentiation, lem:vanishing-density-negligible-deletion}







The set \(E\) is Borel and \(\mathcal H^m(E)<\infty\). Differentiation and finite upper density supply the negligible-deletion estimate.
If \(E\) is weakly linearly approximable at \(\mathcal H^m\)-almost every point of \(E\), then each Borel \(F\subset E\) is weakly linearly approximable at \(\mathcal H^m\)-almost every point of \(F\).
\par\emph{Role: Preiss proof.}
\par Lean declarations: \texttt{Preiss.ae\_weaklyLinearlyApproximableAt\_subset}.
\end{lemma}
\begin{proof}[Planned proof obligations]

\par Source: De Lellis Chapters 4--5; see the rectifiability audit.
\end{proof}

\begin{lemma}[Base point with upper mass and small discarded mass]
\label{lem:packing-base-point}
\lean{Preiss.ae_upperPowerBoundAt_restrict_hausdorff, Preiss.exists_upper_mass_small_sdiff_point}
\uses{gmt:hausdorff-density-outside, gmt:hausdorff-density-upper, lem:vanishing-density-scaled-balls, lem:quantitative-upper-density-bound}







The Borel sets \(G\subset E\) satisfy \(0<\mathcal H^m(G)\leq\mathcal H^m(E)<\infty\), and \(t>0\) is chosen before \(a,r_0\).
Choose \(a\in G\) and \(r_0>0\) so that, for \(0<r<r_0\),
\[
 \mathcal H^m(E\cap B_r(a))\leq C_m r^m,\qquad
 \mathcal H^m((E\setminus G)\cap B_{2r}(a))\leq t r^m.
\]
\par\emph{Role: Preiss proof.}
\par Lean declarations: \texttt{Preiss.ae\_upperPowerBoundAt\_restrict\_hausdorff, Preiss.exists\_upper\_mass\_small\_sdiff\_point}.
\end{lemma}
\begin{proof}[Planned proof obligations]

\par Source: De Lellis Chapters 4--5; see the rectifiability audit.
\end{proof}

\begin{lemma}[Transfer density bounds and flat tangents to finite pieces]
\label{lem:finite-piece-transfer}
\lean{Preiss.ae_density_bounds_and_flat_tangents_restrict_hausdorff}
\uses{lem:absolute-continuity-density-bounds, lem:absolute-continuity-tangent-cones}







The pieces are Borel, finite in \(\mathcal H^m\), and their restrictions are comparable with positive finite constants. Assume the density and flat-tangent hypotheses for locally finite \(\mu\).
For each finite piece \(S_j\), the measure \(\nu_j=\mathcal H^m\!\restriction S_j\) has positive lower and finite upper \(m\)-density almost everywhere. If all tangents of \(\mu\) are \(m\)-flat almost everywhere, the same holds for all tangents of \(\nu_j\) at \(\nu_j\)-almost every point.
\par\emph{Role: Preiss proof.}
\par Lean declarations: \texttt{Preiss.ae\_density\_bounds\_and\_flat\_tangents\_restrict\_hausdorff}.
\end{lemma}
\begin{proof}[Planned proof obligations]

\par Source: De Lellis Chapters 4--5; see the rectifiability audit.
\end{proof}

\begin{lemma}[Nested compact two-sided approximation]
\label{lem:nested-compact-approximation}
\lean{Preiss.exists_nested_compact_twoSided_approximation}
\uses{gmt:density-points, lem:tube-support-upgrade, lem:uniform-compact-selection, lem:weak-approximation-restrict-ae, lem:closed-upper-ball-bound-sets, def:two-sided-approximation, lem:compact-two-sided-general, lem:compact-measurable-exhaustion}







The Borel set \(E\) has positive finite \(\mathcal H^m\)-measure and is weakly linearly approximable almost everywhere. Choose the lower constant \(\delta\) first, then a sufficiently small \(\eta\), then the remaining constants.
Choose compact positive-measure sets \(G\subset F\subset E\) with uniform lower mass on \(F\). At each \(a\in G\) and each small scale \(r\), a plane \(V\) gives tube containment of \(F\cap B_{2r}(a)\), proximity of \((a+V)\cap B_r(a)\) to \(F\), and a positive lower bound for the \(E\)-mass of plane-centred \(\eta r\)-balls.
\par\emph{Role: Preiss proof.}
\par Lean declarations: \texttt{Preiss.TwoSidedApproximation, Preiss.exists\_nested\_compact\_twoSided\_approximation}.
\end{lemma}
\begin{proof}[Planned proof obligations]

Preserve parameter order δ then η then γ then t.
\par Source: De Lellis Chapters 4--5; see the rectifiability audit.
\end{proof}

\begin{lemma}[Pure sets have arbitrarily close points outside fixed cones]
\label{lem:pure-cone-escape}
\lean{Preiss.ae_exists_outside_cone_of_purelyUnrectifiable}
\uses{lem:cone-graph, lem:pure-basic}







The set \(C\) is purely \(m\)-unrectifiable. The exceptional null set may depend on \(V,\alpha\); simultaneous quantification over all planes is unnecessary.
For every fixed \(m\)-plane \(V\) and finite aperture \(\alpha\geq0\), at \(\mathcal H^m\)-almost every \(a\in C\), every \(r>0\) admits
\[
 b\in C\cap B_r(a)\setminus C(a,V,\alpha).
\]
\par\emph{Role: Preiss proof.}
\par Lean declarations: \texttt{Preiss.ae\_exists\_outside\_cone\_of\_purelyUnrectifiable}.
\end{lemma}
\begin{proof}[Planned proof obligations]

\par Source: De Lellis Chapters 4--5; see the rectifiability audit.
\end{proof}

\begin{lemma}[Disjoint projection-hole disks with large radius sum]
\label{lem:projection-hole-cover}
\lean{Preiss.HoleFamily, Preiss.exists_projection_hole_disk_family, Preiss.projected_distance_le_of_lower_mass, Preiss.exists_holeFamily_of_compact, Preiss.holeFamily_mass_lower_bound, Preiss.HoleFamily.disjoint_cylinders, Preiss.mem_half_ball_of_tube_and_projected_distance, Preiss.exists_nearest_projected_witness_in_half_ball, Preiss.projected_distance_le_of_affine_plane_lower_mass}
\uses{lem:nested-compact-approximation, lem:packing-base-point, lem:plane-hausdorff, mathlib:vitali-enlargement}







Use \(P_V(G\cap\overline B_{2\sigma}(a))\), with \(G\) compact and null projection. Plane-centred mass and the small discarded-mass bound give the upper bound on hole radii.
There is a countable family of hole centres \(p_i\) and radii \(\rho_i\) such that the disks \(D_{4\rho_i}(p_i)\) are disjoint, their fivefold enlargements cover the relevant plane disk up to a null set, and
\[
 0<\rho_i\leq\eta\sigma,\qquad \sum_i\rho_i^m\geq c_m\sigma^m.
\]
Each radius is distance to the compact projected set.
\par\emph{Role: Preiss proof.}
\par Lean declarations: \texttt{Preiss.HoleFamily, Preiss.exists\_projection\_hole\_disk\_family, Preiss.projected\_distance\_le\_of\_lower\_mass}.
\end{lemma}
\begin{proof}[Planned proof obligations]

Reuse Vitali.exists\_disjoint\_subfamily\_covering\_enlargement\_ball within the plane.
\par Source: De Lellis Chapters 4--5; see the rectifiability audit.
\end{proof}

\begin{lemma}[Pure weakly approximable sets have every-plane null projection]
\label{lem:pure-wla-null-projection}
\lean{Preiss.measure_projection_eq_zero_of_purelyUnrectifiable_weakApproximation}
\uses{lem:fixed-lower-density-projection-null, lem:closed-weak-approximation-good-set, lem:weak-approximation-good-exhaustion}







The Borel set \(E\) is purely \(m\)-unrectifiable, has finite \(\mathcal H^m\)-measure, and is weakly linearly approximable almost everywhere; \(0<m\leq n\).
For every \(m\)-plane \(V\),
\[
 \mathcal H^m(P_V(E))=0.
\]
This is an every-plane conclusion, not an almost-every-plane projection theorem.
\par\emph{Role: Preiss proof.}
\par Lean declarations: \texttt{Preiss.measure\_projection\_eq\_zero\_of\_purelyUnrectifiable\_weakApproximation}.
\end{lemma}
\begin{proof}[Planned proof obligations]

Also reuse Besicovitch.exists\_disjoint\_closedBall\_covering\_ae and Lipschitz distortion; a fivefold covering alone is not identical.
\par Source: De Lellis Chapters 4--5; see the rectifiability audit.
\end{proof}

\begin{lemma}[Cylinders meeting the compact core have small radius sum]
\label{lem:discarded-cylinder-sum}
\lean{Preiss.tsum_hole_radii_bad_indices_le, Preiss.HoleFamily.exists_disjoint_charging_balls, Preiss.tsum_sdiff_ge_of_small_discarded_sum}
\uses{lem:projection-hole-cover}







Use the selected disjoint cylinders, uniform lower mass \(\delta\) on \(F\), and small mass of \(E\setminus G\). Choose \(t\) after \(\delta\).
For indices \(J\) whose inner cylinders meet \(F\),
\[
 \sum_{i\in J}\rho_i^m\leq C_m t\sigma^m/\delta.
\]
For sufficiently small \(t\), the complementary indices \(K\) retain
\(\sum_{i\in K}\rho_i^m\geq c_m\sigma^m\).
\par\emph{Role: Preiss proof.}
\par Lean declarations: \texttt{Preiss.tsum\_hole\_radii\_bad\_indices\_le}.
\end{lemma}
\begin{proof}[Planned proof obligations]

\par Source: De Lellis Chapters 4--5; see the rectifiability audit.
\end{proof}

\begin{lemma}[Many disjoint lower-mass balls inside each surviving cylinder]
\label{lem:column-packing}
\lean{Preiss.exists_column_disjoint_balls}
\uses{lem:empty-cylinder-direction, lem:nested-compact-approximation, lem:projection-hole-cover}







Use an empty inner cylinder, a compact nearest projected boundary point, two-sided approximation at that point, and sufficiently small \(\eta>0\). Never identify \(\sigma\) with \(\rho_i/(8\eta)\).
At the local scale \(r_i=\rho_i/(8\eta)\), each surviving cylinder contains at least \(c/\eta\) separated plane points and corresponding nearby centres in \(F\). The resulting balls of radius \(\rho_i/8\) are disjoint, lie inside the \(4\rho_i\)-cylinder and the ambient \(B_\sigma(a)\), and satisfy the prescribed lower mass bound.
\par\emph{Role: Preiss proof.}
\par Lean declarations: \texttt{Preiss.exists\_column\_disjoint\_balls}.
\end{lemma}
\begin{proof}[Planned proof obligations]

\par Source: De Lellis Chapters 4--5; see the rectifiability audit.
\end{proof}

\begin{lemma}[Null projections contradict weak approximation and finite upper density]
\label{lem:packing-contradiction}
\lean{Preiss.false_of_null_projections_and_uniform_weakApproximation}
\uses{lem:two-sided-packing-core, lem:nested-compact-approximation, lem:packing-base-point}







Use null projections, nested compact approximation, the selected hole disks and the surviving radius sum. Parameter order is \(\delta\), then \(\eta\ll\delta\), then the other approximation constants and \(t\), and finally \(\sigma\).
Summing the disjoint column balls gives
\[
 \mathcal H^m(E\cap B_\sigma(a))
 \geq c_m\delta\sigma^m/\eta,
\]
contradicting the upper bound \(C_m\sigma^m\) when \(\eta\) is sufficiently small.
\par\emph{Role: Preiss proof.}
\par Lean declarations: \texttt{Preiss.false\_of\_column\_mass\_exceeds\_upper\_bound, Preiss.false\_of\_null\_projections\_and\_uniform\_weakApproximation}.
\end{lemma}
\begin{proof}[Planned proof obligations]

\par Source: De Lellis Chapters 4--5; see the rectifiability audit.
\end{proof}

\begin{theorem}[Marstrand-Mattila set criterion]
\label{thm:mm-set}
\lean{Preiss.countablyRectifiable_of_ae_weaklyLinearlyApproximableAt}
\uses{lem:decomposition, lem:weak-approximation-restrict-ae, lem:packing-contradiction, lem:pure-wla-null-projection, lem:packing-base-point, lem:carrier-closure}







Assume \(0<m\leq n\). The zero-measure case is included and is immediate.
A Borel set \(E\subset\mathbb R^n\) with finite \(\mathcal H^m\)-measure which is weakly linearly approximable at \(\mathcal H^m\)-almost every point is countably \(m\)-rectifiable.
\par\emph{Role: Preiss proof.}
\par Lean declarations: \texttt{Preiss.countablyRectifiable\_of\_ae\_weaklyLinearlyApproximableAt}.
\end{theorem}
\begin{proof}[Planned proof obligations]

\par Source: De Lellis Chapters 4--5; see the rectifiability audit.
\end{proof}

\begin{theorem}[Marstrand-Mattila measure criterion]
\label{thm:mm-measure}
\lean{Preiss.isRectifiableMeasure_of_ae_density_bounds_and_flat_tangents}
\uses{lem:finite-pieces, lem:carrier-calculus, lem:finite-piece-transfer, lem:flat-to-approximation, thm:mm-set}







The measure \(\mu\) is locally finite and \(0<m\leq n\). Quantification over all tangents lies inside the almost-everywhere binder.
If, at \(\mu\)-almost every \(x\),
\[
 0<\Theta_*^m(\mu,x)\leq\Theta^{*m}(\mu,x)<\infty
\]
and every tangent measure is \(m\)-flat, then \(\mu\) is \(m\)-rectifiable.
\par\emph{Role: Preiss proof.}
\par Lean declarations: \texttt{Preiss.isRectifiableMeasure\_of\_ae\_density\_bounds\_and\_flat\_tangents}.
\end{theorem}
\begin{proof}[Planned proof obligations]

\par Source: De Lellis Chapters 4--5; see the rectifiability audit.
\end{proof}

\section{Generic Gaussian and multilinear analysis}

\begin{lemma}[Gaussian polynomial integrability under polynomial growth]
\label{lem:gaussian-polynomial-integrable}
\lean{Preiss.integrable_norm_pow_mul_gaussian_of_ball_bound, Preiss.integrable_norm_pow_mul_gaussian_of_polynomial_growth, Preiss.integral_gaussian_pos_of_ball_bound}
\uses{def:polynomial-growth-bound, lem:annulus-integrability, lem:gaussian-annulus-majorant}



For a regular measure with $\mu(B(0,R))\leq C(1+R)^p$ for all $R>0$, $C,p\geq0$, every $d\in\mathbb N$ and $s>0$ satisfy $\int |z|^d e^{-s|z|^2}\,d\mu<\infty$. If $\mu\ne0$, its Gaussian mass is positive.
\par\emph{Role: Preiss proof.}
\par Lean declarations: \texttt{Preiss.integrable\_norm\_pow\_mul\_gaussian\_of\_polynomial\_growth}.
\end{lemma}
\begin{proof}[Planned proof obligations]
For every $d\in\mathbb N$ and $s>0$,
$\int |z|^d e^{-s|z|^2}\,d\nu<\infty$ and $0<I_\nu(s)<\infty$.
For fixed $s,d$, the tails outside $B(0,R)$ tend to zero uniformly over
$\mathcal U_m$. State the more reusable version for a family satisfying one
common polynomial bound $\mu(B(0,R))\leq C(1+R)^p$.
Use this to extend vague convergence to polynomial-times-Gaussian tests.
It also justifies differentiation with respect to positive $s$ locally in $s$.
\par Source: De Lellis Lemma 7.2, Definition 7.4, Lemma 7.6, Proposition 7.7, pp.70-84.
\end{proof}

\begin{lemma}[Analytic extension of weighted complementary projections]
\label{lem:graded-analytic-projection}
\uses{def:orthogonal-grading, lem:kernel-flag-divided-operators, lem:divided-operator-positive-inverse, lem:weighted-grading-positive-complement}
\lean{Preiss.analyticAt_weightedProjection_zero}


This is a measure-independent linear algebra theorem. Let the finite-dimensional
real inner-product space $X$ have an orthogonal decomposition
$X=\bigoplus_{j=1}^{N}X_j$, with projections $P_j$, and let $V\leq X$.
For $s>0$ put $A(s)=\sum_j s^jP_j$ and
$F_s=A(s)^{-1}(V^\perp)$. Then $X=V\oplus F_s$; the projection $Q_s$
onto $V$ along $F_s$ extends real-analytically to a neighborhood of zero.
In particular its operator norm stays bounded as $s\downarrow0$.
For the proof, orthogonally split $V^\perp$ according to the filtration
$V^\perp\cap\bigoplus_{j\leq k}X_j$ and define a polynomial family which
is identity on $V$ and equals $P_k+sP_{k-1}+\cdots+s^{k-1}P_1$ on the
$k$th filtration summand. Its value at zero is invertible, by the
least-nonzero-degree argument, and $Q_s=P_V A_s^{-1}$ for that polynomial
family. In the least-degree argument the lower-degree vector belongs to the
sum of all earlier filtration summands, not necessarily only the immediately
preceding one.
\par\emph{Role: Preiss proof.}
\par Lean declarations: \texttt{Preiss.IsOrthogonalGrading, Preiss.analyticAt\_weightedProjection\_zero, Preiss.weightedComplement, Preiss.weightedGrading}.
\end{lemma}
\begin{proof}[Planned proof obligations]
This is a measure-independent linear algebra theorem. Let the finite-dimensional
real inner-product space $X$ have an orthogonal decomposition
$X=\bigoplus_{j=1}^{N}X_j$, with projections $P_j$, and let $V\leq X$.
For $s>0$ put $A(s)=\sum_j s^jP_j$ and
$F_s=A(s)^{-1}(V^\perp)$. Then $X=V\oplus F_s$; the projection $Q_s$
onto $V$ along $F_s$ extends real-analytically to a neighborhood of zero.
In particular its operator norm stays bounded as $s\downarrow0$.
For the proof, orthogonally split $V^\perp$ according to the filtration
$V^\perp\cap\bigoplus_{j\leq k}X_j$ and define a polynomial family which
is identity on $V$ and equals $P_k+sP_{k-1}+\cdots+s^{k-1}P_1$ on the
$k$th filtration summand. Its value at zero is invertible, by the
least-nonzero-degree argument, and $Q_s=P_V A_s^{-1}$ for that polynomial
family. In the least-degree argument the lower-degree vector belongs to the
sum of all earlier filtration summands, not necessarily only the immediately
preceding one.
\par Corrections and boundary conditions: Generic finite-dimensional replacement for the source tensor argument; filtration decomposition must include all earlier summands.
\par Source: De Lellis Lemma 7.2, Definition 7.4, Lemma 7.6, Proposition 7.7, pp.70-84.
\end{proof}

\begin{lemma}[Greatest crossing with an interval bound]
\label{lem:greatest-crossing}
\lean{Preiss.exists_greatest_crossing}



If $0<a<b$, $f:[a,b]\to\mathbb R$ is continuous and $f(a)>\varepsilon>f(b)$, there is $c\in(a,b)$ such that $f(c)=\varepsilon$ and $f(t)\leq\varepsilon$ for all $t\in[c,b]$.
\par\emph{Role: Preiss proof.}
\par Lean declarations: \texttt{Preiss.exists\_greatest\_crossing}.
\end{lemma}
\begin{proof}[Planned proof obligations]
After S08 normalization, suppose $\alpha$ is a flat dimensional tangent and
$\zeta$ is a nonflat dimensional tangent, all dimensional tangents being unit
uniform. Let $\chi$ be the blow-down of $\zeta$. Rigidity and the gap give
$F(\chi)>\varepsilon$. Choose $r_i\to0$ with blow-ups tending to $\alpha$
and $s_i<r_i$ with blow-ups tending to $\chi$; subsequence selection makes
the ordering possible. For sufficiently large $i$, the continuous scale
function satisfies $F(D_{x,r_i}^m\mu)<\varepsilon$ and
$F(D_{x,s_i}^m\mu)>\varepsilon$.
Choose $\sigma_i\in[s_i,r_i]$ with energy exactly $\varepsilon$ and energy
at most $\varepsilon$ on the entire interval $[\sigma_i,r_i]$.
This is a greatest-crossing/last-crossing assertion, not merely the intermediate
value theorem. Extract a nonzero unit uniform limit $\xi$ at the $\sigma_i$.
\par Corrections and boundary conditions: The ordinary intermediate value theorem alone does not give the full later interval bound.
\par Source: De Lellis Lemma 6.20 and Section 6.4, pp.66-69.
\end{proof}

\begin{definition}[Finite Euclidean polynomial feature map]
\label{def:polynomial-feature}
\uses{}
\lean{Preiss.PolynomialFeatureIndex, Preiss.PolynomialFeatureSpace, Preiss.polynomialFeature, Preiss.polynomialFeatureBlock, Preiss.polynomialFeatureProjection}


For $1\leq k\leq N$, index a Euclidean block by all $\iota:\operatorname{Fin}k\to\operatorname{Fin}n$ and set $\Phi_k(x)_\iota=\sqrt{2^k/k!}\prod_j x_{\iota(j)}$. The block pairing is $\langle\Phi_k(x),\Phi_k(y)\rangle=(2^k/k!)\langle x,y\rangle^k$. The direct sum has orthogonal coordinate projections $P_k$.
\par\emph{Role: Preiss proof.}
\par Lean declarations: \texttt{Preiss.PolynomialFeatureIndex, Preiss.PolynomialFeatureSpace, Preiss.polynomialFeature, Preiss.polynomialFeatureBlock, Preiss.polynomialFeatureProjection}.
\end{definition}

\begin{lemma}[Top eigenspaces minimize normal covariance trace]
\label{lem:spectral-minimizing-plane}
\lean{Preiss.trace_orthogonalComplement_minimized_topEigenspace}
\uses{def:plane}




For a positive semidefinite symmetric form $B$ on $\mathbb R^n$, choose an
orthonormal eigenbasis with $\alpha_1\geq\cdots\geq\alpha_n\geq0$.
The span $V$ of the first $m$ eigenvectors minimizes
$\operatorname{tr}(B|W^\perp)$ over $m$-planes $W$.
For G02 this trace is a fixed positive multiple of the Gaussian distance
energy. This is a reusable finite-dimensional variational spectral lemma.
\par\emph{Role: Preiss proof.}
\par Lean declarations: \texttt{Preiss.trace\_orthogonalComplement\_minimized\_topEigenspace}.
\end{lemma}
\begin{proof}[Planned proof obligations]
For a positive semidefinite symmetric form $B$ on $\mathbb R^n$, choose an
orthonormal eigenbasis with $\alpha_1\geq\cdots\geq\alpha_n\geq0$.
The span $V$ of the first $m$ eigenvectors minimizes
$\operatorname{tr}(B|W^\perp)$ over $m$-planes $W$.
For G02 this trace is a fixed positive multiple of the Gaussian distance
energy. This is a reusable finite-dimensional variational spectral lemma.
\par Source: De Lellis Corollary 8.4, Lemmas 8.6-8.7, Proposition 8.5, pp.85-94; author errata for the dimension-two shortcut.
\end{proof}

\begin{lemma}[Gaussian weighted polynomials determine growing measures]
\label{lem:moment-determination}
\lean{Preiss.measure_eq_of_gaussian_polynomial_moments_eq}
\uses{lem:gaussian-polynomial-integrable, def:polynomial-growth-bound, lem:complex-mgf-moment-determination, lem:gaussian-weight-finite, lem:gaussian-weight-exponential}
Regular measures with a common polynomial growth bound are equal if their
$e^{-|z|^2}$-weighted polynomial moments agree.
\par\emph{Role: Preiss proof.}
\end{lemma}
\begin{proof}[Proof sketch]
Weight the two measures by $e^{-\|z\|^2}$. Polynomial growth makes the weighted measures finite. For every vector $w$ and every real $t$, the inequality
\[
t\langle z,w\rangle-\|z\|^2\leq |t|^2\|w\|^2/2-\|z\|^2/2
\]
gives integrability of the exponential of the projection under each weighted measure. The assumed polynomial equalities give equality of all natural moments of every projection. The entire complex moment generating functions have identical derivatives at zero, so analytic uniqueness identifies them. Evaluating at the imaginary unit identifies the characteristic functions of the finite weighted measures. Mathlib's characteristic-function uniqueness theorem identifies these measures, and cancellation of the strictly positive finite Gaussian density identifies the original measures. This avoids expanding a Gaussian at an arbitrary rate without domination.
\par Source: De Lellis, Chapters 7--8, with a uniqueness argument through finite Gaussian weighting and Mathlib's analytic moment-generating-function and characteristic-function APIs.
\end{proof}

\begin{lemma}[Uniform Gaussian polynomial tails]
\label{lem:gaussian-uniform-tails}
\lean{Preiss.exists_uniform_gaussian_tail_radius, Preiss.gaussian_polynomial_uniform_tail}
\uses{def:polynomial-growth-bound, lem:gaussian-polynomial-integrable, lem:gaussian-tail-series-bound}







For fixed $C,p\geq0$, $d\in\mathbb N$, $s>0$, the supremum of $\int_{|z|>R}|z|^d e^{-s|z|^2}\,d\mu$ over all regular measures satisfying $\mu(B(0,t))\leq C(1+t)^p$ tends to zero as $R\to\infty$.
\par\emph{Role: Preiss proof.}
\par Lean declarations: \texttt{Preiss.gaussian\_polynomial\_uniform\_tail}.
\end{lemma}
\begin{proof}[Planned proof obligations]

For every $d\in\mathbb N$ and $s>0$,
$\int |z|^d e^{-s|z|^2}\,d\nu<\infty$ and $0<I_\nu(s)<\infty$.
For fixed $s,d$, the tails outside $B(0,R)$ tend to zero uniformly over
$\mathcal U_m$. State the more reusable version for a family satisfying one
common polynomial bound $\mu(B(0,R))\leq C(1+R)^p$.
Use this to extend vague convergence to polynomial-times-Gaussian tests.
It also justifies differentiation with respect to positive $s$ locally in $s$.
\par Corrections and boundary conditions: Explicit uniform tail control is required to pass unbounded Gaussian tests through vague limits.
\par Source: De Lellis Lemma 7.2, Definition 7.4, Lemma 7.6, Proposition 7.7, pp.70-84.
\end{proof}

\begin{lemma}[Differentiation of Gaussian polynomial integrals]
\label{lem:gaussian-parameter-derivative}
\lean{Preiss.hasDerivAt_integral_polynomial_mul_gaussian, Preiss.hasDerivAt_integral_polynomial_gaussian_of_ball_bound}
\uses{def:polynomial-growth-bound, lem:gaussian-polynomial-integrable, lem:polynomial-gaussian-integrability}







Under polynomial growth, $s\mapsto\int P(z)e^{-s|z|^2}\,d\mu$ is differentiable for $s>0$ with derivative $-\int |z|^2P(z)e^{-s|z|^2}\,d\mu$; the identity iterates, with domination uniform on compact positive parameter intervals.
\par\emph{Role: Preiss proof.}
\par Lean declarations: \texttt{Preiss.hasDerivAt\_integral\_polynomial\_mul\_gaussian}.
\end{lemma}
\begin{proof}[Planned proof obligations]

For every $d\in\mathbb N$ and $s>0$,
$\int |z|^d e^{-s|z|^2}\,d\nu<\infty$ and $0<I_\nu(s)<\infty$.
For fixed $s,d$, the tails outside $B(0,R)$ tend to zero uniformly over
$\mathcal U_m$. State the more reusable version for a family satisfying one
common polynomial bound $\mu(B(0,R))\leq C(1+R)^p$.
Use this to extend vague convergence to polynomial-times-Gaussian tests.
It also justifies differentiation with respect to positive $s$ locally in $s$.
\par Source: De Lellis Lemma 7.2, Definition 7.4, Lemma 7.6, Proposition 7.7, pp.70-84.
\end{proof}

\begin{definition}[Countable-test vague metric specification]
\label{def:vague-metrization}
\lean{Preiss.RegularEuclideanMeasure, Preiss.vagueTestDistance, Preiss.VagueMetrization}
\uses{gmt:vague}
A regular Euclidean measure is a positive Borel measure equipped with regularity. A vague metrization consists of compactly supported continuous tests $f_j$, a metric on these measures whose distance is $\sum_{j\geq0}2^{-(j+1)}\min(1,|\int f_j\,d\mu-\int f_j\,d\nu|)$, cutoff tests majorizing $1$ on each positive-radius closed ball and nonnegative everywhere, and equivalence between metric sequential convergence and vague convergence for every sequence of regular measures.
\end{definition}

\begin{lemma}[Existence of a countable-test vague metric]
\label{lem:vague-countable-metric}
\lean{Preiss.exists_vagueMetrization}
\uses{def:vague-metrization, lem:countable-coordinate-metric, lem:dense-compact-tests-exist, def:vague-metrization-from-tests}
For every natural ambient dimension $n$, there exists a countable-test vague metrization with the preceding specification.
\end{lemma}
\begin{proof}[Planned construction]
Choose countably many uniformly dense tests on compact support strata and add nonnegative cutoff tests controlling every bounded ball. Prove that the weighted truncated distance separates regular measures, satisfies the metric axioms, and detects vague convergence. The cutoff coordinates supply uniform compact mass bounds when passing from dense tests to arbitrary compactly supported continuous tests. This existence proof remains a separate obligation.
\end{proof}

\begin{lemma}[Consequences of a supplied vague metric]
\label{lem:vague-metric-properties}
\lean{Preiss.vagueMetric_tendsto_iff, Preiss.vagueMetric_separates, Preiss.vagueMetric_finite_cover_of_compact}
\uses{def:vague-metrization}
For supplied $M:\mathrm{VagueMetrization}(n)$, metric convergence is equivalent to vague convergence; its countable-test distance separates measures; and every compact family has a finite cover by metric balls of every positive radius, with centers in that family.
\end{lemma}
\begin{proof}[Proof sketch]
Use the supplied convergence and distance fields, metric separation, and finite ball covers of compact sets. None of these arguments constructs a metrization.
\end{proof}


\begin{lemma}[Measurable hulls preserve disjoint holes]
\label{lem:measurable-hull-holes}
\lean{Preiss.measure_toMeasurable_inter_eq_zero_of_disjoint}
\uses{mathlib:measurable-hull}







For arbitrary $D$ with $\mu D<\infty$, put $H=\operatorname{toMeasurable}_\mu D$.
Then for every measurable $U$ disjoint from $D$, $\mu(H\cap U)=0$.
This is an immediate use of mathlib's
\texttt{Measure.measure\_toMeasurable\_inter}; an s-finite version is also
available. In particular $\mu(H\setminus\overline D)=0$.
This statement requires no measurability of $D$.
\par\emph{Role: Preiss proof.}
\par Lean declarations: \texttt{Preiss.measure\_toMeasurable\_inter\_eq\_zero\_of\_disjoint}.
\end{lemma}
\begin{proof}[Planned proof obligations]

For arbitrary $D$ with $\mu D<\infty$, put $H=\operatorname{toMeasurable}_\mu D$.
Then for every measurable $U$ disjoint from $D$, $\mu(H\cap U)=0$.
This is an immediate use of mathlib's
\texttt{Measure.measure\_toMeasurable\_inter}; an s-finite version is also
available. In particular $\mu(H\setminus\overline D)=0$.
This statement requires no measurability of $D$.
\par Corrections and boundary conditions: New proof repair: arbitrary sets of bad centers are not asserted Borel. Use finite-outer-mass localization and Measure.measure\_toMeasurable\_inter.
\par Source: De Lellis Lemma 6.6, Theorem 6.8, pp.57-58,63-65; Corollary 3.9 and corrected Lemma 3.8, pp.15,25-26; author errata.
\end{proof}

\begin{lemma}[Wick fourth moment of linear Gaussian tests]
\label{lem:gaussian-wick-fourth}
\lean{Preiss.gaussian_integral_inner_pow_even, Preiss.gaussian_integral_inner_sq_mul_inner_sq, Preiss.gaussian_integral_four_inner}
\uses{mathlib:standard-gaussian, lem:real-gaussian-even-moments, lem:gaussian-projection-integrability}





For vectors $y,z$ in a finite-dimensional Euclidean space with standard covariance-one Gaussian probability $\gamma$, $\int\langle y,v\rangle^2\langle z,v\rangle^2\,d\gamma(v)=|y|^2|z|^2+2\langle y,z\rangle^2$.
\par\emph{Role: Preiss proof.}
\par Lean declarations: \texttt{Preiss.gaussian\_integral\_inner\_pow\_even, Preiss.gaussian\_integral\_inner\_sq\_mul\_inner\_sq}.
\end{lemma}
\begin{proof}[Planned proof obligations]

For $y,z\in V$,
$\int\langle y,v\rangle^2\langle z,v\rangle^2\,d\gamma
=|y|^2|z|^2+2\langle y,z\rangle^2$.
Combining this with R07 yields, for $w\in V^\perp$,
\[
 \int\langle y,v\rangle^2\langle\omega(v,v),w\rangle\,d\gamma(v)
 =2\langle\omega(y,y),w\rangle.
\]
Justify limits by Gaussian quartic domination and Fubini.
\par Source: De Lellis Proposition 9.1, Lemmas 9.2,9.4,9.9,9.10 and Proposition 9.5, pp.95-109.
\end{proof}

\begin{definition}[Bilinear trace and restriction]
\label{def:bilinear-trace}
\lean{Preiss.bilinearTrace, Preiss.restrictBilinear, Preiss.momentBilinear}


For a continuous bilinear form on a finite-dimensional real inner product space, take the trace of its Riesz-represented linear map. Restriction to a subspace and conversion of a two-linear moment form to a bilinear form are concrete data operations.
\end{definition}

\begin{lemma}[Approximation of support points under vague convergence]
\label{lem:vague-support-approximation}
\lean{Preiss.exists_support_sequence_of_vague}
\uses{gmt:portmanteau}







If regular $\mu_i\rightharpoonup\mu$ and $x\in\operatorname{supp}\mu$, there are, for all sufficiently large $i$, $x_i\in\operatorname{supp}\mu_i$ with $x_i\to x$.
\par\emph{Role: Preiss proof.}
\par Lean declarations: \texttt{Preiss.exists\_support\_sequence\_of\_vague}.
\end{lemma}
\begin{proof}[Planned proof obligations]

If regular $\mu_i\rightharpoonup\mu$ and $x\in\operatorname{supp}\mu$, then
there exist $x_i\in\operatorname{supp}\mu_i$ with $x_i\to x$, after discarding
finitely many terms if necessary. Use the open-ball portmanteau inequality.
Conversely, if $x_i\to x$, $x_i\in\operatorname{supp}\mu_i$, and the measures
have a common positive lower support-centred ball bound at each radius, then
$x\in\operatorname{supp}\mu$. Bare vague convergence does not give this converse.
\par Source: De Lellis Chapters 3 and 6; general supporting analysis.
\end{proof}

\begin{lemma}[Limits of support points under lower ball bounds]
\label{lem:vague-support-limit}
\lean{Preiss.mem_support_of_vague_lower_ball_bound}
\uses{gmt:portmanteau}







If regular $\mu_i\rightharpoonup\mu$, $x_i\in\operatorname{supp}\mu_i$, $x_i\to x$, and for every $r>0$ there is $c_r>0$ with $\mu_i(B(x_i,r))\geq c_r$ for every sufficiently large $i$, then $x\in\operatorname{supp}\mu$.
\par\emph{Role: Preiss proof.}
\par Lean declarations: \texttt{Preiss.mem\_support\_of\_vague\_lower\_ball\_bound}.
\end{lemma}
\begin{proof}[Planned proof obligations]

If regular $\mu_i\rightharpoonup\mu$ and $x\in\operatorname{supp}\mu$, then
there exist $x_i\in\operatorname{supp}\mu_i$ with $x_i\to x$, after discarding
finitely many terms if necessary. Use the open-ball portmanteau inequality.
Conversely, if $x_i\to x$, $x_i\in\operatorname{supp}\mu_i$, and the measures
have a common positive lower support-centred ball bound at each radius, then
$x\in\operatorname{supp}\mu$. Bare vague convergence does not give this converse.
\par Corrections and boundary conditions: The positive lower ball hypothesis is necessary; bare vague convergence does not preserve support.
\par Source: De Lellis Chapters 3 and 6; general supporting analysis.
\end{proof}

\begin{lemma}[Vague convergence with uniformly moving compact tests]
\label{lem:vague-moving-tests}
\lean{Preiss.vague_integral_moving_compact_test}
\uses{gmt:compact-test-perturbation, gmt:portmanteau, gmt:vague}








If $\mu_i\rightharpoonup\mu$ vaguely, all measures are regular, $K$ is compact,
the continuous tests $f_i,f$ vanish outside $K$, and $f_i\to f$ uniformly,
then $\int f_i\,d\mu_i\to\int f\,d\mu$. First derive a uniform mass bound on
$K$ using a compactly supported cutoff equal to one on $K$.
\par\emph{Role: Preiss proof.}
\par Lean declarations: \texttt{Preiss.vague\_integral\_moving\_compact\_test}.
\end{lemma}
\begin{proof}[Planned proof obligations]

If $\mu_i\rightharpoonup\mu$ vaguely, all measures are regular, $K$ is compact,
the continuous tests $f_i,f$ vanish outside $K$, and $f_i\to f$ uniformly,
then $\int f_i\,d\mu_i\to\int f\,d\mu$. First derive a uniform mass bound on
$K$ using a compactly supported cutoff equal to one on $K$.
\par Source: De Lellis Chapters 3 and 6; general supporting analysis.
\end{proof}

\begin{lemma}[Diagonal selection for vague convergence]
\label{lem:vague-diagonal}
\lean{Preiss.vague_diagonal_selection}
\uses{lem:vague-countable-metric}






For a countable determining compact-test family, if $\mu_{i,j}\rightharpoonup\mu_i$ as $j\to\infty$ and $\mu_i\rightharpoonup\mu$, select $j(i)\geq i$ such that $\mu_{i,j(i)}\rightharpoonup\mu$; additional finitely many eventual constraints at stage $i$ can be imposed. Assume regularity and the compact mass bounds needed by the determining family.
\par\emph{Role: Preiss proof.}
\par Lean declarations: \texttt{Preiss.vague\_diagonal\_selection}.
\end{lemma}
\begin{proof}[Planned proof obligations]

Construct or reuse a metric on regular positive locally finite Euclidean
measures which detects vague convergence on families with common compact mass
bounds. Include enough cutoff tests to control masses. The necessary outputs are
sequential convergence equivalence, separation, a diagonal selection lemma,
and finite covers of a compact family by metric balls. This can live in a
general vague-convergence module rather than the Preiss theorem file.
\par Source: De Lellis Chapters 3 and 6; general supporting analysis.
\end{proof}

\begin{lemma}[Trace as an orthonormal diagonal sum]
\label{lem:bilinear-trace-basis}
\lean{Preiss.bilinearTrace_eq_sum_diagonal}
\uses{def:bilinear-trace}


The trace of a continuous bilinear form is the sum of its diagonal evaluations on any orthonormal basis.
\end{lemma}

\begin{lemma}[Vague convergence of Gaussian polynomial integrals]
\label{lem:vague-gaussian-polynomials}
\lean{Preiss.vague_integral_polynomial_mul_gaussian}
\uses{def:polynomial-growth-bound, lem:gaussian-uniform-tails, lem:vague-moving-tests, lem:vague-uniform-norm-tails, lem:polynomial-gaussian-norm-tails}







If $\mu_i\rightharpoonup\mu$ regularly with one common polynomial growth bound, then for every real polynomial $P$ and $s>0$, $\int P(z)e^{-s|z|^2}\,d\mu_i\to\int P(z)e^{-s|z|^2}\,d\mu$.
\par\emph{Role: Preiss proof.}
\par Lean declarations: \texttt{Preiss.vague\_integral\_polynomial\_mul\_gaussian}.
\end{lemma}
\begin{proof}[Planned proof obligations]

For every $d\in\mathbb N$ and $s>0$,
$\int |z|^d e^{-s|z|^2}\,d\nu<\infty$ and $0<I_\nu(s)<\infty$.
For fixed $s,d$, the tails outside $B(0,R)$ tend to zero uniformly over
$\mathcal U_m$. State the more reusable version for a family satisfying one
common polynomial bound $\mu(B(0,R))\leq C(1+R)^p$.
Use this to extend vague convergence to polynomial-times-Gaussian tests.
It also justifies differentiation with respect to positive $s$ locally in $s$.
\par Corrections and boundary conditions: Vague convergence alone is insufficient.
\par Source: De Lellis Lemma 7.2, Definition 7.4, Lemma 7.6, Proposition 7.7, pp.70-84.
\end{proof}

\begin{lemma}[Gaussian quadratic form second moment and trace bound]
\label{lem:gaussian-quadratic-trace-bound}
\lean{Preiss.gaussian_integral_bilinear_eq_trace, Preiss.gaussian_integral_bilinear_sq, Preiss.gaussian_integral_bilinear_sq_ge_trace_sq}
\uses{def:bilinear-trace, lem:bilinear-trace-basis, lem:gaussian-wick-fourth, mathlib:standard-gaussian, lem:gaussian-projection-integrability}







For a symmetric bilinear form $A$ on an $m$-dimensional Euclidean space, $m>0$, $\int A(v,v)^2\,d\gamma=(\operatorname{tr}A)^2+2\|A\|_{HS}^2\geq(1+2/m)(\int A(v,v)\,d\gamma)^2$.
\par\emph{Role: Preiss proof.}
\par Lean declarations: \texttt{Preiss.gaussian\_integral\_bilinear\_eq\_trace, Preiss.gaussian\_integral\_bilinear\_sq, Preiss.gaussian\_integral\_bilinear\_sq\_ge\_trace\_sq}.
\end{lemma}
\begin{proof}[Planned proof obligations]

For every symmetric bilinear form $A$ on an $m$-dimensional Euclidean space,
\[
 \int A(v,v)^2\,d\gamma=(\operatorname{tr}A)^2+
 2\|A\|_{\mathrm{HS}}^2
 \geq(1+2/m)\left(\int A(v,v)\,d\gamma\right)^2.
\]
This is a reusable Gaussian/spectral theorem independent of uniform measures.
Combining it with R13 gives
$\int\widehat b(v,v)\,d\gamma\geq-8m|b|^2/(m+2)$.
R08 and R04 therefore give
$\operatorname{tr}b_2^{(2)}\geq16|b|^2/(m+2)$.
Since that trace is zero, $b=0$; R03 makes the support planar, and plane-supported
uniform classification finishes rigidity. This is the actual content hidden
inside the current single trace-rigidity node.
\par Source: De Lellis Proposition 9.1, Lemmas 9.2,9.4,9.9,9.10 and Proposition 9.5, pp.95-109.
\end{proof}

\begin{lemma}[Operator norm remainder from finitely many support features]
\label{lem:moment-feature-remainder}
\lean{Preiss.featureMoment_operator_remainder}
\uses{def:moments, def:polynomial-feature, def:unit-uniform, lem:moment-feature-annihilation, lem:polynomial-feature-identities, lem:taylor-one, lem:finite-dimensional-operator-bigO}







The norm polynomials $|x|^{2j}$ are linear functionals of $\Phi(x)$ for $j\leq q$. Choose a finite basis of $V$ from support features. The pointwise remainders at those basis points imply $\|L_s|V-\sum_{j=1}^q s^j\ell_j|V/j!\|=O(s^{q+1/2})$ as $s\downarrow0$.
\par\emph{Role: Preiss proof.}
\par Lean declarations: \texttt{Preiss.exists\_normPolynomial\_feature\_functionals, Preiss.featureMoment\_operator\_remainder}.
\end{lemma}
\begin{proof}[Planned proof obligations]

For fixed $q$, take the finite sum of symmetric tensor degrees $1$ through
$2q$, with the weighted inner product from Section 7.3, and let $V$ be spanned
by the vectors $(x,x^{\otimes2},\ldots,x^{\otimes2q})$ for support points $x$.
The moment functional $b_s$ vanishes on $F_s$ and M04 determines
$b_s|V$ to order $q$. Choose a finite basis from these spanning support vectors:
its finitely many remainder estimates give an operator-norm
$O(s^{q+1/2})$ bound on $V$. Composing with the bounded $Q_s$ and using M05
gives the Taylor expansion of $b_s$.
Do not implement the source's claim that the unrescaled symmetrized coordinate
monomials are orthonormal: for example
$|e_1\odot e_2|^2=1/2$ for orthogonal unit vectors under its displayed inner
product. Use the inherited tensor inner product or insert multinomial weights.
A finite polynomial-coordinate model is also possible, provided its pairing
realizes $\langle u^{\otimes k},v^{\otimes k}\rangle=\langle u,v\rangle^k$.
\par Corrections and boundary conditions: Use a finite support-feature basis to obtain a genuine linear operator-norm estimate.
\par Source: De Lellis Lemma 7.2, Definition 7.4, Lemma 7.6, Proposition 7.7, pp.70-84.
\end{proof}

\section{Uniform moments and expansions}

\begin{definition}[Continuous symmetric Gaussian moment forms]
\label{def:moments}
\lean{Preiss.MomentForm, Preiss.gaussianMass, Preiss.gaussianMoment, Preiss.gaussianMomentForm, Preiss.momentKernel}





For $s>0$, define the continuous symmetric $k$-linear form $b_{k,s}(u_1,\ldots,u_k)=\frac{(2s)^k}{k! I_\nu(s)}\int e^{-s|z|^2}\prod_{i=1}^k\langle z,u_i\rangle\,d\nu(z)$. Prove its diagonal evaluation agrees with the scalar Gaussian moment; the degree-zero form equals $1$ for nonzero measures with finite Gaussian mass.
\par\emph{Role: Preiss proof.}
\par Lean declarations: \texttt{Preiss.MomentForm, Preiss.gaussianMass, Preiss.gaussianMoment, Preiss.gaussianMomentForm, Preiss.gaussianMomentForm\_permute, Preiss.momentBilinear, Preiss.momentKernel}.
\end{definition}

\begin{lemma}[Feature moment functional and weighted annihilator]
\label{lem:moment-feature-annihilation}
\lean{Preiss.featureMoment_apply_feature, Preiss.featureMoment_apply_feature_of_growth}
\uses{def:moments, def:polynomial-feature, def:unit-uniform, lem:polynomial-feature-identities, def:feature-moment, lem:gaussian-polynomial-integrable, lem:gaussian-feature-integrability}







Let $V=\operatorname{span}\Phi(\operatorname{supp}\nu)$ and $A(s)=\sum_{k=1}^{2q}s^kP_k$. The functional $L_s(u)=I(s)^{-1}\int\langle A(s)u,\Phi(z)\rangle e^{-s|z|^2}\,d\nu$ is continuous linear, has $L_s(\Phi(x))=\sum_{k=1}^{2q}b_{k,s}(x^k)$, and vanishes on $A(s)^{-1}(V^\perp)$.
\par\emph{Role: Preiss proof.}
\par Lean declarations: \texttt{Preiss.featureMoment, Preiss.featureMoment\_apply\_feature, Preiss.featureMoment\_vanishes\_weightedComplement, Preiss.supportFeatureSpan}.
\end{lemma}
\begin{proof}[Planned proof obligations]
For fixed $q$, take the finite sum of symmetric tensor degrees $1$ through
$2q$, with the weighted inner product from Section 7.3, and let $V$ be spanned
by the vectors $(x,x^{\otimes2},\ldots,x^{\otimes2q})$ for support points $x$.
The moment functional $b_s$ vanishes on $F_s$ and M04 determines
$b_s|V$ to order $q$. Choose a finite basis from these spanning support vectors:
its finitely many remainder estimates give an operator-norm
$O(s^{q+1/2})$ bound on $V$. Composing with the bounded $Q_s$ and using M05
gives the Taylor expansion of $b_s$.
Do not implement the source's claim that the unrescaled symmetrized coordinate
monomials are orthonormal: for example
$|e_1\odot e_2|^2=1/2$ for orthogonal unit vectors under its displayed inner
product. Use the inherited tensor inner product or insert multinomial weights.
A finite polynomial-coordinate model is also possible, provided its pairing
realizes $\langle u^{\otimes k},v^{\otimes k}\rangle=\langle u,v\rangle^k$.
\par Source: De Lellis Lemma 7.2, Definition 7.4, Lemma 7.6, Proposition 7.7, pp.70-84.
\end{proof}

\begin{definition}[Gaussian moment coefficient families]
\label{def:moment-coefficients}
\lean{Preiss.MomentCoefficientFamily, Preiss.momentTaylorPolynomial, Preiss.IsMomentExpansion}
\uses{def:moments}








A coefficient family assigns a continuous $k$-linear form $B_k^{(j)}$ to each
pair of natural numbers $k,j$. Its Taylor polynomial of order $q$ is
$\sum_{j=1}^q s^j B_k^{(j)}/j!$. The expansion predicate requires symmetry
and an operator-norm remainder $o(s^q)$ as $s\downarrow0$, for every $k,q>0$.
Existence, uniqueness in positive degrees, degree vanishings, and support
identities are separate theorems. No coefficient is selected from an
unproved existence theorem inside this definition.
\end{definition}

\begin{lemma}[Evaluation of Gaussian moment forms]
\label{lem:moment-form-evaluation}
\lean{Preiss.momentKernel_apply, Preiss.gaussianMomentForm_apply, Preiss.momentBilinear_apply}
\uses{def:moments}


If the Gaussian-weighted multilinear kernel is Bochner integrable, evaluation of the bundled moment form agrees with its scalar integral formula.
\end{lemma}

\begin{lemma}[Radial pushforward distribution and integrals]
\label{lem:radial}
\lean{Preiss.uniform_radial_pushforward}
\uses{def:unit-uniform, gmt:uniform-radial, lem:unit-uniform-adapter}




For $\nu\in\mathcal U_m$ and $y\in\operatorname{supp}\nu$, identify the
pushforward under $z\mapsto|z-y|$ with that of
$\omega_m^{-1}\mathcal L^m$ under the norm. State nonnegative Borel integrals
as lintegrals, then derive signed Bochner-integral identities only with their
integrability hypotheses. GMT already gives equality of continuous radial
integrals on closed balls. The full pushforward distribution and unbounded
Gaussian consequences are new adapters.
\par\emph{Role: Preiss proof.}
\par Lean declarations: \texttt{Preiss.uniform\_radial\_pushforward}.
\end{lemma}
\begin{proof}[Planned proof obligations]

For $\nu\in\mathcal U_m$ and $y\in\operatorname{supp}\nu$, identify the
pushforward under $z\mapsto|z-y|$ with that of
$\omega_m^{-1}\mathcal L^m$ under the norm. State nonnegative Borel integrals
as lintegrals, then derive signed Bochner-integral identities only with their
integrability hypotheses. GMT already gives equality of continuous radial
integrals on closed balls. The full pushforward distribution and unbounded
Gaussian consequences are new adapters.
\par Corrections and boundary conditions: State lintegral identity first; signed or Bochner identities require integrability. Merge the two existing consecutive proof blocks under this label.
\par Source: De Lellis Lemma 7.2, Definition 7.4, Lemma 7.6, Proposition 7.7, pp.70-84.
\end{proof}

\begin{lemma}[Moment covariance, scalar invariance and dilation]
\label{lem:moment-covariance}
\lean{Preiss.gaussianMomentForm_dilate, Preiss.gaussianMomentForm_smul, Preiss.gaussianMomentForm_map_linearIsometry}
\uses{def:dimensional-blowup, def:moments, def:unit-uniform, lem:blowup-composition, lem:moment-form-evaluation, lem:moment-kernel-calculus}








Package the current scalar \texttt{gaussianMoment} as a continuous symmetric
$k$-multilinear form, prove the evaluation formula, scalar-invariance,
orthogonal covariance, and the scaling law
\[
 b^{D^m_{0,r}\nu}_{k,s}=r^k b^\nu_{k,s/r^2}.
\]
The definition with $k=0$ is the constant one for uniform measures; Taylor
coefficient assertions below start at $k=1$.
\par\emph{Role: Preiss proof.}
\par Lean declarations: \texttt{Preiss.gaussianMomentForm\_dilate, Preiss.gaussianMomentForm\_map\_linearIsometry, Preiss.gaussianMomentForm\_smul}.
\end{lemma}
\begin{proof}[Planned proof obligations]

Package the current scalar \texttt{gaussianMoment} as a continuous symmetric
$k$-multilinear form, prove the evaluation formula, scalar-invariance,
orthogonal covariance, and the scaling law
\[
 b^{D^m_{0,r}\nu}_{k,s}=r^k b^\nu_{k,s/r^2}.
\]
The definition with $k=0$ is the constant one for uniform measures; Taylor
coefficient assertions below start at $k=1$.
\par Source: De Lellis Lemma 7.2, Definition 7.4, Lemma 7.6, Proposition 7.7, pp.70-84.
\end{proof}

\begin{lemma}[Uniform operator norm bounds for Gaussian moments]
\label{lem:moment-norm-bound}
\lean{Preiss.integrable_gaussianMomentKernel, Preiss.norm_gaussianMomentForm_le, Preiss.IsUnitUniform.hasPolynomialGrowthBound, Preiss.integrable_gaussianMomentKernel_of_integrable_norm, Preiss.gaussianMass_unitUniform}
\uses{def:moments, def:unit-uniform, lem:moment-form-evaluation, lem:radial, lem:moment-kernel-calculus, lem:gaussian-ball-power-parameter, lem:gaussian-power-envelope}








For $\nu\in\mathcal U_m$, $k\geq1$ and $s>0$, $\|b_{k,s}\|\leq C_m 2^k k^{k/2}s^{k/2}/k!$, where $C_m$ depends only on $m$.
\par\emph{Role: Preiss proof.}
\par Lean declarations: \texttt{Preiss.integrable\_gaussianMomentKernel, Preiss.norm\_gaussianMomentForm\_le}.
\end{lemma}
\begin{proof}[Planned proof obligations]

Prove the norm bound
$|b_{k,s}(u_1,\ldots,u_k)|\leq
C_m 2^k k^{k/2}s^{k/2}\prod_i|u_i|/k!$.
For $x$ in the support, the absolutely convergent series
$\sum_{k\geq0} b_{k,s}(x^k)=e^{s|x|^2}$ follows by translating the
Gaussian. A valid dominating function is
$e^{-s|z|^2+2s|z||x|}$, bounded by a constant times
$e^{-(s/2)|z|^2}$, not the expression with a negative cross term printed in
the source. Taylor truncation gives, for each $q\geq1$,
\[
 \left|\sum_{k=1}^{2q}b_{k,s}(x^k)
 -\sum_{j=1}^{q}\frac{s^j|x|^{2j}}{j!}\right|
 \leq C_{m,q}(s|x|^2)^{q+1/2}.
\]
For $s|x|^2\geq1$ use the finite sum and the norm bound; for small values use
the convergent exponential series with its remainder bound.
\par Source: De Lellis Lemma 7.2, Definition 7.4, Lemma 7.6, Proposition 7.7, pp.70-84.
\end{proof}

\begin{lemma}[Comparison of homogeneous ball and Gaussian energies]
\label{lem:homogeneous-energy-radial}
\lean{Preiss.integral_gaussian_mul_homogeneous}
\uses{def:conical, def:plane, lem:radial, lem:homogeneous-ball-scaling, lem:gaussian-ball-power}




If $\nu$ is regular and $m$-homogeneous and $h\geq0$ is measurable and $a$-homogeneous with $m+a>0$ and integrable on the unit closed ball, then the Gaussian-weighted integrand is integrable and $\int e^{-|z|^2}h(z)\,d\nu=\Gamma((m+a)/2+1)\int_{\overline B(0,1)}h(z)\,d\nu$. In particular $h(z)=\operatorname{dist}(z,V)^2$ gives the constant $\Gamma(m/2+2)$.
\par\emph{Role: Preiss proof.}
\par Lean declarations: \texttt{Preiss.integral\_gaussian\_mul\_homogeneous}.
\end{lemma}
\begin{proof}[Planned proof obligations]

For every plane $V$, the distance-squared measure has homogeneous degree
$m+2$, so
$\int e^{-|z|^2}\operatorname{dist}(z,V)^2\,d\nu
=C_m\int_{\overline B(0,1)}\operatorname{dist}(z,V)^2\,d\nu$
for a finite positive constant $C_m$ obtained by radial integration.
State and prove this for nonnegative homogeneous integrands more generally.
It is the missing explicit bridge from the theorem's unit-ball energy to the
Gaussian energy used in the spectral proof. The implementation applies the
general centered-ball power-law formula to $h\nu$: exponential tails and
nonnegative Fubini reduce the integral to Euler's Gamma formula.
\par Corrections and boundary conditions: Explicit bridge between the gap theorem unit-ball energy and the spectral Gaussian energy.
\par Source: De Lellis Corollary 8.4, Lemmas 8.6-8.7, Proposition 8.5, pp.85-94; author errata for the dimension-two shortcut.
\end{proof}

\begin{lemma}[Absolutely convergent Gaussian translation series]
\label{lem:gaussian-translation-series}
\lean{Preiss.hasSum_gaussianMoment_diagonal, Preiss.integral_gaussian_center_unitUniform, Preiss.integral_gaussian_exp_inner_unitUniform}
\uses{def:moments, def:unit-uniform, lem:gaussian-polynomial-integrable, lem:moment-form-evaluation, lem:moment-norm-bound, lem:radial, lem:gaussian-ball-power-parameter, lem:integrated-gaussian-exponential-series}







For $\nu\in\mathcal U_m$, $x\in\operatorname{supp}\nu$, $s>0$, the absolutely convergent series satisfies $\sum_{k=0}^{\infty}b_{k,s}(x^k)=e^{s|x|^2}$.
\par\emph{Role: Preiss proof.}
\par Lean declarations: \texttt{Preiss.hasSum\_gaussianMoment\_diagonal}.
\end{lemma}
\begin{proof}[Planned proof obligations]

Prove the norm bound
$|b_{k,s}(u_1,\ldots,u_k)|\leq
C_m 2^k k^{k/2}s^{k/2}\prod_i|u_i|/k!$.
For $x$ in the support, the absolutely convergent series
$\sum_{k\geq0} b_{k,s}(x^k)=e^{s|x|^2}$ follows by translating the
Gaussian. A valid dominating function is
$e^{-s|z|^2+2s|z||x|}$, bounded by a constant times
$e^{-(s/2)|z|^2}$, not the expression with a negative cross term printed in
the source. Taylor truncation gives, for each $q\geq1$,
\[
 \left|\sum_{k=1}^{2q}b_{k,s}(x^k)
 -\sum_{j=1}^{q}\frac{s^j|x|^{2j}}{j!}\right|
 \leq C_{m,q}(s|x|^2)^{q+1/2}.
\]
For $s|x|^2\geq1$ use the finite sum and the norm bound; for small values use
the convergent exponential series with its remainder bound.
\par Corrections and boundary conditions: Use domination exp(-s|z|\textasciicircum{}2+2s|z||x|) <= C exp(-(s/2)|z|\textasciicircum{}2); the source has an invalid sign in the displayed dominator.
\par Source: De Lellis Lemma 7.2, Definition 7.4, Lemma 7.6, Proposition 7.7, pp.70-84.
\end{proof}

\begin{lemma}[Quantitative finite moment remainder]
\label{lem:taylor-one}
\lean{Preiss.gaussianMoment_taylor_remainder}
\uses{def:moments, def:unit-uniform, lem:gaussian-translation-series, lem:moment-norm-bound, lem:scalar-gaussian-series-remainder}







For every $m,q\geq1$ there is $C_{m,q}\geq0$ such that for every $\nu\in\mathcal U_m$, support point $x$ and $s>0$, $|\sum_{k=1}^{2q} b_{k,s}(x^k)-\sum_{j=1}^q s^j|x|^{2j}/j!|\leq C_{m,q}(s|x|^2)^{q+1/2}$.
\par\emph{Role: Preiss proof.}
\par Lean declarations: \texttt{Preiss.gaussianMoment\_taylor\_remainder}.
\end{lemma}
\begin{proof}[Planned proof obligations]

Prove the norm bound
$|b_{k,s}(u_1,\ldots,u_k)|\leq
C_m 2^k k^{k/2}s^{k/2}\prod_i|u_i|/k!$.
For $x$ in the support, the absolutely convergent series
$\sum_{k\geq0} b_{k,s}(x^k)=e^{s|x|^2}$ follows by translating the
Gaussian. A valid dominating function is
$e^{-s|z|^2+2s|z||x|}$, bounded by a constant times
$e^{-(s/2)|z|^2}$, not the expression with a negative cross term printed in
the source. Taylor truncation gives, for each $q\geq1$,
\[
 \left|\sum_{k=1}^{2q}b_{k,s}(x^k)
 -\sum_{j=1}^{q}\frac{s^j|x|^{2j}}{j!}\right|
 \leq C_{m,q}(s|x|^2)^{q+1/2}.
\]
For $s|x|^2\geq1$ use the finite sum and the norm bound; for small values use
the convergent exponential series with its remainder bound.
\par Source: De Lellis Lemma 7.2, Definition 7.4, Lemma 7.6, Proposition 7.7, pp.70-84.
\end{proof}

\begin{lemma}[Coherent operator-norm moment Taylor coefficients]
\label{lem:taylor-two}
\lean{Preiss.exists_gaussianMoment_taylor_coefficients, Preiss.gaussianMomentCoefficient_unique}
\uses{def:moment-coefficients, lem:finite-moment-taylor-jets, lem:coherent-finite-jets, lem:asymptotic-taylor-vanishing}







For $\nu\in\mathcal U_m$ there are symmetric continuous forms $b_k^{(j)}$ for $k,j\geq1$ such that for all $q\geq1$ and $k\geq1$, $b_{k,s}=\sum_{j=1}^q s^j b_k^{(j)}/j!+o(s^q)$ in operator norm as $s\downarrow0$. These coefficients are unique and independent of truncation order.
\par\emph{Role: Preiss proof.}
\par Lean declarations: \texttt{Preiss.exists\_gaussianMoment\_taylor\_coefficients, Preiss.gaussianMomentCoefficient\_unique}.
\end{lemma}
\begin{proof}[Planned proof obligations]

There exist symmetric continuous forms $b_k^{(j)}$, $k,j\geq1$, such that for
every $q\geq1$,
\[
 b_{k,s}=\sum_{j=1}^{q}\frac{s^j}{j!}b_k^{(j)}+o(s^q),\qquad
 b_k^{(j)}=0\ (k>2j),\qquad
 \sum_{k=1}^{2q}b_k^{(q)}(x^k)=|x|^{2q}\ (x\in\operatorname{supp}\nu).
\]
Uniqueness of finite Taylor coefficients makes choices from different truncation
orders coherent. If $k>2q$, M04 gives $o(s^q)$ directly. The operator-norm
statement is stronger and safer than merely pointwise convergence on the support.
\par Source: De Lellis Lemma 7.2, Definition 7.4, Lemma 7.6, Proposition 7.7, pp.70-84.
\end{proof}

\begin{lemma}[Degree restriction on Taylor coefficients]
\label{lem:moment-coefficient-vanishing}
\lean{Preiss.gaussianMomentCoefficient_eq_zero_of_two_mul_lt}
\uses{def:moment-coefficients, lem:asymptotic-taylor-vanishing, lem:moment-operator-growth}







For $k,j\geq1$ with $k>2j$, $b_k^{(j)}=0$.
\par\emph{Role: Preiss proof.}
\par Lean declarations: \texttt{Preiss.gaussianMomentCoefficient\_eq\_zero\_of\_two\_mul\_lt}.
\end{lemma}
\begin{proof}[Planned proof obligations]

There exist symmetric continuous forms $b_k^{(j)}$, $k,j\geq1$, such that for
every $q\geq1$,
\[
 b_{k,s}=\sum_{j=1}^{q}\frac{s^j}{j!}b_k^{(j)}+o(s^q),\qquad
 b_k^{(j)}=0\ (k>2j),\qquad
 \sum_{k=1}^{2q}b_k^{(q)}(x^k)=|x|^{2q}\ (x\in\operatorname{supp}\nu).
\]
Uniqueness of finite Taylor coefficients makes choices from different truncation
orders coherent. If $k>2q$, M04 gives $o(s^q)$ directly. The operator-norm
statement is stronger and safer than merely pointwise convergence on the support.
\par Source: De Lellis Lemma 7.2, Definition 7.4, Lemma 7.6, Proposition 7.7, pp.70-84.
\end{proof}

\begin{lemma}[Vanishing total second Taylor trace]
\label{lem:moment-second-total-trace}
\lean{Preiss.trace_gaussianMomentForm_two, Preiss.trace_secondMomentCoefficient_eq_zero}
\uses{lem:moment-trace-operator, lem:gaussian-radial-second-moment, lem:asymptotic-taylor-vanishing, def:moment-coefficients}






For every $s>0$, $\operatorname{tr}b_{2,s}=ms$. Consequently $\operatorname{tr}b_2^{(2)}=0$.
\par\emph{Role: Preiss proof.}
\par Lean declarations: \texttt{Preiss.trace\_gaussianMomentForm\_two, Preiss.trace\_secondMomentCoefficient\_eq\_zero}.
\end{lemma}
\begin{proof}[Planned proof obligations]

The exact radial identity $\operatorname{tr}b_{2,s}=ms$ and M07 imply
$\operatorname{tr}b_2^{(2)}=0$.
R03 and the first-moment expansion give
$\operatorname{tr}(b_2^{(2)}|V^\perp)=8|b|^2
=2\|b_1^{(1)}\|^2$.
These are separate, comparatively short lemmas.
\par Source: De Lellis Proposition 9.1, Lemmas 9.2,9.4,9.9,9.10 and Proposition 9.5, pp.95-109.
\end{proof}

\begin{lemma}[Polynomial support identities for Taylor coefficients]
\label{lem:moment-support-polynomial}
\lean{Preiss.sum_gaussianMomentCoefficient_eq_norm_pow}
\uses{def:moment-coefficients, lem:moment-coefficient-vanishing, lem:taylor-one, lem:asymptotic-taylor-vanishing}







For every $q\geq1$ and $x\in\operatorname{supp}\nu$, $\sum_{k=1}^{2q}b_k^{(q)}(x^k)=|x|^{2q}$.
\par\emph{Role: Preiss proof.}
\par Lean declarations: \texttt{Preiss.sum\_gaussianMomentCoefficient\_eq\_norm\_pow}.
\end{lemma}
\begin{proof}[Planned proof obligations]

There exist symmetric continuous forms $b_k^{(j)}$, $k,j\geq1$, such that for
every $q\geq1$,
\[
 b_{k,s}=\sum_{j=1}^{q}\frac{s^j}{j!}b_k^{(j)}+o(s^q),\qquad
 b_k^{(j)}=0\ (k>2j),\qquad
 \sum_{k=1}^{2q}b_k^{(q)}(x^k)=|x|^{2q}\ (x\in\operatorname{supp}\nu).
\]
Uniqueness of finite Taylor coefficients makes choices from different truncation
orders coherent. If $k>2q$, M04 gives $o(s^q)$ directly. The operator-norm
statement is stronger and safer than merely pointwise convergence on the support.
\par Source: De Lellis Lemma 7.2, Definition 7.4, Lemma 7.6, Proposition 7.7, pp.70-84.
\end{proof}

\begin{lemma}[Odd and even leading moment limits]
\label{lem:scaled-moment-limits}
\lean{Preiss.gaussianMoment_scaled_odd_limit, Preiss.gaussianMoment_scaled_even_limit}
\uses{def:moment-coefficients, lem:moment-coefficient-vanishing, lem:asymptotic-taylor-vanishing}





For odd $k\geq1$, $s^{-k/2}b_{k,s}\to0$. For $j\geq1$, $s^{-j}b_{2j,s}\to b_{2j}^{(j)}/j!$, in operator norm as $s\downarrow0$.
\par\emph{Role: Preiss proof.}
\par Lean declarations: \texttt{Preiss.gaussianMoment\_scaled\_even\_limit, Preiss.gaussianMoment\_scaled\_odd\_limit}.
\end{lemma}
\begin{proof}[Planned proof obligations]

For odd $k$, $s^{-k/2}b_{k,s}\to0$; for $k=2j$,
$s^{-j}b_{2j,s}\to b_{2j}^{(j)}/j!$.
The factorial is essential; the displayed summary on printed page 74 omits it,
whereas the calculation in Chapter 9 includes it correctly.
Use M03 with $s=r^{-2}$ to deduce the limit of every Gaussian-weighted
polynomial moment of $D^m_{0,r}\nu$ as $r\to\infty$.
\par Corrections and boundary conditions: Retain the factorial omitted in the summary on printed p.74.
\par Source: De Lellis Lemma 7.2, Definition 7.4, Lemma 7.6, Proposition 7.7, pp.70-84.
\end{proof}

\section{Uniform compactness and dimension reduction}

\begin{lemma}[Uniform ball laws are closed under vague limits]
\label{lem:uniform-vague-closed}
\lean{Preiss.isUnitUniform_of_vague}
\uses{def:unit-uniform, gmt:portmanteau, lem:vague-support-approximation, lem:vague-support-limit}






If $\nu_i\in\mathcal U_m$ and $\nu_i\rightharpoonup\nu$ with $\nu$ regular, then $\nu\in\mathcal U_m$.
\par\emph{Role: Preiss proof.}
\par Lean declarations: \texttt{Preiss.isUnitUniform\_of\_vague}.
\end{lemma}
\begin{proof}[Planned proof obligations]

A vaguely convergent sequence in $\mathcal U_m$ has a limit in $\mathcal U_m$.
The family is sequentially compact by upstream compact mass bounds.
First approximate support centres using S04; squeeze open and closed radii to
recover the exact ball law, rather than assuming spheres of an arbitrary limit
already have zero mass. The origin's positive ball mass excludes the zero limit.
Also prove the variant where the measures are uniform with the same constant
and their supports approach the origin, without requiring the origin in every
support. This is needed in the bounded-hole argument R05 below.
\par Source: De Lellis Chapters 3 and 6; general supporting analysis.
\end{proof}

\begin{lemma}[Uniform compactness with centers approaching the support]
\label{lem:uniform-near-support-compactness}
\lean{Preiss.exists_vague_subsequence_uniform_of_dist_support_tendsto_zero}
\uses{def:uniform-constant, lem:uniform-constant-one, lem:uniform-scaling, lem:uniform-compactness, lem:vague-moving-blowups, lem:blowup-composition}
Let $\nu_i$ be regular measures with nonempty support and support-centered closed-ball law $\nu_i(\overline B(y,r))=r^m$ for every support center $y$ and $r>0$. If $\operatorname{dist}(0,\operatorname{supp}\nu_i)\to0$, some subsequence converges vaguely to a member of $\mathcal U_m$. This includes $m=0$.
\end{lemma}
\begin{proof}[Proof sketch]
Choose a nearest support point $y_i$ to the origin; closedness and properness give existence because each support is nonempty. The distance hypothesis implies $y_i\to0$. Recenter each measure by $D^m_{y_i,1}$, whose unit-uniform ball law follows from blow-up transport. Extract a unit-uniform vague limit by centered compactness, then undo translations using moving blow-up continuity and the composition identity. No initial-index truncation or positive-dimensional hypothesis is needed.
\end{proof}

\begin{lemma}[Hyperplane transport with an explicit isometry]
\label{lem:hyperplane-transport-map}
\uses{def:uniform-constant, lem:closed-map-support, lem:proper-map-regularity}
\lean{Preiss.uniform_hyperplane_transport, Preiss.uniformWithConstant_transport, Preiss.uniformWithConstant_map_linearIsometry, Preiss.isUnitUniform_map_linearIsometry}
\uses{def:dimensional-blowup, def:plane, gmt:hyperplane-transport, lem:blowup-composition}






Given an isometric identification $L:\mathbb R^{k-1}\simeq w^\perp$, return
the transported measure together with the equality of its pushforward with
$\lambda$, its ball law, support correspondence and blow-up covariance.
The existential theorem
\texttt{SUniformAux.exists\_isSUniform\_of\_hyperplane}
does not export enough data by itself.
\par\emph{Role: Preiss proof.}
\par Lean declarations: \texttt{Preiss.isUnitUniform\_map\_linearIsometry, Preiss.uniformWithConstant\_map\_linearIsometry, Preiss.uniformWithConstant\_transport, Preiss.uniform\_hyperplane\_transport}.
\end{lemma}
\begin{proof}[Planned proof obligations]

Given an isometric identification $L:\mathbb R^{k-1}\simeq w^\perp$, return
the transported measure together with the equality of its pushforward with
$\lambda$, its ball law, support correspondence and blow-up covariance.
The existential theorem
\texttt{SUniformAux.exists\_isSUniform\_of\_hyperplane}
does not export enough data by itself.
\par Source: De Lellis Lemma 6.6, Theorem 6.8, pp.57-58,63-65; Corollary 3.9 and corrected Lemma 3.8, pp.15,25-26; author errata.
\end{proof}

\begin{lemma}[Sequential compactness of centered unit uniform measures]
\label{lem:uniform-compactness}
\lean{Preiss.exists_vague_subsequence_unitUniform}
\uses{def:unit-uniform, lem:uniform-vague-closed, thm:compactness}






Every sequence in $\mathcal U_m$ has a subsequence converging vaguely to an element of $\mathcal U_m$.
\par\emph{Role: Preiss proof.}
\par Lean declarations: \texttt{Preiss.exists\_vague\_subsequence\_unitUniform}.
\end{lemma}
\begin{proof}[Planned proof obligations]

A vaguely convergent sequence in $\mathcal U_m$ has a limit in $\mathcal U_m$.
The family is sequentially compact by upstream compact mass bounds.
First approximate support centres using S04; squeeze open and closed radii to
recover the exact ball law, rather than assuming spheres of an arbitrary limit
already have zero mass. The origin's positive ball mass excludes the zero limit.
Also prove the variant where the measures are uniform with the same constant
and their supports approach the origin, without requiring the origin in every
support. This is needed in the bounded-hole argument R05 below.
\par Source: De Lellis Chapters 3 and 6; general supporting analysis.
\end{proof}

\begin{lemma}[Finite metric covers of uniform families]
\label{lem:uniform-finite-metric-covers}
\lean{Preiss.unitUniform_totallyBounded}
\uses{def:unit-uniform, def:vague-metrization, lem:uniform-compactness, lem:vague-moving-blowups, lem:blowup-composition, lem:uniform-scaling}






For a supplied vague metrization and every $\varepsilon>0$ and $0<c\leq C<\infty$, the centered uniform measures with ball constants in $[c,C]$ admit a finite cover by vague-metric balls of radius $\varepsilon$. This includes the family $\mathcal U_m$ and all its support-centered translates.
\par\emph{Role: Preiss proof.}
\par Lean declarations: \texttt{Preiss.unitUniform\_totallyBounded}.
\end{lemma}
\begin{proof}[Proof sketch]
Extract a convergent subsequence of ball constants in $[c,C]$, normalize by their positive reciprocals, and use compactness of unit-uniform measures. Moving scalar convergence recovers convergence of the original measures. The supplied metric's convergence field gives sequential compactness, hence total boundedness. This proof does not use the existence of a metrization.
\end{proof}

\begin{lemma}[Existence of uniform blow-down limits]
\label{lem:infinity-compactness}
\uses{lem:uniform-scaling, lem:uniform-compactness}
\lean{Preiss.exists_tangentAtInfinity}
\uses{def:tangent-at-infinity, def:unit-uniform, lem:blowup-composition, lem:uniform-compactness}






For every $\nu\in\mathcal U_m$ and every sequence $r_i\to\infty$ of positive radii, a subsequence of $D^m_{0,r_i}\nu$ converges vaguely to a member of $\mathcal U_m$.
\par\emph{Role: Preiss proof.}
\par Lean declarations: \texttt{Preiss.exists\_tangentAtInfinity}.
\end{lemma}
\begin{proof}[Planned proof obligations]

Use S05 to obtain subsequential uniform limits, M02 and M08 to identify their
Gaussian moments, and M09 to identify the limits. The compactness principle
``every sequence escaping to infinity has a convergent subsequence with the
same limit'' gives full convergence. A separate output is that every positive
sequence $r_i\to\infty$ converges to this same measure. Scaling by a fixed
factor preserves the set of limits at infinity, so S02--S03 imply conicality.
\par Source: De Lellis Lemma 7.2, Definition 7.4, Lemma 7.6, Proposition 7.7, pp.70-84.
\end{proof}

\begin{lemma}[Common tangent normalization at a density point]
\label{lem:dimensional-tangent-ball-constant}
\lean{Preiss.dimensionalTangent_ball_constant}
\uses{def:dimensional-tangent, def:positive-density, def:unit-uniform, gmt:continuity-sets, gmt:normalized-tangent, lem:tangent-ray, thm:doubling, lem:dimensional-tangent-origin-ball-law}






At $x$ with density $q=\lim_{r\downarrow0}\mu(B(x,r))/r^m\in(0,\infty)$, if all arbitrary-scalar tangents are uniform, every dimensional tangent contains $0$ in its support and has ball law $\nu(B(y,r))=qr^m$.
\par\emph{Role: Preiss proof.}
\par Lean declarations: \texttt{Preiss.dimensionalTangent\_ball\_constant}.
\end{lemma}
\begin{proof}[Planned proof obligations]

Let $q_x=\lim_{r\downarrow0}\mu(B(x,r))/r^m\in(0,\infty)$.
If every arbitrary-scalar tangent is uniform, every dimensional tangent
has ball constant exactly $q_x$ and contains zero in its support.
Use doubling to get the support assertion, the uniform sphere-null property,
and convergence on continuity balls to determine the constant. Dividing the
original measure by $q_x$ therefore puts its dimensional tangents in
$\mathcal U_m$. Existing arbitrary-scalar and dimensional tangent sets are rays
of one another, not literally equal. A pointwise Part B theorem must use this
adapter before applying a universal flatness threshold.
\par Source: De Lellis Chapters 3 and 6; general supporting analysis.
\end{proof}

\begin{lemma}[Reduction to unit uniform dimensional tangents]
\label{lem:density-normalization}
\lean{Preiss.density_normalization_unitUniform}
\uses{def:dimensional-tangent, def:positive-density, def:unit-uniform, lem:blowup-composition, lem:dimensional-tangent-ball-constant, lem:flat-positive-scalar, lem:dimensional-tangent-operations, lem:uniform-scaling}






Under S08a, replacing $\mu$ by $q^{-1}\mu$ identifies its dimensional tangent family with $q^{-1}\operatorname{Tan}_m(\mu,x)\subset\mathcal U_m$ and preserves flatness of tangents.
\par\emph{Role: Preiss proof.}
\par Lean declarations: \texttt{Preiss.density\_normalization\_unitUniform}.
\end{lemma}
\begin{proof}[Planned proof obligations]

Let $q_x=\lim_{r\downarrow0}\mu(B(x,r))/r^m\in(0,\infty)$.
If every arbitrary-scalar tangent is uniform, every dimensional tangent
has ball constant exactly $q_x$ and contains zero in its support.
Use doubling to get the support assertion, the uniform sphere-null property,
and convergence on continuity balls to determine the constant. Dividing the
original measure by $q_x$ therefore puts its dimensional tangents in
$\mathcal U_m$. Existing arbitrary-scalar and dimensional tangent sets are rays
of one another, not literally equal. A pointwise Part B theorem must use this
adapter before applying a universal flatness threshold.
\par Corrections and boundary conditions: A universal gap threshold applies after one common normalization, not independently normalized arbitrary-scalar tangents.
\par Source: De Lellis Chapters 3 and 6; general supporting analysis.
\end{proof}

\begin{lemma}[Half-space reduction retaining the tangent witness]
\label{lem:halfspace-tangent-witness}
\lean{Preiss.exists_halfSpace_dimensionalTangent}
\uses{gmt:support-not-full, gmt:uniform-halfspace, gmt:uniform-tangents, lem:dimensional-tangent-ball-constant, lem:unit-uniform-tangent-existence, gmt:tangent-open-null}






For uniform $\nu$ with non-full support there exist a support point $y$,
a nonzero vector $e$, and a tangent $\lambda$ at $y$ supported in
$\{z:\langle z,e\rangle\geq0\}$. Preserve the exact normalization when
choosing the tangent. GMT's \texttt{SUniformAux.exists\_halfSpace} proves
this in its body but exports only the existence of the new uniform measure.
The adapter should reuse that proof pattern and retain its witness.
\par\emph{Role: Preiss proof.}
\par Lean declarations: \texttt{Preiss.exists\_halfSpace\_dimensionalTangent}.
\end{lemma}
\begin{proof}[Planned proof obligations]

For uniform $\nu$ with non-full support there exist a support point $y$,
a nonzero vector $e$, and a tangent $\lambda$ at $y$ supported in
$\{z:\langle z,e\rangle\geq0\}$. Preserve the exact normalization when
choosing the tangent. GMT's \texttt{SUniformAux.exists\_halfSpace} proves
this in its body but exports only the existence of the new uniform measure.
The adapter should reuse that proof pattern and retain its witness.
\par Corrections and boundary conditions: Existing upstream existential statement discards the tangent witness; preserve it and its normalization.
\par Source: De Lellis Lemma 6.6, Theorem 6.8, pp.57-58,63-65; Corollary 3.9 and corrected Lemma 3.8, pp.15,25-26; author errata.
\end{proof}

\begin{lemma}[Corrected hyperplane reduction with lineage]
\label{lem:hyperplane-tangent-lineage}
\lean{Preiss.exists_hyperplane_eq_or_dimensionalTangent}
\uses{gmt:barycenter-quadratic-bound, gmt:quadratic-support-tangent, gmt:uniform-hyperplane, gmt:zero-barycenter-hyperplane, lem:unit-uniform-tangent-existence}






For a centred uniform half-space measure $\nu$, there exist $w\ne0$ and a
uniform $\lambda$ supported in $w^\perp$ such that either $\lambda=\nu$
or $\lambda$ is a tangent to $\nu$ at zero, up to the explicitly fixed
positive normalization. The new $w$ need not be $e$.
The proof of GMT's \texttt{SUniformAux.exists\_hyperplane} already has
exactly these two branches. Expose
\texttt{SUniformAux.abs\_inner\_integral\_le},
\texttt{SUniformAux.support\_tangent\_subset\_of\_estimate}, and
\texttt{SUniformAux.support\_subset\_of\_integral\_inner\_eq\_zero}
as upstream prerequisites. The obsolete statement that every tangent lies in
the original half-space boundary must not enter the graph.
\par\emph{Role: Preiss proof.}
\par Lean declarations: \texttt{Preiss.exists\_hyperplane\_eq\_or\_dimensionalTangent}.
\end{lemma}
\begin{proof}[Planned proof obligations]

For a centred uniform half-space measure $\nu$, there exist $w\ne0$ and a
uniform $\lambda$ supported in $w^\perp$ such that either $\lambda=\nu$
or $\lambda$ is a tangent to $\nu$ at zero, up to the explicitly fixed
positive normalization. The new $w$ need not be $e$.
The proof of GMT's \texttt{SUniformAux.exists\_hyperplane} already has
exactly these two branches. Expose
\texttt{SUniformAux.abs\_inner\_integral\_le},
\texttt{SUniformAux.support\_tangent\_subset\_of\_estimate}, and
\texttt{SUniformAux.support\_subset\_of\_integral\_inner\_eq\_zero}
as upstream prerequisites. The obsolete statement that every tangent lies in
the original half-space boundary must not enter the graph.
\par Corrections and boundary conditions: Author erratum: the new hyperplane need not be the original half-space boundary. Preserve equality-or-tangent lineage.
\par Source: De Lellis Lemma 6.6, Theorem 6.8, pp.57-58,63-65; Corollary 3.9 and corrected Lemma 3.8, pp.15,25-26; author errata.
\end{proof}

\begin{lemma}[Codimension reduction retaining orbit membership]
\label{lem:uniform-reduction}
\lean{Preiss.uniform_dimension_reduction}
\uses{def:unit-uniform, gmt:support-not-full, lem:halfspace-tangent-witness, lem:hyperplane-tangent-lineage, lem:hyperplane-transport-map, lem:pseudotangent-transitive, lem:dimensional-tangent-bounded-orbit, lem:plane-isometric-parametrization}






If $\nu\in\mathcal U_m$ in ambient dimension $k>m$, there is a unit uniform measure on a $(k-1)$-dimensional Euclidean space whose isometric pushforward belongs to the support-centered scaling orbit closure of $\nu$.
\par\emph{Role: Preiss proof.}
\par Lean declarations: \texttt{Preiss.uniform\_dimension\_reduction}.
\end{lemma}
\begin{proof}[Planned proof obligations]

Induct on ambient codimension using A01--A05. At dimension $m$, use
\texttt{mattila\_14\_11}, Hausdorff/Lebesgue normalization and isometric
pushforward to identify the terminal uniform measure with a flat one.
There is no need for analytic support varieties, a full Marstrand integer
dimension theorem, or classification of all uniform measures.
\par Source: De Lellis Lemma 6.6, Theorem 6.8, pp.57-58,63-65; Corollary 3.9 and corrected Lemma 3.8, pp.15,25-26; author errata.
\end{proof}

\begin{lemma}[Classification of plane-supported uniform measures]
\label{lem:plane-supported-uniform-flat}
\lean{Preiss.unitUniform_eq_flat_of_support_subset_plane, Preiss.uniformWithConstant_eq_flat_of_support_subset_plane}
\uses{lem:plane-isometric-parametrization, lem:unit-flat-unit-uniform, lem:uniform-full-dimension-unique, lem:hyperplane-transport-map, lem:uniform-scaling}






If $\nu\in\mathcal U_m$ and $\operatorname{supp}\nu\subset V$ for an $m$-dimensional linear plane $V$, then $\nu=\sigma_V$. More generally a positive uniform ball constant determines the corresponding scalar multiple of $\sigma_V$.
\par\emph{Role: Preiss proof.}
\par Lean declarations: \texttt{Preiss.uniformWithConstant\_eq\_flat\_of\_support\_subset\_plane, Preiss.unitUniform\_eq\_flat\_of\_support\_subset\_plane}.
\end{lemma}
\begin{proof}[Planned proof obligations]

Induct on ambient codimension using A01--A05. At dimension $m$, use
\texttt{mattila\_14\_11}, Hausdorff/Lebesgue normalization and isometric
pushforward to identify the terminal uniform measure with a flat one.
There is no need for analytic support varieties, a full Marstrand integer
dimension theorem, or classification of all uniform measures.
\par Source: De Lellis Lemma 6.6, Theorem 6.8, pp.57-58,63-65; Corollary 3.9 and corrected Lemma 3.8, pp.15,25-26; author errata.
\end{proof}

\begin{theorem}[Flat pseudotangent of a uniform measure]
\label{thm:flat-pseudotangent}
\lean{Preiss.uniform_has_flat_pseudotangent, Preiss.unitUniform_exists_flat_uniformPseudotangent}
\uses{def:flat, def:unit-uniform, lem:plane-supported-uniform-flat, lem:pseudotangent-transitive, lem:uniform-reduction, lem:unit-flat-support-regularity, lem:unit-flat-unit-uniform}


For $0<m\leq n$ and an $m$-uniform measure $\nu$, there are an $m$-plane $V$, support centres $a_i$, and radii $0<r_i\leq1$ such that $D^m_{a_i,r_i}\nu\rightharpoonup\sigma_V$, where $\sigma_V$ is the unit-ball normalized flat measure on $V$. The centres may move; no assertion that $r_i\to0$ is made.
\end{theorem}
\begin{proof}[Planned proof]
Induct on ambient codimension using A01--A05. At dimension $m$, use
\texttt{mattila\_14\_11}, Hausdorff/Lebesgue normalization and isometric
pushforward to identify the terminal uniform measure with a flat one.
There is no need for analytic support varieties, a full Marstrand integer
dimension theorem, or classification of all uniform measures.
\end{proof}

\section{Blow-downs and the conical gap}

\begin{theorem}[Uniqueness and full convergence of the blow-down]
\label{thm:unique-infinity}
\lean{Preiss.unique_tangent_at_infinity, Preiss.blowdown_sequence_converges}
\uses{lem:unit-uniform-infinity-existence, lem:unit-uniform-infinity-unique, lem:uniform-compactness}




Every $\nu\in\mathcal U_m$ admits a unique $\lambda\in\mathcal U_m$ such that $D^m_{0,r}\nu\rightharpoonup\lambda$ as $r\to\infty$. In particular every positive divergent radius sequence has this same limit.
\par\emph{Role: Preiss proof.}
\par Lean declarations: \texttt{Preiss.blowdown\_sequence\_converges, Preiss.unique\_tangent\_at\_infinity}.
\end{theorem}
\begin{proof}[Planned proof obligations]

Use S05 to obtain subsequential uniform limits, M02 and M08 to identify their
Gaussian moments, and M09 to identify the limits. The compactness principle
``every sequence escaping to infinity has a convergent subsequence with the
same limit'' gives full convergence. A separate output is that every positive
sequence $r_i\to\infty$ converges to this same measure. Scaling by a fixed
factor preserves the set of limits at infinity, so S02--S03 imply conicality.
\par Source: De Lellis Lemma 7.2, Definition 7.4, Lemma 7.6, Proposition 7.7, pp.70-84.
\end{proof}

\begin{lemma}[Elementary eigenvalue rigidity from one nearby support point]
\label{lem:conical-eigenvalue-rigidity}
\lean{Preiss.conical_covariance_eq_projection_of_near_eigenvector}
\uses{lem:quadratic-near-basis-rigidity}




For a PSD form $B$ of trace $m$, a support set satisfying $B(x,x)=|x|^2$, and an ordered orthonormal eigenbasis, suppose there is a support point within $\delta$ of the $m$th eigenvector, $0<\delta<1/m$. Then the top $m$ eigenvalues are $1$ and the others $0$; the support is contained in their span.
\par\emph{Role: Preiss proof.}
\par Lean declarations: \texttt{Preiss.conical\_covariance\_eq\_projection\_of\_near\_eigenvector}.
\end{lemma}
\begin{proof}[Planned proof obligations]

Under G02, assume the top-$m$ eigenspace has a support point within distance
$\delta$ of its $m$th unit eigenvector, where $0<\delta<1/m$.
If $\alpha_m<1$, the trace gives
$\alpha_i-1\leq(m-1)(1-\alpha_m)$ for $i<m$ and
$\alpha_i\leq\alpha_m$ for $i\geq m$.
The support identity then gives
$0\leq(1-\alpha_m)((m-1)^2\delta^2-(1-\delta)^2)<0$.
Thus $\alpha_m\geq1$. Trace $m$ forces the first $m$ eigenvalues to be one
and the rest zero. Hence support is contained in $V$ and $\nu=\sigma_V$.
G03--G06 give the conical gap for all positive $m$ without the flawed
two-dimensional shortcut in the original printed proof.
\par Source: De Lellis Corollary 8.4, Lemmas 8.6-8.7, Proposition 8.5, pp.85-94; author errata for the dimension-two shortcut.
\end{proof}

\begin{corollary}[The unique blow-down is conical]
\label{cor:conical}
\lean{Preiss.tangent_at_infinity_conical}
\uses{lem:unit-uniform-infinity-unique, lem:vague-moving-blowups, lem:blowup-composition}




For the unique blow-down $\lambda$ of $\nu\in\mathcal U_m$, $D^m_{0,r}\lambda=\lambda$ for every $r>0$.
\par\emph{Role: Preiss proof.}
\par Lean declarations: \texttt{Preiss.tangent\_at\_infinity\_conical}.
\end{corollary}
\begin{proof}[Planned proof obligations]

Use S05 to obtain subsequential uniform limits, M02 and M08 to identify their
Gaussian moments, and M09 to identify the limits. The compactness principle
``every sequence escaping to infinity has a convergent subsequence with the
same limit'' gives full convergence. A separate output is that every positive
sequence $r_i\to\infty$ converges to this same measure. Scaling by a fixed
factor preserves the set of limits at infinity, so S02--S03 imply conicality.
\par Source: De Lellis Lemma 7.2, Definition 7.4, Lemma 7.6, Proposition 7.7, pp.70-84.
\end{proof}

\begin{theorem}[Uniform conical flatness gap]
\label{thm:gap}
\lean{Preiss.conical_flatness_gap}
\uses{lem:ordered-bilinear-spectrum, lem:covariance-minimizing-plane, lem:gaussian-bilinear-covariance, lem:conical-gaussian-ball-energy, lem:small-energy-no-holes, lem:conical-covariance, lem:conical-eigenvalue-rigidity}




For $1\leq m\leq n$ there is $\varepsilon_{m,n}>0$ such that if $\nu\in\mathcal U_m$ is conical and for some $m$-plane $V$, $\int_{B(0,1)}\operatorname{dist}(z,V)^2\,d\nu\leq\varepsilon_{m,n}$, then $\nu$ is flat.
\par\emph{Role: Preiss proof.}
\par Lean declarations: \texttt{Preiss.conical\_flatness\_gap}.
\end{theorem}
\begin{proof}[Planned proof obligations]

Under G02, assume the top-$m$ eigenspace has a support point within distance
$\delta$ of its $m$th unit eigenvector, where $0<\delta<1/m$.
If $\alpha_m<1$, the trace gives
$\alpha_i-1\leq(m-1)(1-\alpha_m)$ for $i<m$ and
$\alpha_i\leq\alpha_m$ for $i\geq m$.
The support identity then gives
$0\leq(1-\alpha_m)((m-1)^2\delta^2-(1-\delta)^2)<0$.
Thus $\alpha_m\geq1$. Trace $m$ forces the first $m$ eigenvalues to be one
and the rest zero. Hence support is contained in $V$ and $\nu=\sigma_V$.
G03--G06 give the conical gap for all positive $m$ without the flawed
two-dimensional shortcut in the original printed proof.
\par Corrections and boundary conditions: Use the uniform spectral argument for every positive m; no dependence on the erroneous printed m=2 shortcut or unnecessary low-dimensional classification.
\par Source: De Lellis Corollary 8.4, Lemmas 8.6-8.7, Proposition 8.5, pp.85-94; author errata for the dimension-two shortcut.
\end{proof}

\section{Flat-at-infinity rigidity}

\begin{definition}[Normal first-moment vector]
\label{def:rigidity-normal-vector}
\lean{Preiss.IsFirstMomentNormalVector, Preiss.exists_firstMomentNormalVector}
\uses{def:moment-coefficients, def:plane}


The interface passes this data explicitly and records its defining identity in a predicate; existence is a separate theorem.



For $\nu\in\mathcal U_m$ with blow-down $\sigma_V$, there is a unique $b\in V^\perp$ satisfying $b_1^{(1)}(x)=2\langle b,x\rangle$ for every $x$.
\par\emph{Role: Preiss proof.}
\par Lean declarations: \texttt{Preiss.IsFirstMomentNormalVector}.
\end{definition}

\begin{definition}[Normal valued bilinear map from the third coefficient]
\label{def:rigidity-omega}
\lean{Preiss.RigidityOmega, Preiss.IsRigidityOmega}
\uses{def:moment-coefficients, def:plane, def:rigidity-normal-vector}


The interface passes this data explicitly and records its defining identity in a predicate; existence is a separate theorem.




Define the continuous symmetric bilinear map $\omega:V\times V\to V^\perp$ by $\langle\omega(u,v),w\rangle=3b_3^{(2)}(u,v,w)-4\langle u,v\rangle\langle b,w\rangle$ for every normal $w$.
\par\emph{Role: Preiss proof.}
\par Lean declarations: \texttt{Preiss.IsRigidityOmega, Preiss.RigidityOmega}.
\end{definition}

\begin{definition}[Adjusted symmetric quadratic form on the limiting plane]
\label{def:rigidity-adjusted-form}
\lean{Preiss.RigidityAdjustedForm, Preiss.IsRigidityAdjustedForm}
\uses{def:moment-coefficients, def:plane, def:rigidity-normal-vector, def:rigidity-omega}


The interface passes this data explicitly and records its defining identity in a predicate; existence is a separate theorem.




Define the continuous symmetric bilinear form $\widehat b(u,v)=b_2^{(2)}(u,v)+\langle\omega(u,v),b\rangle$ on $V$.
\par\emph{Role: Preiss proof.}
\par Lean declarations: \texttt{Preiss.IsRigidityAdjustedForm, Preiss.RigidityAdjustedForm}.
\end{definition}

\begin{lemma}[Even Taylor coefficients from a flat blow-down]
\label{lem:flat-infinity-moments}
\lean{Preiss.evenMomentCoefficient_eq_projection_norm}
\uses{lem:blowdown-even-moment-limit, lem:vague-gaussian-moments, lem:flat-gaussian-even-moments}







For every $j\geq1$, $b_{2j}^{(j)}(x^{2j})=|Px|^{2j}$ for every $x$.
Use M02, M03 and the Gaussian moment formula of the flat limit, including the
factorial from M08.
\par\emph{Role: Preiss proof.}
\par Lean declarations: \texttt{Preiss.evenMomentCoefficient\_eq\_projection\_norm}.
\end{lemma}
\begin{proof}[Planned proof obligations]

For every $j\geq1$, $b_{2j}^{(j)}(x^{2j})=|Px|^{2j}$ for every $x$.
Use M02, M03 and the Gaussian moment formula of the flat limit, including the
factorial from M08.
\par Source: De Lellis Proposition 9.1, Lemmas 9.2,9.4,9.9,9.10 and Proposition 9.5, pp.95-109.
\end{proof}

\begin{lemma}[Odd top coefficients vanish on the flat limiting plane]
\label{lem:odd-coefficients-plane-zero}
\lean{Preiss.oddMomentCoefficient_restrict_plane_eq_zero}
\uses{lem:flat-infinity-moments, lem:moment-support-polynomial, lem:vague-support-approximation, lem:dimensional-blowup-support, lem:leading-homogeneous-term, lem:symmetric-multilinear-diagonal-injective}







For $j\geq1$, the restriction of $b_{2j-1}^{(j)}$ to $V$ is zero.
For a nonzero $v\in V$, use the flat blow-down and S04 to find support points
$x_i$ with $|x_i|\to\infty$ and $x_i/|x_i|\to v/|v|$.
The support polynomial identity, R01 and the fact that all remaining terms have
lower degree show $b_{2j-1}^{(j)}(v^{2j-1})\geq0$.
Repeat with $-v$ and use polarization to obtain the full multilinear restriction.
\par\emph{Role: Preiss proof.}
\par Lean declarations: \texttt{Preiss.oddMomentCoefficient\_restrict\_plane\_eq\_zero}.
\end{lemma}
\begin{proof}[Planned proof obligations]

For $j\geq1$, the restriction of $b_{2j-1}^{(j)}$ to $V$ is zero.
For a nonzero $v\in V$, use the flat blow-down and S04 to find support points
$x_i$ with $|x_i|\to\infty$ and $x_i/|x_i|\to v/|v|$.
The support polynomial identity, R01 and the fact that all remaining terms have
lower degree show $b_{2j-1}^{(j)}(v^{2j-1})\geq0$.
Repeat with $-v$ and use polarization to obtain the full multilinear restriction.
\par Source: De Lellis Proposition 9.1, Lemmas 9.2,9.4,9.9,9.10 and Proposition 9.5, pp.95-109.
\end{proof}

\begin{lemma}[Cylinder identity and bounded normal coordinates]
\label{lem:uniform-cylinder-identity}
\lean{Preiss.support_normal_sq_eq_inner}
\uses{lem:moment-support-polynomial, lem:flat-infinity-moments, lem:normal-coordinate-identities}








With $Q=1-P_V$, every $x\in\operatorname{supp}\nu$ satisfies $|Qx|^2=2\langle b,x\rangle$, $|Qx-b|=|b|$, and $|Qx|\leq2|b|$.
\par\emph{Role: Preiss proof.}
\par Lean declarations: \texttt{Preiss.support\_normal\_sq\_eq\_inner}.
\end{lemma}
\begin{proof}[Planned proof obligations]

Riesz representation and R02 give $b\in V^\perp$ with
$b_1^{(1)}(x)=2\langle b,x\rangle$. The first support identity gives
$|Qx|^2=2\langle b,x\rangle$ on the support, and hence $|Qx|\leq2|b|$.
Also $|Qx-b|=|b|$ there. Keep these identities as named lemmas.
\par Source: De Lellis Proposition 9.1, Lemmas 9.2,9.4,9.9,9.10 and Proposition 9.5, pp.95-109.
\end{proof}

\begin{lemma}[Mixed Gaussian moment limit defining omega]
\label{lem:rigidity-mixed-gaussian-limit}
\lean{Preiss.mixed_gaussian_limit_eq_omega}
\uses{lem:mixed-centered-gaussian-limit}







For all ambient $v,w$, $\lim_{s\downarrow0}\frac{8s}{I(s)}\int e^{-s|z|^2}\langle z,v\rangle^2\langle z-b,w\rangle\,d\nu=3b_3^{(2)}(v,v,w)-4b_2^{(1)}(v,v)\langle b,w\rangle$. For $v\in V,w\in V^\perp$ this equals $\langle\omega(v,v),w\rangle$.
\par\emph{Role: Preiss proof.}
\par Lean declarations: \texttt{Preiss.mixed\_gaussian\_limit\_eq\_omega}.
\end{lemma}
\begin{proof}[Planned proof obligations]

For $v,w\in\mathbb R^n$, prove
\[
 \lim_{s\downarrow0}\frac{8s}{I(s)}\int e^{-s|z|^2}
 \langle z,v\rangle^2\langle z-b,w\rangle\,d\nu
 =3b_3^{(2)}(v,v,w)-4b_2^{(1)}(v,v)\langle b,w\rangle.
\]
For $v\in V,w\in V^\perp$ this is $\langle\omega(v,v),w\rangle$.
Further prove
\[
 \lim_{s\downarrow0}\frac{s}{I(s)}\int e^{-s|z|^2}
 |Pz|^2\langle z-b,w\rangle\,d\nu=0.
\]
To justify the latter, the analogous $|z|^2$ expression has a limit by the
first formula and a finite orthonormal-basis sum. Differentiating the Gaussian
first moment and applying the appropriate one-sided l'Hopital argument
identifies this limit as zero; the $|Qz|^2$ term tends to zero by R03 and M02.
Record differentiability and domination hypotheses before using l'Hopital.
\par Source: De Lellis Proposition 9.1, Lemmas 9.2,9.4,9.9,9.10 and Proposition 9.5, pp.95-109.
\end{proof}

\begin{lemma}[Normal second Taylor trace]
\label{lem:moment-normal-trace}
\lean{Preiss.trace_secondMomentCoefficient_normal_eq}
\uses{lem:cylinder-second-coefficient-trace, lem:uniform-cylinder-identity}






$\operatorname{tr}(b_2^{(2)}|V^\perp)=8|b|^2=2\|b_1^{(1)}\|^2$.
\par\emph{Role: Preiss proof.}
\par Lean declarations: \texttt{Preiss.trace\_secondMomentCoefficient\_normal\_eq}.
\end{lemma}
\begin{proof}[Planned proof obligations]

The exact radial identity $\operatorname{tr}b_{2,s}=ms$ and M07 imply
$\operatorname{tr}b_2^{(2)}=0$.
R03 and the first-moment expansion give
$\operatorname{tr}(b_2^{(2)}|V^\perp)=8|b|^2
=2\|b_1^{(1)}\|^2$.
These are separate, comparatively short lemmas.
\par Source: De Lellis Proposition 9.1, Lemmas 9.2,9.4,9.9,9.10 and Proposition 9.5, pp.95-109.
\end{proof}

\begin{lemma}[Bounded distance from the limiting plane to the support]
\label{lem:flat-infinity-bounded-holes}
\lean{Preiss.exists_uniform_dist_plane_support_bound}
\uses{lem:rigidity-normal-vector-exists, lem:uniform-cylinder-identity, lem:bounded-tube-plane-to-support}







There exists $R_0<\infty$ with
$\operatorname{dist}(v,\operatorname{supp}\nu)<R_0$ for every $v\in V$.
If not, choose $v_i$ with distances $r_i\to\infty$, a nearest support point,
and its projection $z_i\in V$. The measures $D^m_{z_i,r_i}\nu$ have
supports approaching the origin and lie in shrinking tubes about $V$ by R03.
S05 gives a unit uniform limit supported in $V$, hence $\sigma_V$.
The original holes give a unit hole centred at a bounded point of $V$, a
contradiction. The centres $z_i$ need not lie in the support, which is why the
near-support variant of S05 is essential.
\par\emph{Role: Preiss proof.}
\par Lean declarations: \texttt{Preiss.exists\_uniform\_dist\_plane\_support\_bound}.
\end{lemma}
\begin{proof}[Planned proof obligations]

There exists $R_0<\infty$ with
$\operatorname{dist}(v,\operatorname{supp}\nu)<R_0$ for every $v\in V$.
If not, choose $v_i$ with distances $r_i\to\infty$, a nearest support point,
and its projection $z_i\in V$. The measures $D^m_{z_i,r_i}\nu$ have
supports approaching the origin and lie in shrinking tubes about $V$ by R03.
S05 gives a unit uniform limit supported in $V$, hence $\sigma_V$.
The original holes give a unit hole centred at a bounded point of $V$, a
contradiction. The centres $z_i$ need not lie in the support, which is why the
near-support variant of S05 is essential.
\par Corrections and boundary conditions: The projected rescaling centers need not be in support; invoke the near-support compactness variant.
\par Source: De Lellis Proposition 9.1, Lemmas 9.2,9.4,9.9,9.10 and Proposition 9.5, pp.95-109.
\end{proof}

\begin{lemma}[Vanishing normal mixed cubic coefficients]
\label{lem:rigidity-cubic-vanishing}
\lean{Preiss.thirdMomentCoefficient_normal_vanishes}
\uses{lem:bounded-tube-cubic-vanishing, lem:uniform-cylinder-identity, lem:odd-coefficients-plane-zero, lem:flat-infinity-moments, lem:symmetric-cubic-binomial, lem:moment-support-polynomial}






For $v\in V,w\in V^\perp$, $b_3^{(2)}(v,v,v)=b_3^{(2)}(v,w,w)=b_3^{(2)}(w,w,w)=0$. Polarization gives all corresponding mixed restrictions.
\par\emph{Role: Preiss proof.}
\par Lean declarations: \texttt{Preiss.thirdMomentCoefficient\_normal\_vanishes}.
\end{lemma}
\begin{proof}[Planned proof obligations]

For $v\in V,w\in V^\perp$,
$b_3^{(2)}(v,v,v)=b_3^{(2)}(v,w,w)=b_3^{(2)}(w,w,w)=0$.
The first is R02; the others follow by R03's bounded normal coordinates and
Gaussian estimates in the defining limits. Therefore, on the support,
\[
 b_1^{(2)}(z)+b_2^{(2)}(z,z)+3b_3^{(2)}(Pz,Pz,Qz)
 =|Qz|^4+2|Pz|^2|Qz|^2.
\]
Use the order-two support identity and R01; do not drop mixed cubic terms
without first establishing these vanishings.
\par Source: De Lellis Proposition 9.1, Lemmas 9.2,9.4,9.9,9.10 and Proposition 9.5, pp.95-109.
\end{proof}

\begin{lemma}[Vanishing contracted Gaussian limit]
\label{lem:rigidity-contracted-limit-zero}
\lean{Preiss.contracted_gaussian_limit_eq_zero}
\uses{lem:bounded-tube-projected-contraction, lem:uniform-cylinder-identity}






For every normal $w$, $\lim_{s\downarrow0}\frac{s}{I(s)}\int e^{-s|z|^2}|P_Vz|^2\langle z-b,w\rangle\,d\nu=0$.
\par\emph{Role: Preiss proof.}
\par Lean declarations: \texttt{Preiss.contracted\_gaussian\_limit\_eq\_zero}.
\end{lemma}
\begin{proof}[Planned proof obligations]

For $v,w\in\mathbb R^n$, prove
\[
 \lim_{s\downarrow0}\frac{8s}{I(s)}\int e^{-s|z|^2}
 \langle z,v\rangle^2\langle z-b,w\rangle\,d\nu
 =3b_3^{(2)}(v,v,w)-4b_2^{(1)}(v,v)\langle b,w\rangle.
\]
For $v\in V,w\in V^\perp$ this is $\langle\omega(v,v),w\rangle$.
Further prove
\[
 \lim_{s\downarrow0}\frac{s}{I(s)}\int e^{-s|z|^2}
 |Pz|^2\langle z-b,w\rangle\,d\nu=0.
\]
To justify the latter, the analogous $|z|^2$ expression has a limit by the
first formula and a finite orthonormal-basis sum. Differentiating the Gaussian
first moment and applying the appropriate one-sided l'Hopital argument
identifies this limit as zero; the $|Qz|^2$ term tends to zero by R03 and M02.
Record differentiability and domination hypotheses before using l'Hopital.
\par Corrections and boundary conditions: Establish the derivative-quotient limit first from mixed moments, then apply the one-sided lHopital argument; record domination for Gaussian differentiation.
\par Source: De Lellis Proposition 9.1, Lemmas 9.2,9.4,9.9,9.10 and Proposition 9.5, pp.95-109.
\end{proof}

\begin{lemma}[Fourth degree support identity]
\label{lem:rigidity-fourth-support-identity}
\lean{Preiss.fourth_order_support_identity}
\uses{lem:bounded-tube-cubic-vanishing, lem:uniform-cylinder-identity, lem:odd-coefficients-plane-zero, lem:flat-infinity-moments, lem:symmetric-cubic-binomial, lem:moment-support-polynomial}






For every $z\in\operatorname{supp}\nu$, $b_1^{(2)}(z)+b_2^{(2)}(z,z)+3b_3^{(2)}(Pz,Pz,Qz)=|Qz|^4+2|Pz|^2|Qz|^2$.
\par\emph{Role: Preiss proof.}
\par Lean declarations: \texttt{Preiss.fourth\_order\_support\_identity}.
\end{lemma}
\begin{proof}[Planned proof obligations]

For $v\in V,w\in V^\perp$,
$b_3^{(2)}(v,v,v)=b_3^{(2)}(v,w,w)=b_3^{(2)}(w,w,w)=0$.
The first is R02; the others follow by R03's bounded normal coordinates and
Gaussian estimates in the defining limits. Therefore, on the support,
\[
 b_1^{(2)}(z)+b_2^{(2)}(z,z)+3b_3^{(2)}(Pz,Pz,Qz)
 =|Qz|^4+2|Pz|^2|Qz|^2.
\]
Use the order-two support identity and R01; do not drop mixed cubic terms
without first establishing these vanishings.
\par Source: De Lellis Proposition 9.1, Lemmas 9.2,9.4,9.9,9.10 and Proposition 9.5, pp.95-109.
\end{proof}

\begin{lemma}[Zero Gaussian mean of omega]
\label{lem:rigidity-omega-trace-zero}
\lean{Preiss.integral_rigidityOmega_eq_zero}
\uses{def:moment-coefficients, def:plane, def:rigidity-normal-vector, def:rigidity-omega, def:tangent-at-infinity, def:unit-uniform, lem:flat-gaussian-even-moments, lem:rigidity-contracted-limit-zero, lem:rigidity-mixed-gaussian-limit}






For standard covariance-one Gaussian probability $\gamma_V$ on $V$, $\int\omega(v,v)\,d\gamma_V(v)=0$. Equivalently its inner product with every fixed normal vector has integral zero.
\par\emph{Role: Preiss proof.}
\par Lean declarations: \texttt{Preiss.integral\_rigidityOmega\_eq\_zero}.
\end{lemma}
\begin{proof}[Planned proof obligations]

Let $\gamma$ be the standard covariance-one Gaussian probability measure on
$V$. Prove $\int\langle\omega(v,v),b\rangle\,d\gamma(v)=0$ using R07,
Fubini and domination by a constant times $|v|^2$.
Then
$\operatorname{tr}(b_2^{(2)}|V)=\int\widehat b(v,v)\,d\gamma(v)$.
The same argument works against each fixed normal vector, not only $b$.
\par Source: De Lellis Proposition 9.1, Lemmas 9.2,9.4,9.9,9.10 and Proposition 9.5, pp.95-109.
\end{proof}

\begin{lemma}[Pointwise square bound for the adjusted form]
\label{lem:rigidity-pointwise-quadratic-bound}
\lean{Preiss.adjustedForm_sq_le_norm_omega_sq}
\uses{lem:rigidity-pointwise-from-growth, lem:rigidity-error-linear-growth}







For all $v\in V$,
$\widehat b(v,v)^2\leq |b|^2|\omega(v,v)|^2$.
Choose support points $t_i v+v_i+w_i$ within bounded distance of $t_i v$ using
R05. Pass to a convergent subsequence of bounded normal components $w_i$.
Divide R09 by $t_i^2$ and use R03 to obtain
$b_2^{(2)}(v,v)+\langle\omega(v,v),w\rangle=0$ and
$|w|^2=2\langle b,w\rangle$.
Thus $|w-b|=|b|$ and
$\widehat b(v,v)+\langle\omega(v,v),w-b\rangle=0$; Cauchy--Schwarz proves
the inequality.
\par\emph{Role: Preiss proof.}
\par Lean declarations: \texttt{Preiss.adjustedForm\_sq\_le\_norm\_omega\_sq}.
\end{lemma}
\begin{proof}[Planned proof obligations]

For all $v\in V$,
$\widehat b(v,v)^2\leq |b|^2|\omega(v,v)|^2$.
Choose support points $t_i v+v_i+w_i$ within bounded distance of $t_i v$ using
R05. Pass to a convergent subsequence of bounded normal components $w_i$.
Divide R09 by $t_i^2$ and use R03 to obtain
$b_2^{(2)}(v,v)+\langle\omega(v,v),w\rangle=0$ and
$|w|^2=2\langle b,w\rangle$.
Thus $|w-b|=|b|$ and
$\widehat b(v,v)+\langle\omega(v,v),w-b\rangle=0$; Cauchy--Schwarz proves
the inequality.
\par Source: De Lellis Proposition 9.1, Lemmas 9.2,9.4,9.9,9.10 and Proposition 9.5, pp.95-109.
\end{proof}

\begin{lemma}[Vanishing lower order Gaussian error]
\label{lem:rigidity-rescaled-error}
\lean{Preiss.adjustedForm_gaussian_error_tendsto_zero}
\uses{lem:rigidity-error-limit-from-growth, lem:rigidity-error-linear-growth}







R09 and R03 imply that
$\widehat b(Pz,Pz)+\langle\omega(Pz,Pz),Qz-b\rangle$
has absolute value at most $C(1+|z|)$ on the support.
Hence its integral against $sI(s)^{-1}e^{-s|z|^2}\nu$ tends to zero.
This is equation (9.34); it is a necessary bridge to the global trace bound.
\par\emph{Role: Preiss proof.}
\par Lean declarations: \texttt{Preiss.adjustedForm\_gaussian\_error\_tendsto\_zero, Preiss.rigidityError, Preiss.rigidityError\_le\_linear\_growth}.
\end{lemma}
\begin{proof}[Planned proof obligations]

R09 and R03 imply that
$\widehat b(Pz,Pz)+\langle\omega(Pz,Pz),Qz-b\rangle$
has absolute value at most $C(1+|z|)$ on the support.
Hence its integral against $sI(s)^{-1}e^{-s|z|^2}\nu$ tends to zero.
This is equation (9.34); it is a necessary bridge to the global trace bound.
\par Source: De Lellis Proposition 9.1, Lemmas 9.2,9.4,9.9,9.10 and Proposition 9.5, pp.95-109.
\end{proof}

\begin{lemma}[Tangential trace is the mean adjusted form]
\label{lem:rigidity-gaussian-trace}
\lean{Preiss.trace_tangent_eq_integral_adjustedForm}
\uses{def:moment-coefficients, def:plane, def:rigidity-adjusted-form, def:rigidity-normal-vector, def:rigidity-omega, def:tangent-at-infinity, def:unit-uniform, lem:flat-gaussian-even-moments, lem:rigidity-omega-trace-zero, def:bilinear-trace}







$\operatorname{tr}(b_2^{(2)}|V)=\int\widehat b(v,v)\,d\gamma_V(v)$.
\par\emph{Role: Preiss proof.}
\par Lean declarations: \texttt{Preiss.trace\_tangent\_eq\_integral\_adjustedForm}.
\end{lemma}
\begin{proof}[Planned proof obligations]

Let $\gamma$ be the standard covariance-one Gaussian probability measure on
$V$. Prove $\int\langle\omega(v,v),b\rangle\,d\gamma(v)=0$ using R07,
Fubini and domination by a constant times $|v|^2$.
Then
$\operatorname{tr}(b_2^{(2)}|V)=\int\widehat b(v,v)\,d\gamma(v)$.
The same argument works against each fixed normal vector, not only $b$.
\par Source: De Lellis Proposition 9.1, Lemmas 9.2,9.4,9.9,9.10 and Proposition 9.5, pp.95-109.
\end{proof}

\begin{lemma}[Gaussian reproducing identity for omega]
\label{lem:rigidity-omega-reproducing}
\lean{Preiss.integral_inner_sq_mul_omega}
\uses{def:moment-coefficients, def:plane, def:rigidity-normal-vector, def:rigidity-omega, def:tangent-at-infinity, def:unit-uniform, lem:gaussian-wick-fourth, lem:rigidity-omega-trace-zero}







For $y\in V,w\in V^\perp$, $\int\langle y,v\rangle^2\langle\omega(v,v),w\rangle\,d\gamma_V(v)=2\langle\omega(y,y),w\rangle$.
\par\emph{Role: Preiss proof.}
\par Lean declarations: \texttt{Preiss.integral\_inner\_sq\_mul\_omega}.
\end{lemma}
\begin{proof}[Planned proof obligations]

For $y,z\in V$,
$\int\langle y,v\rangle^2\langle z,v\rangle^2\,d\gamma
=|y|^2|z|^2+2\langle y,z\rangle^2$.
Combining this with R07 yields, for $w\in V^\perp$,
\[
 \int\langle y,v\rangle^2\langle\omega(v,v),w\rangle\,d\gamma(v)
 =2\langle\omega(y,y),w\rangle.
\]
Justify limits by Gaussian quartic domination and Fubini.
\par Source: De Lellis Proposition 9.1, Lemmas 9.2,9.4,9.9,9.10 and Proposition 9.5, pp.95-109.
\end{proof}

\begin{lemma}[Integral quadratic inequality]
\label{lem:rigidity-integral-quadratic-bound}
\lean{Preiss.integral_adjustedForm_sq_le}
\uses{def:moment-coefficients, def:plane, def:rigidity-adjusted-form, def:rigidity-normal-vector, def:rigidity-omega, def:tangent-at-infinity, def:unit-uniform, lem:rigidity-mixed-gaussian-limit, lem:rigidity-omega-reproducing, lem:rigidity-pointwise-quadratic-bound, lem:rigidity-rescaled-error, lem:vague-gaussian-polynomials}






Integrate R10, express $\int|\omega(v,v)|^2\,d\gamma$ through R07 and R12,
use R11, and pass to the flat blow-down using M02. The result is
\[
 \int\widehat b(v,v)^2\,d\gamma(v)
 \leq-8|b|^2\int\widehat b(v,v)\,d\gamma(v).
\]
The inequality must remain an inequality throughout. The equality sign at
the start of the second long display on printed page 108 loses the earlier
inequality; its final statement (9.41) has the correct direction.
\par\emph{Role: Preiss proof.}
\par Lean declarations: \texttt{Preiss.integral\_adjustedForm\_sq\_le}.
\end{lemma}
\begin{proof}[Planned proof obligations]

Integrate R10, express $\int|\omega(v,v)|^2\,d\gamma$ through R07 and R12,
use R11, and pass to the flat blow-down using M02. The result is
\[
 \int\widehat b(v,v)^2\,d\gamma(v)
 \leq-8|b|^2\int\widehat b(v,v)\,d\gamma(v).
\]
The inequality must remain an inequality throughout. The equality sign at
the start of the second long display on printed page 108 loses the earlier
inequality; its final statement (9.41) has the correct direction.
\par Corrections and boundary conditions: Preserve the inequality through the calculation; the p.108 intermediate equality is too strong.
\par Source: De Lellis Proposition 9.1, Lemmas 9.2,9.4,9.9,9.10 and Proposition 9.5, pp.95-109.
\end{proof}

\begin{lemma}[Positive lower bound for the total second coefficient trace]
\label{lem:trace-rigidity}
\lean{Preiss.integral_adjustedForm_lower_bound, Preiss.secondMomentCoefficient_trace_lower_bound}
\uses{def:bilinear-trace, def:moment-coefficients, def:plane, def:rigidity-adjusted-form, def:rigidity-normal-vector, def:tangent-at-infinity, def:unit-uniform, lem:bilinear-trace-basis, lem:gaussian-quadratic-trace-bound, lem:moment-normal-trace, lem:rigidity-gaussian-trace, lem:rigidity-integral-quadratic-bound}






For a unit uniform measure with flat blow-down, $\operatorname{tr}b_2^{(2)}\geq16|b|^2/(m+2)$. More precisely $\int\widehat b(v,v)\,d\gamma_V\geq-8m|b|^2/(m+2)$.
\par\emph{Role: Preiss proof.}
\par Lean declarations: \texttt{Preiss.integral\_adjustedForm\_lower\_bound, Preiss.secondMomentCoefficient\_trace\_lower\_bound}.
\end{lemma}
\begin{proof}[Planned proof obligations]

For every symmetric bilinear form $A$ on an $m$-dimensional Euclidean space,
\[
 \int A(v,v)^2\,d\gamma=(\operatorname{tr}A)^2+
 2\|A\|_{\mathrm{HS}}^2
 \geq(1+2/m)\left(\int A(v,v)\,d\gamma\right)^2.
\]
This is a reusable Gaussian/spectral theorem independent of uniform measures.
Combining it with R13 gives
$\int\widehat b(v,v)\,d\gamma\geq-8m|b|^2/(m+2)$.
R08 and R04 therefore give
$\operatorname{tr}b_2^{(2)}\geq16|b|^2/(m+2)$.
Since that trace is zero, $b=0$; R03 makes the support planar, and plane-supported
uniform classification finishes rigidity. This is the actual content hidden
inside the current single trace-rigidity node.
\par Source: De Lellis Proposition 9.1, Lemmas 9.2,9.4,9.9,9.10 and Proposition 9.5, pp.95-109.
\end{proof}

\begin{theorem}[Flatness follows from flatness at infinity]
\label{thm:infinity-rigidity}
\lean{Preiss.flat_of_flat_at_infinity}
\uses{def:flat, def:plane, def:tangent-at-infinity, def:unit-uniform, lem:moment-second-total-trace, lem:plane-supported-uniform-flat, lem:rigidity-forms-exist, lem:rigidity-normal-vector-exists, lem:trace-rigidity, lem:uniform-cylinder-identity}







If $\nu\in\mathcal U_m$ has a flat tangent at infinity, then $\nu$ is flat: the trace bounds force $b=0$, the cylinder identity forces support in $V$, and plane-supported uniform classification gives $\nu=\sigma_V$.
\par\emph{Role: Preiss proof.}
\par Lean declarations: \texttt{Preiss.flat\_of\_flat\_at\_infinity}.
\end{theorem}
\begin{proof}[Planned proof obligations]

For every symmetric bilinear form $A$ on an $m$-dimensional Euclidean space,
\[
 \int A(v,v)^2\,d\gamma=(\operatorname{tr}A)^2+
 2\|A\|_{\mathrm{HS}}^2
 \geq(1+2/m)\left(\int A(v,v)\,d\gamma\right)^2.
\]
This is a reusable Gaussian/spectral theorem independent of uniform measures.
Combining it with R13 gives
$\int\widehat b(v,v)\,d\gamma\geq-8m|b|^2/(m+2)$.
R08 and R04 therefore give
$\operatorname{tr}b_2^{(2)}\geq16|b|^2/(m+2)$.
Since that trace is zero, $b=0$; R03 makes the support planar, and plane-supported
uniform classification finishes rigidity. This is the actual content hidden
inside the current single trace-rigidity node.
\par Source: De Lellis Proposition 9.1, Lemmas 9.2,9.4,9.9,9.10 and Proposition 9.5, pp.95-109.
\end{proof}

\section{Preiss theorem assembly}

\begin{theorem}[Existence of a flat tangent almost everywhere]
\label{thm:part-a}
\lean{Preiss.ae_exists_flat_tangent}
\uses{lem:uniform-adapter, lem:recenter, lem:radius-power-density, lem:dimensional-tangent-compactness, lem:dimensional-tangent-ball-constant, lem:flat-tangent-from-uniform-recentering, lem:unit-flat-support-regularity, lem:flat-positive-scalar}







Intersect the a.e. sets for uniformity, density and A09. Choose any tangent,
use A06, recenter every approximating term using A09, and use S07 for the
nonzero flat limit. This is \texttt{ae\_exists\_flat\_tangent}.
\par\emph{Role: Preiss proof.}
\par Lean declarations: \texttt{Preiss.ae\_exists\_flat\_tangent}.
\end{theorem}
\begin{proof}[Planned proof obligations]

Intersect the a.e. sets for uniformity, density and A09. Choose any tangent,
use A06, recenter every approximating term using A09, and use S07 for the
nonzero flat limit. This is \texttt{ae\_exists\_flat\_tangent}.
\par Source: De Lellis Lemma 6.6, Theorem 6.8, pp.57-58,63-65; Corollary 3.9 and corrected Lemma 3.8, pp.15,25-26; author errata.
\end{proof}

\begin{lemma}[First-crossing sequences and their uniform tangent limit]
\label{lem:first-crossing}
\lean{Preiss.exists_first_crossing_sequences}
\uses{cor:conical, def:dimensional-tangent, def:first-crossing-data, def:flat, def:flatness, def:flatness-gap, def:positive-density, def:tangent-at-infinity, def:unit-uniform, lem:density-normalization, lem:dimensional-tangent-compactness, lem:flatness-continuous, lem:flatness-minimum, lem:flatness-radius-continuous, lem:greatest-crossing, lem:tangent-blowdown, thm:gap, thm:infinity-rigidity, thm:unique-infinity, def:flatness-cutoff}




After S08 normalization, suppose $\alpha$ is a flat dimensional tangent and
$\zeta$ is a nonflat dimensional tangent, all dimensional tangents being unit
uniform. Let $\chi$ be the blow-down of $\zeta$. Rigidity and the gap give
$F(\chi)>\varepsilon$. Choose $r_i\to0$ with blow-ups tending to $\alpha$
and $s_i<r_i$ with blow-ups tending to $\chi$; subsequence selection makes
the ordering possible. For sufficiently large $i$, the continuous scale
function satisfies $F(D_{x,r_i}^m\mu)<\varepsilon$ and
$F(D_{x,s_i}^m\mu)>\varepsilon$.
Choose $\sigma_i\in[s_i,r_i]$ with energy exactly $\varepsilon$ and energy
at most $\varepsilon$ on the entire interval $[\sigma_i,r_i]$.
This is a greatest-crossing/last-crossing assertion, not merely the intermediate
value theorem. Extract a nonzero unit uniform limit $\xi$ at the $\sigma_i$.
\par\emph{Role: Preiss proof.}
\par Lean declarations: \texttt{Preiss.exists\_first\_crossing\_sequences}.
\end{lemma}
\begin{proof}[Planned proof obligations]

After S08 normalization, suppose $\alpha$ is a flat dimensional tangent and
$\zeta$ is a nonflat dimensional tangent, all dimensional tangents being unit
uniform. Let $\chi$ be the blow-down of $\zeta$. Rigidity and the gap give
$F(\chi)>\varepsilon$. Choose $r_i\to0$ with blow-ups tending to $\alpha$
and $s_i<r_i$ with blow-ups tending to $\chi$; subsequence selection makes
the ordering possible. For sufficiently large $i$, the continuous scale
function satisfies $F(D_{x,r_i}^m\mu)<\varepsilon$ and
$F(D_{x,s_i}^m\mu)>\varepsilon$.
Choose $\sigma_i\in[s_i,r_i]$ with energy exactly $\varepsilon$ and energy
at most $\varepsilon$ on the entire interval $[\sigma_i,r_i]$.
This is a greatest-crossing/last-crossing assertion, not merely the intermediate
value theorem. Extract a nonzero unit uniform limit $\xi$ at the $\sigma_i$.
\par Source: De Lellis Lemma 6.20 and Section 6.4, pp.66-69.
\end{proof}

\begin{lemma}[Divergence of the crossing scale ratios]
\label{lem:first-crossing-ratio}
\lean{Preiss.first_crossing_ratio_tendsto_atTop}
\uses{lem:vague-limit-unique, lem:vague-moving-blowups, lem:plane-scaling}




Prove $r_i/\sigma_i\to\infty$. If the ratios have a bounded subsequence,
extract a further limit $R\in[1,\infty)$. S02--S03 give
$D^m_{0,R}\xi=\alpha$. Invert the dilation to conclude $\xi$ flat,
contradicting $F(\xi)=\varepsilon>0$.
\par\emph{Role: Preiss proof.}
\par Lean declarations: \texttt{Preiss.first\_crossing\_ratio\_tendsto\_atTop}.
\end{lemma}
\begin{proof}[Planned proof obligations]

Prove $r_i/\sigma_i\to\infty$. If the ratios have a bounded subsequence,
extract a further limit $R\in[1,\infty)$. S02--S03 give
$D^m_{0,R}\xi=\alpha$. Invert the dilation to conclude $\xi$ flat,
contradicting $F(\xi)=\varepsilon>0$.
\par Source: De Lellis Lemma 6.20 and Section 6.4, pp.66-69.
\end{proof}

\begin{lemma}[Crossing limit has flat blow-down]
\label{lem:first-crossing-flat-infinity}
\lean{Preiss.first_crossing_limit_flat_at_infinity}
\uses{cor:conical, lem:unit-uniform-infinity-existence, lem:flatness-continuous}




For the crossing limit $\xi$, every $R\geq1$ satisfies $F(D^m_{0,R}\xi)\leq\varepsilon$. Hence its unique conical blow-down $\psi$ has $F(\psi)\leq\varepsilon$ and is flat by the conical gap.
\par\emph{Role: Preiss proof.}
\par Lean declarations: \texttt{Preiss.first\_crossing\_limit\_flat\_at\_infinity}.
\end{lemma}
\begin{proof}[Planned proof obligations]

For every fixed $R\geq1$, eventually $\sigma_i R\in[\sigma_i,r_i]$.
S02--S03 and C04 give $F(D^m_{0,R}\xi)\leq\varepsilon$.
Let $R\to\infty$. M10 and continuity give $F(\psi)\leq\varepsilon$ for
the conical blow-down $\psi$ of $\xi$. Apply the conical gap, then rigidity,
to conclude $\xi$ flat, the final contradiction.
Keep the centre $x$ fixed throughout; the printed scale function sometimes
uses zero where $x$ is intended. Every use of $F$ on a dilated measure must
retain the factor $R^{-m}$.
\par Corrections and boundary conditions: Retain the fixed original center x and the full dimensional mass factor R\textasciicircum{}\{-m\}.
\par Source: De Lellis Lemma 6.20 and Section 6.4, pp.66-69.
\end{proof}

\begin{theorem}[One flat tangent implies all tangents are flat]
\label{thm:part-b}
\lean{Preiss.all_tangents_flat_of_one}
\uses{def:dimensional-tangent, def:flat, def:positive-density, lem:density-normalization, lem:first-crossing, lem:first-crossing-flat-infinity, lem:flatness-cutoff-exists, lem:flatness-gap-cutoff, lem:tangent-ray, thm:infinity-rigidity, lem:radius-power-density}







At a point $x$ with positive finite existing $m$-density, suppose every tangent is $m$-uniform and at least one tangent is flat. Then every tangent at $x$ is flat. First use the common density normalization for dimensional tangents, then the first-crossing contradiction, and finally transfer to arbitrary-scalar tangents by the ray adapter.
\par\emph{Role: Preiss proof.}
\par Lean declarations: \texttt{Preiss.all\_tangents\_flat\_of\_one}.
\end{theorem}
\begin{proof}[Planned proof obligations]

For every fixed $R\geq1$, eventually $\sigma_i R\in[\sigma_i,r_i]$.
S02--S03 and C04 give $F(D^m_{0,R}\xi)\leq\varepsilon$.
Let $R\to\infty$. M10 and continuity give $F(\psi)\leq\varepsilon$ for
the conical blow-down $\psi$ of $\xi$. Apply the conical gap, then rigidity,
to conclude $\xi$ flat, the final contradiction.
Keep the centre $x$ fixed throughout; the printed scale function sometimes
uses zero where $x$ is intended. Every use of $F$ on a dilated measure must
retain the factor $R^{-m}$.
\par Source: De Lellis Lemma 6.20 and Section 6.4, pp.66-69.
\end{proof}

\begin{theorem}[Preiss]
\label{thm:preiss}
\lean{Preiss.preiss}
\uses{def:positive-density, def:rectifiable, lem:excess-dimension, lem:uniform-adapter, lem:zero-dimension, lem:zero-measure, thm:mm-measure, thm:part-a, thm:part-b}










Let $m,n\in\N$ and let $\mu$ be a locally finite Borel measure on $\R^n$.
If $0<\Theta_*^m(\mu,x)=\Theta^{*m}(\mu,x)<\infty$ for
$\mu$-almost every $x$, then $\mu$ is an $m$-rectifiable measure.
\end{theorem}
\begin{proof}[Assembly, still dependent on unproved obligations]

For $0<m\leq n$, intersect the finitely many full-measure sets, apply
Parts A and B pointwise, and use the measure criterion. Handle $m=0$ and
$m>n$ separately. The Lean assembly already expresses this dependency
structure, but inherits the unproved obligations' axioms.
\end{proof}

\section{Companion characterization}

\begin{lemma}[Lipschitz maps take null domain sets to Hausdorff-null sets]
\label{lem:lipschitz-null-image}
\lean{Preiss.hausdorffMeasure_image_eq_zero_of_lipschitz_null}
\uses{lem:plane-hausdorff}







Here \(A\subset\mathbb R^m\). Use the Hausdorff/Lebesgue normalization on the domain and Lipschitz distortion.
If \(f:A\to\mathbb R^n\) is Lipschitz and \(\mathcal L^m(A)=0\), then \(\mathcal H^m(f(A))=0\).
\par\emph{Role: Rectifiability characterization.}
\par Lean declarations: \texttt{Preiss.hausdorffMeasure\_image\_eq\_zero\_of\_lipschitz\_null}.
\end{lemma}
\begin{proof}[Planned proof obligations]

\par Source: De Lellis Chapters 4--5; see the rectifiability audit.
\end{proof}

\begin{lemma}[Affine graph area formula]
\label{lem:graph-affine-area}
\lean{Preiss.linearJacobian, Preiss.map_jacobian_hausdorff_affine_injection}
\uses{lem:plane-hausdorff}







The spaces are finite-dimensional Euclidean spaces and the affine linear part is injective. State an equality of measures, hence a formula on every Borel restriction.
For an injective affine graph parametrization, pushforward of its Jacobian-weighted source Haar measure equals Hausdorff measure on the affine graph, with the fixed normalization factor.
\par\emph{Role: Rectifiability characterization.}
\par Lean declarations: \texttt{Preiss.linearJacobian, Preiss.map\_jacobian\_hausdorff\_affine\_injection}.
\end{lemma}
\begin{proof}[Planned proof obligations]

\par Source: De Lellis Chapters 4--5; see the rectifiability audit.
\end{proof}

\begin{lemma}[Graphs are closed with finite bounded measure]
\label{lem:graph-borel-finite}
\lean{Preiss.isClosed_and_measure_closedBall_lt_top}
\uses{lem:graph-parametrization, lem:plane-hausdorff}







\(\Gamma=\{v+f(v):v\in V\}\), where \(\dim V=m\) and \(f:V\to V^\perp\) is Lipschitz.
Every full Lipschitz \(m\)-graph \(\Gamma\) is closed, and
\[
 \mathcal H^m(\Gamma\cap\overline B_R(x))<\infty
\]
for every \(x\in\mathbb R^n\) and finite radius \(R\).
\par\emph{Role: Rectifiability characterization.}
\par Lean declarations: \texttt{Preiss.isClosed\_and\_measure\_closedBall\_lt\_top}.
\end{lemma}
\begin{proof}[Planned proof obligations]

\par Source: De Lellis Chapters 4--5; see the rectifiability audit.
\end{proof}

\begin{lemma}[Full-rank differentiability locus maps into countably many graphs]
\label{lem:full-rank-graph-pieces}
\lean{Preiss.exists_countable_graph_cover_fullRank_locus}
\uses{lem:cone-graph, mathlib:derivative-partitions}








Use countable derivative-control pieces and a countable family of approximating full-rank matrices. For a nonempty rank-\(m\) locus, \(m\leq n\).
The image of the full-rank differentiability locus of a Lipschitz \(f:\mathbb R^m\to\mathbb R^n\) is covered by countably many full Lipschitz \(m\)-graphs.
\par\emph{Role: Rectifiability characterization.}
\par Lean declarations: \texttt{Preiss.exists\_countable\_graph\_cover\_fullRank\_locus}.
\end{lemma}
\begin{proof}[Planned proof obligations]

\par Source: De Lellis Chapters 4--5; see the rectifiability audit.
\end{proof}

\begin{lemma}[Rank-deficient derivative image is null]
\label{lem:rank-deficient-image-null}
\lean{Preiss.hausdorffMeasure_image_rank_lt_eq_zero}
\uses{lem:lipschitz-null-image}







Work with measurable pieces or suitable measurable envelopes. This is the rank-deficient-image lemma, not an invocation of an already available full area formula.
For a Lipschitz \(f:\mathbb R^m\to\mathbb R^n\), the image of its differentiability locus where \(\operatorname{rank}(Df)<m\) is \(\mathcal H^m\)-null.
\par\emph{Role: Rectifiability characterization.}
\par Lean declarations: \texttt{Preiss.hausdorffMeasure\_image\_rank\_lt\_eq\_zero}.
\end{lemma}
\begin{proof}[Planned proof obligations]

De Lellis Lemma4.11; do not claim the full area formula already exists.
\par Source: De Lellis Chapters 4--5; see the rectifiability audit.
\end{proof}

\begin{lemma}[Area formula for Lipschitz graph parametrizations]
\label{lem:graph-area}
\lean{Preiss.map_withDensity_jacobian_eq_graph_hausdorff}
\uses{lem:graph-affine-area, lem:lipschitz-null-image, mathlib:derivative-partitions}








Define the Jacobian on the differentiability locus and choose a measurable extension on its null complement. The map \(F\) is injective.
For a Lipschitz graph parametrization \(F(v)=v+f(v)\), Hausdorff measure on its graph is the pushforward of the measurable Jacobian-weighted domain measure, with the normalization established in the affine case.
\par\emph{Role: Rectifiability characterization.}
\par Lean declarations: \texttt{Preiss.map\_withDensity\_jacobian\_eq\_graph\_hausdorff}.
\end{lemma}
\begin{proof}[Planned proof obligations]

\par Source: De Lellis Chapters 4--5; see the rectifiability audit.
\end{proof}

\begin{lemma}[Fixed-dimensional blow-up transfer under density change]
\label{lem:normalized-density-transfer}
\lean{Preiss.tendsto_integral_blowup_sub_density_mul}
\uses{gmt:compact-test-perturbation, gmt:density-bounds, gmt:density-points, gmt:scaled-discarded-mass, mathlib:besicovitch-differentiation}







The measures are locally finite; \(f\geq0\) is locally \(\mu\)-integrable; \(x\) is a differentiation point and has finite upper \(m\)-density for \(\mu\).
If \(\lambda=f\mu\), then at good differentiation points \(x\), for every compactly supported continuous \(\varphi\),
\[
 \int\varphi\,dD^m_{x,r}\lambda-
 f(x)\int\varphi\,dD^m_{x,r}\mu\longrightarrow0.
\]
The exact multiplier \(f(x)\) is retained.
\par\emph{Role: Rectifiability characterization.}
\par Lean declarations: \texttt{Preiss.tendsto\_integral\_blowup\_sub\_density\_mul}.
\end{lemma}
\begin{proof}[Planned proof obligations]

Reuse GMT MatchingTangents norm\_integral\_cc\_le, weaklyConverges\_of\_tendsto\_integral\_sub, tendsto\_mul\_measure\_diff\_zero. Finite upper density addresses author erratum for Prop3.12.
\par Source: De Lellis Chapters 4--5; see the rectifiability audit.
\end{proof}

\begin{lemma}[Lipschitz-image and graph rectifiability agree]
\label{lem:graphs}
\lean{Preiss.countablyRectifiable_iff_lipschitzGraph_cover}
\uses{lem:domain-extension, lem:full-rank-graph-pieces, lem:graph-rectifiable, lem:lipschitz-null-image, lem:rank-deficient-image-null}







Natural dimensions \(m,n\) are arbitrary. For \(m>n\), use \(\mathcal H^m(\mathbb R^n)=0\) and the empty graph family. Measurability is not intrinsic to this set predicate.
A set \(S\subset\mathbb R^n\) is countably \(m\)-rectifiable if and only if there is a countable family of full Lipschitz \(m\)-graphs \(\Gamma_j\) such that
\[
 \mathcal H^m(S\setminus\bigcup_j\Gamma_j)=0.
\]
\par\emph{Role: Rectifiability characterization.}
\par Lean declarations: \texttt{Preiss.countablyRectifiable\_iff\_lipschitzGraph\_cover}.
\end{lemma}
\begin{proof}[Planned proof obligations]

\par Source: De Lellis Chapters 4--5; see the rectifiability audit.
\end{proof}

\begin{lemma}[Blow-ups of graph measure converge to its derivative plane]
\label{lem:graph-blowup}
\lean{Preiss.ae_graph_hausdorff_blowup_limit}
\uses{lem:graph-area, lem:graph-borel-finite, lem:plane-hausdorff}







Assume \(m>0\). Local finiteness follows from the graph measure API. Choose \(V_x\) before quantifying over positive sequences of radii.
For \(\mathcal H^m\)-almost every \(x\) of a full Lipschitz graph \(\Gamma\), one plane \(V_x\) satisfies
\[
 D^m_{x,r}(\mathcal H^m\!\restriction\Gamma)
 \rightharpoonup\mathcal H^m\!\restriction V_x
 \quad(r\downarrow0).
\]
\par\emph{Role: Rectifiability characterization.}
\par Lean declarations: \texttt{Preiss.ae\_graph\_hausdorff\_blowup\_limit}.
\end{lemma}
\begin{proof}[Planned proof obligations]

\par Source: De Lellis Chapters 4--5; see the rectifiability audit.
\end{proof}

\begin{lemma}[Density representation on a sigma-finite carrier]
\label{lem:representation}
\lean{Preiss.rectifiable_iff_density_representation}
\uses{lem:borel-carrier}







\(\mu\) is locally finite. Radon--Nikodym is applied on the carrier restriction, not on ambient \(\mathcal H^m\) when \(m<n\).
The measure \(\mu\) is \(m\)-rectifiable if and only if
\[
 \mu=(\mathcal H^m\!\restriction S)\mathbin{\mathrm{withDensity}} f
\]
for a Borel countably \(m\)-rectifiable set \(S\), a sigma-finite restriction \(\mathcal H^m\!\restriction S\), and a measurable \(f:\mathbb R^n\to[0,\infty]\).
\par\emph{Role: Rectifiability characterization.}
\par Lean declarations: \texttt{Preiss.rectifiable\_iff\_density\_representation}.
\end{lemma}
\begin{proof}[Planned proof obligations]

\par Source: De Lellis Chapters 4--5; see the rectifiability audit.
\end{proof}

\begin{theorem}[Rectifiable measures have a single normalized flat limit]
\label{thm:rectifiable-blowup}
\lean{Preiss.rectifiable_has_flat_blowup_limit}
\uses{lem:graph-blowup, lem:graphs, lem:normalized-density-transfer, lem:representation}







The measure \(\mu\) is locally finite and \(m\)-rectifiable, with \(0<m\leq n\).
At \(\mu\)-almost every \(x\) there exist one \(m\)-plane \(V_x\) and one \(0<c_x<\infty\) such that
\[
 D^m_{x,r}\mu\rightharpoonup c_x\mathcal H^m\!\restriction V_x
 \quad(r\downarrow0).
\]
The same \(V_x,c_x\) work for every positive sequence \(r_j\to0\).
\par\emph{Role: Rectifiability characterization.}
\par Lean declarations: \texttt{Preiss.rectifiable\_has\_flat\_blowup\_limit}.
\end{theorem}
\begin{proof}[Planned proof obligations]

\par Source: De Lellis Chapters 4--5; see the rectifiability audit.
\end{proof}

\begin{corollary}[Rectifiable density and common tangent plane]
\label{cor:rectifiable-tangent-plane}
\lean{Preiss.rectifiable_has_density_and_tangent_plane}
\uses{gmt:continuity-sets, lem:plane-hausdorff, lem:tangent-ray, thm:rectifiable-blowup}







The measure \(\mu\) is locally finite and \(m\)-rectifiable, with \(0<m\leq n\). Tangents use GMT's arbitrary positive scalar normalization.
At \(\mu\)-almost every \(x\), the \(m\)-density exists and is positive and finite, and there is one \(m\)-plane \(V_x\) such that
\[
 \operatorname{Tan}(\mu,x)
 =\{c\,\mathcal H^m\!\restriction V_x:0<c<\infty\}.
\]
\par\emph{Role: Rectifiability characterization.}
\par Lean declarations: \texttt{Preiss.rectifiable\_has\_density\_and\_tangent\_plane}.
\end{corollary}
\begin{proof}[Planned proof obligations]

\par Source: De Lellis Chapters 4--5; see the rectifiability audit.
\end{proof}

\begin{theorem}[Full tangent characterization]
\label{thm:rectifiable-iff-flat}
\lean{Preiss.rectifiable_iff_density_bounds_and_flat_tangents, Preiss.rectifiable_iff_density_and_tangent_plane}
\uses{cor:rectifiable-tangent-plane, thm:mm-measure}







The measure \(\mu\) is locally finite and \(0<m\leq n\). No density-limit existence or common-plane hypothesis is assumed on the right.
The measure \(\mu\) is \(m\)-rectifiable if and only if positive lower and finite upper \(m\)-density hold almost everywhere and every tangent is \(m\)-flat almost everywhere. On the right, the plane may depend on the individual tangent; rectifiability subsequently supplies a common plane almost everywhere.
Equivalently, at \(\mu\)-almost every \(x\), the positive finite
\(m\)-density exists and there is a single \(m\)-dimensional linear plane
\(V_x\) such that
\[
 \operatorname{Tan}(\mu,x)=
 \{c\,\mathcal H^m\!\restriction V_x:0<c<\infty\}.
\]
The second Lean declaration records this common-plane equivalence explicitly;
its assembly uses the same necessity and sufficiency results.
\par\emph{Role: Rectifiability characterization.}
\par Lean declarations: \texttt{Preiss.rectifiable\_iff\_density\_bounds\_and\_flat\_tangents}.
\end{theorem}
\begin{proof}[Planned proof obligations]

\par Source: De Lellis Chapters 4--5; see the rectifiability audit.
\end{proof}

\section{Explicit analytical data and their existence}

\begin{lemma}[Permutation symmetry of Gaussian moment forms]
\label{lem:moment-form-symmetry}
\lean{Preiss.gaussianMomentForm_permute}
\uses{def:moments}
Gaussian moment forms are invariant under every permutation of their arguments. The total Bochner definition is symmetric even when a nonintegrable kernel has integral zero.
\end{lemma}

\begin{definition}[A common polynomial growth bound]
\label{def:polynomial-growth-bound}
\lean{Preiss.HasPolynomialGrowthBound}
A common polynomial growth bound consists of finite real constants $C,p\geq0$ such that $\mu(\overline B(0,R))\leq C(1+R)^p$ for every $R>0$.
\end{definition}

\begin{definition}[An admissible flatness cutoff]
\label{def:flatness-cutoff}
\lean{Preiss.IsFlatnessCutoff}
A cutoff $\phi\in C_c(\mathbb R^n,\mathbb R)$ is admissible when $0\leq\phi\leq1$, it equals one on the closed unit ball, and its topological support lies in the open ball of radius two.
\end{definition}

\begin{lemma}[Existence of an admissible cutoff]
\label{lem:flatness-cutoff-exists}
\lean{Preiss.exists_flatnessCutoff}
\uses{def:flatness-cutoff}
An admissible flatness cutoff exists in every Euclidean dimension.
\end{lemma}

\begin{definition}[A cutoff flatness gap threshold]
\label{def:flatness-gap}
\lean{Preiss.IsFlatnessGap}
\uses{def:conical, def:flat, def:flatness, def:unit-uniform}
A positive number $\varepsilon$ is a cutoff gap threshold if every unit $m$-uniform conical measure with cutoff flatness at most $\varepsilon$ is flat.
\end{definition}

\begin{lemma}[Transfer the conical gap to cutoff flatness]
\label{lem:flatness-gap-cutoff}
\lean{Preiss.exists_flatnessGap}
\uses{thm:gap, lem:flatness-minimum}
For $0<m\leq n$ and an admissible cutoff, there exists a positive cutoff gap threshold.
\end{lemma}

\begin{definition}[Quantified crossing scale data]
\label{def:first-crossing-data}
\lean{Preiss.FirstCrossingData}
\uses{def:dimensional-blowup, def:flatness, gmt:vague}
Crossing data record positive scales $\sigma_i\leq r_i\to0$, vague limits of the two blow-up sequences, energy exactly $\varepsilon$ at every $\sigma_i$ and at its limit, and energy at most $\varepsilon$ on every interval $[\sigma_i,r_i]$. Divergence of $r_i/\sigma_i$ and flatness at infinity are separate conclusions.
\end{definition}

\begin{lemma}[Existence of the normal first-moment vector]
\label{lem:rigidity-normal-vector-exists}
\lean{Preiss.exists_firstMomentNormalVector}
\uses{lem:odd-coefficients-plane-zero, lem:first-moment-normal-from-vanishing}
For a unit uniform measure with a flat blow-down on $V$ and an explicit moment expansion, there exists $b\in V^\perp$ representing the first Taylor coefficient: $B_1^{(1)}(x)=2\langle b,x\rangle$.
\end{lemma}

\begin{lemma}[Existence of the rigidity bilinear forms]
\label{lem:rigidity-forms-exist}
\lean{Preiss.exists_rigidityOmega_and_adjustedForm}
\uses{lem:rigidity-forms-from-symmetry, def:moment-coefficients}
Given the moment expansion, flat blow-down plane, and normal first-moment vector, there exist the normal-valued symmetric bilinear form $\omega$ and adjusted real symmetric bilinear form satisfying their exact defining identities.
\end{lemma}

\begin{lemma}[Inner products and grading of polynomial features]
\label{lem:polynomial-feature-identities}
\lean{Preiss.polynomialFeatureProjection_apply, Preiss.polynomialFeature_inner, Preiss.polynomialFeatureProjection_isOrthogonalGrading, Preiss.polynomialFeatureProjection_inner, Preiss.polynomialFeature_weighted_inner}
\uses{def:polynomial-feature, def:orthogonal-grading, lem:orthonormal-span-projection}
The coordinate polynomial feature map recovers the weighted powers of $\langle x,y\rangle$ through inner products. Its homogeneous coordinate blocks form an orthogonal grading.
\end{lemma}

\section{Uniform constants and normalization}

\begin{definition}[Uniform measures with a specified ball constant]
\label{def:uniform-constant}
\lean{Preiss.IsUniformWithConstant}
An uncentered uniform measure with constant $q$ is regular, has nonempty support, satisfies $q>0$, and gives every closed support-centered ball of radius $r>0$ mass $qr^m$. Nonempty support excludes the zero measure.
\end{definition}

\begin{lemma}[Centered unit normalization]
\label{lem:uniform-constant-one}
\lean{Preiss.isUnitUniform_iff_constant_one}
\uses{def:uniform-constant, def:unit-uniform}
A measure is unit uniform exactly when it is uniform with constant one and its support contains the origin.
\end{lemma}

\begin{lemma}[Scaling uniform measures and their ball constants]
\label{lem:uniform-scaling}
\lean{Preiss.uniformWithConstant_dimensionalBlowup, Preiss.isUnitUniform_dimensionalBlowup, Preiss.uniformWithConstant_smul}
\uses{def:uniform-constant, def:dimensional-blowup, lem:uniform-constant-one, lem:blowup-composition, lem:dimensional-blowup-support}
Dimensional blow-ups preserve a specified uniform ball constant. Centering at a support point preserves unit uniformity. Multiplication by a positive finite scalar multiplies the ball constant by the same scalar.
\end{lemma}

\section{Generic density limit values}
\begin{lemma}[Values of existing dimensional densities]
\label{lem:density-limit-values}
\lean{Preiss.dimensional_density_eq_of_tendsto, Preiss.lower_eq_upper_of_hasDensity, Preiss.density_eq_lower_of_hasDensity, Preiss.density_eq_upper_of_hasDensity}
\uses{def:density-api}
For an outer measure on a pseudo extended metric space and any real dimension parameter, convergence of the dimensional density ratio identifies its density value with the limit. Whenever the density exists, its value equals both the lower and upper densities.
\end{lemma}

\begin{lemma}[Continuity, homogeneity, and norm bounds for moment kernels]
\label{lem:moment-kernel-calculus}
\lean{Preiss.continuous_momentKernel, Preiss.momentKernel_smul, Preiss.norm_momentKernel_le, Preiss.norm_momentKernel}
\uses{def:moments}
The degree-$k$ multilinear moment kernel depends continuously on its vector argument,
is homogeneous of degree $k$ under real scalar multiplication, and has operator norm
at most the $k$th power of the vector norm. These assertions include degree zero.
\end{lemma}
\begin{proof}
Use the continuous multilinear product of the inner-product functionals, together
with their operator-norm bound.
\end{proof}

\section{Reusable upper Hausdorff comparison}

\begin{lemma}[Hausdorff comparison from diameter bounds]
\label{lem:diameter-hausdorff-bound}
\lean{Preiss.measure_le_smul_hausdorff_of_diameter_bound}
\uses{mathlib:hausdorff-mass-bound}
On a metric space, suppose $0<c<\infty$, $\varepsilon>0$, and
$\mu(T)\le c\,\operatorname{diam}(T)^m$ whenever the extended diameter
is at most $\varepsilon$. Then $\mu\le c\mathcal H^m$.
The dimension $m$ is natural, including zero.
\end{lemma}
\begin{proof}
Apply the mass-distribution bound to $c^{-1}\mu$, then multiply by $c$.
\end{proof}

\begin{lemma}[Hausdorff comparison from upper ball bounds]
\label{lem:ball-upper-hausdorff-bound}
\lean{Preiss.measure_restrict_le_hausdorff_of_uniform_ball_upper_bound}
\uses{lem:diameter-hausdorff-bound}
Let $S$ be Borel in a metric space, and let $b,R>0$. If
$\mu(B(x,r))\le b r^m$ for every $x\in S$ and $0<r<R$, then
$\mu\restriction S\le b\mathcal H^m\restriction S$.
Neither properness nor local finiteness is required.
\end{lemma}
\begin{proof}
For a set of diameter below $R/2$ meeting $S$, choose a point in that
intersection and enclose the set in balls whose radii decrease to its diameter.
Continuity of the power bound gives the diameter estimate, including diameter
zero. Restrict the resulting measure inequality to $S$.
\end{proof}

\begin{lemma}[Closed upper-density strata]
\label{lem:upper-density-localization}
\lean{Preiss.exists_closedBall_upper_bound_of_upperDensity_lt_top, Preiss.exists_goodBallSet_mem_of_closedBall_upper_bound, Preiss.exists_goodBallSet_mem_of_upperDensity_lt_top}
\uses{def:density-api, gmt:goodballs}
A point of finite upper $m$-density has a finite positive upper power bound
on sufficiently small centred closed balls. Such a bound places the point
in one of the upstream closed good-ball sets with natural constant and radius
indices. No positive lower density is assumed.
\end{lemma}
\begin{proof}
First bound the density ratio strictly above its limsup. A ball of radius $r$
containing $x$ lies in the ball centred at $x$ of radius $2r$. Enlarge the
power constant to a natural number and decrease the radius to $1/(k+1)$.
These countably many closed strata cover all points of finite upper density.
\end{proof}

\section{Annular estimates for Gaussian integrals}

\begin{lemma}[Integrability from summable annulus bounds]
\label{lem:annulus-integrability}
\lean{Preiss.integrable_of_summable_annulus_bounds}
Let $E$ be a proper normed additive group with its Borel structure, and let $\mu$ be finite on every positive-radius closed ball centred at zero. Let $f:E\to F$ be continuous, where $F$ is a normed additive group. If $M_j\ge0$, $\|f(x)\|\le M_j$ on $j\le\|x\|\le j+1$, and $\sum_j M_j\mu(\overline B(0,j+1))$ is summable as a real series, then $f$ is Bochner integrable.
\end{lemma}
\begin{proof}
Integrate on compact unit annuli, use the finite ball masses, and apply summability of the integral bounds to their countable cover.
\end{proof}

\begin{lemma}[Integral bound from a countable measurable cover]
\label{lem:nonnegative-integral-cover-bound}
\lean{Preiss.integral_iUnion_le_tsum_bounds}
For any measure on a measurable space, let $f$ be a nonnegative integrable real function and let $S_j$ be measurable sets. If a summable real sequence $B_j$ bounds $\int_{S_j}f$, then $\int_{\bigcup_j S_j}f\le\sum_j B_j$. The sets may overlap.
\end{lemma}
\begin{proof}
Disjointify the measurable cover and use nonnegativity before summing.
\end{proof}

\begin{lemma}[A summable common Gaussian annulus majorant]
\label{lem:gaussian-annulus-majorant}
\lean{Preiss.summable_gaussian_annulus_majorant}
For $C\ge0$, $p\in\R$, $d\in\mathbb N$, and $s>0$, the sequence $B_j=(j+1)^d e^{-s j^2}C(j+2)^p$ is summable. Bound the polynomial factor by a natural power and use $j^2\ge j$ to compare with a polynomial times the geometric sequence $(e^{-s})^j$.
\end{lemma}

\begin{lemma}[Gaussian integrals on unit annuli]
\label{lem:gaussian-annulus-integral}
\lean{Preiss.gaussian_annulus_integral_bound}
Suppose $C\ge0$ and $\mu(\overline B(0,R))\le C(1+R)^p$ for all $R>0$. For $j,d\in\mathbb N$ and $s>0$, the integral of $\|z\|^d e^{-s\|z\|^2}$ over $j\le\|z\|\le j+1$ is at most $B_j=(j+1)^d e^{-s j^2}C(j+2)^p$. No regularity hypothesis is needed.
\end{lemma}

\begin{lemma}[A common series controls every Gaussian tail]
\label{lem:gaussian-tail-series-bound}
\lean{Preiss.gaussian_tail_le_annulus_sum}
\uses{lem:gaussian-polynomial-integrable, lem:nonnegative-integral-cover-bound, lem:gaussian-annulus-majorant, lem:gaussian-annulus-integral}
Under the polynomial ball bound above, for $N\in\mathbb N$ and $R\ge N$, the integral of $\|z\|^d e^{-s\|z\|^2}$ outside $\overline B(0,R)$ is at most $\sum_{j\ge0}B_{j+N}$. The bound is independent of the measure.
\end{lemma}

\begin{lemma}[Integrability of all polynomial Gaussian weights]
\label{lem:polynomial-gaussian-integrability}
\lean{Preiss.integrable_norm_pow_mul_polynomial_mul_gaussian}
\uses{lem:gaussian-polynomial-integrable}

For a measure mu on Euclidean n with mu(closedBall(0,R))<=ofReal(C(1+R)^p) at every R>0, where C>=0 and p is real, every multivariate real polynomial P, natural d, and s>0 satisfy Bochner integrability of norm(z)^d P(z) exp(-s norm(z)^2). The proof inducts on P using closure under constants, addition, and multiplication by a coordinate; norm(z_i)<=norm(z) supplies the inductive domination.
\end{lemma}

\begin{lemma}[Weighted ball scaling for conical measures]
\label{lem:homogeneous-ball-scaling}
\lean{Preiss.integral_homogeneous_closedBall, Preiss.integrableOn_homogeneous_closedBall}
\uses{def:conical, def:dimensional-blowup}

For a conical measure $\mu$ and a measurable real-valued function $h$ homogeneous of arbitrary real degree $a$, for every $r>0$, $\int_{\overline B(0,r)}h\,d\mu=r^{m+a}\int_{\overline B(0,1)}h\,d\mu$. No positivity or integrability assumption is needed for this scaling identity; both sides use the Bochner nonintegrable-value convention. Under the same scaling hypotheses, integrability of h on the unit closed ball implies integrability on every positive-radius closed ball.
\end{lemma}

\begin{lemma}[Small covers from a lower ball bound]
\label{lem:lower-ball-small-cover}
\lean{Preiss.exists_small_cover_of_ball_lower_bound}
\uses{mathlib:vitali-enlargement}

Let $X$ be a separable metric Borel space, $A\subseteq U$ with $U$ open, and $a,R>0$. Assume $a r^m\leq\mu(B(x,r))$ for all $x\in A$ and $0<r<R$. For every $\varepsilon>0$, there is a countable cover of $A$ by fivefold enlargements of disjoint balls in $U$, each enlarged ball having diameter at most $\varepsilon$, and total diameter-to-the-$m$ cost at most $(10^m/a)\mu(U)$.
\end{lemma}

\begin{lemma}[Lower ball bounds control Hausdorff measure in open supersets]
\label{lem:lower-ball-open-comparison}
\lean{Preiss.hausdorff_le_open_mass_of_ball_lower_bound}
\uses{lem:lower-ball-small-cover}

Under the lower ball bound above, $\mathcal H^m(A)\leq (10^m/a)\mu(U)$ for every open superset $U$ of $A$. No measurability of $A$ or regularity of $\mu$ is needed.
\end{lemma}

\begin{lemma}[Hausdorff comparison by outer regularity]
\label{lem:lower-ball-hausdorff-bound}
\lean{Preiss.hausdorff_le_measure_of_uniform_ball_lower_bound}
\uses{lem:lower-ball-open-comparison}

If $\mu$ is outer regular and the same uniform lower ball bound holds at every point of an arbitrary set $A$, then $\mathcal H^m(A)\leq (10^m/a)\mu(A)$.
\end{lemma}

\begin{lemma}[Restricted lower ball comparison]
\label{lem:lower-ball-restrict-comparison}
\lean{Preiss.hausdorff_restrict_le_measure_of_uniform_ball_lower_bound}
\uses{lem:lower-ball-hausdorff-bound}

For a Borel stratum $S$ and an outer regular measure satisfying the uniform lower ball bound at every point of $S$, $\mathcal H^m\!\restriction S\leq (10^m/a)\mu\!\restriction S$.
\end{lemma}

\begin{lemma}[Support under homeomorphisms and nonzero scalar multiplication]
\label{lem:measure-support-transport}
\lean{Preiss.measure_support_map_homeomorph, Preiss.measure_support_smul}

For a homeomorphism $e:X\to Y$ of Borel topological spaces and any measure $\mu$ on $X$, $\operatorname{spt}(e_\#\mu)=e(\operatorname{spt}\mu)$. For a measure on any topological measurable space and any $c\in[0,\infty]$ with $c\ne0$, $\operatorname{spt}(c\mu)=\operatorname{spt}\mu$. Neither assertion requires regularity or local finiteness.
\end{lemma}

\begin{lemma}[Support and regularity of dimensional blow-ups]
\label{lem:dimensional-blowup-support}
\lean{Preiss.dimensionalBlowup_support, Preiss.regular_dimensionalBlowup}
\uses{def:dimensional-blowup, lem:measure-support-transport}

For every $m,n\in\mathbb N$, measure $\mu$ on $\mathbb R^n$, point $a$, and radius $r>0$, the dimensional blow-up $\mu_{a,r}=r^{-m}(T_{a,r})_\#\mu$ has support $T_{a,r}(\operatorname{spt}\mu)$, where $T_{a,r}(y)=(y-a)/r$. If $\mu$ is regular, then $\mu_{a,r}$ is regular.
\end{lemma}

\begin{lemma}[Rational good sets from positive finite density bounds]
\label{lem:density-bounds-good-set}
\lean{Preiss.exists_goodSet_mem_of_density_bounds}
\uses{def:positive-density, gmt:goodsets}

If $0<\Theta_*^m(\mu,x)$ and $\Theta^{m,*}(\mu,x)<\infty$, then there are positive rational numbers $p,q$ and $k\in\mathbb N$ such that, for every $0<r<1/(k+1)$, $p r^m\leq\mu(\overline B(x,r))\leq q r^m$. Equivalently, $x\in\mathrm{goodSet}(m,p,q,k,\mu)$ using the existing FormalizingGMT definition.
\end{lemma}

\begin{lemma}[Exact blow-up invariance of flat plane measures]
\label{lem:flat-plane-blowup}
\lean{Preiss.dimensionalBlowup_flatOn}
\uses{def:dimensional-blowup, lem:plane-hausdorff, lem:blowup-composition}

For an $m$-plane $V$, every scalar $c\in[0,\infty]$, every $x\in V$, and $r>0$, $D^m_{x,r}(c\,\mathcal H^m\!\restriction V)=c\,\mathcal H^m\!\restriction V$. This includes $m=0$ and does not require $c$ positive or finite.
\end{lemma}

\begin{lemma}[Nonnegativity of distance energy and flatness]
\label{lem:flatness-basic-nonneg}
\lean{Preiss.flatnessEnergy_nonneg, Preiss.flatness_nonneg_of_nonneg}
\uses{def:flatness}

For every nonnegative compactly supported continuous cutoff and every measure, the squared-distance energy of every plane is nonnegative, and so is its infimum defining flatness. No regularity or minimizing-plane theorem is needed; when the plane type is empty, the real infimum convention gives zero.
\end{lemma}

\begin{lemma}[Vanishing distance energy on a supporting plane]
\label{lem:flatness-basic-zero}
\lean{Preiss.flatnessEnergy_restrict_plane, Preiss.flatness_eq_zero_of_isFlat_of_nonneg}
\uses{def:flatness, def:flat, lem:flatness-basic-nonneg}

The squared-distance energy on $V$ of $c\,\mathcal H^m\!\restriction V$ is zero for any scalar $c\in[0,\infty]$. Consequently every flat measure has zero flatness for every nonnegative cutoff, without regularity or existence of a global minimizer.
\end{lemma}

\begin{lemma}[Gaussian integral from a centered power ball law]
\label{lem:gaussian-ball-power}
\lean{Preiss.lintegral_gaussian_of_closedBall_power, Preiss.integral_gaussian_of_closedBall_power}
\uses{lem:exponential-tail-integral}

For an s-finite positive measure $\rho$ on a normed additive commutative group with its Borel measurable structure, $C\geq0$, and $d>0$, if $\rho(\overline B(0,r))=Cr^d$ for every $r>0$, then $e^{-\|z\|^2}$ is integrable and its integral is $C\Gamma(d/2+1)$. The centered ball law suffices.
\end{lemma}

\begin{lemma}[The nonnegative exponential tail integral]
\label{lem:exponential-tail-integral}
\lean{Preiss.lintegral_exp_neg_Ici}

For every real a, the Lebesgue lintegral of ofReal(exp(-t)) over Ici(a) equals ofReal(exp(-a)).
\end{lemma}

\begin{lemma}[Lower density from a lower power bound]
\label{lem:lower-density-from-power}
\lean{Preiss.lower_density_ge_of_ball_power_lower_bound}
\uses{def:density-api}

Let $\mu$ be an outer measure on a pseudometric space, $m\in\mathbb N$, $a\geq0$, and $R>0$. If $\mu(B(x,r))\geq a r^m$ for every $0<r<R$, then $\Theta_*^m(\mu,x)\geq a/2^m$. The factor $2^m$ comes from the diameter-normalized definition of density.
\end{lemma}

\begin{lemma}[All exponential moments make power moments determining]
\label{lem:complex-mgf-moment-determination}
\lean{Preiss.complexMGF_eq_of_all_moments_eq}

Let $X$ and $Y$ be real functions on arbitrary measure spaces $(\Omega,\mu)$ and $(\Omega^\prime,\nu)$. Suppose $\exp(tX)$ and $\exp(tY)$ are integrable for every real $t$. If $\int X^k\,d\mu=\int Y^k\,d\nu$ for every natural $k$, then their complex moment generating functions agree at every complex argument.
\end{lemma}

\begin{lemma}[Gaussian weighting gives a finite measure]
\label{lem:gaussian-weight-finite}
\lean{Preiss.isFiniteMeasure_withDensity_gaussian}
\uses{def:polynomial-growth-bound, lem:gaussian-polynomial-integrable}

If $\mu$ on $\mathbb R^n$ obeys $\operatorname{HasPolynomialGrowthBound}(\mu,C,p)$, then the measure $e^{-\|z\|^2}\mu$ is finite.
\end{lemma}

\begin{lemma}[Every projection has all exponential moments after Gaussian weighting]
\label{lem:gaussian-weight-exponential}
\lean{Preiss.integrable_exp_inner_mul_gaussian, Preiss.integrable_exp_inner_withDensity_gaussian, Preiss.integrable_exp_inner_mul_gaussian_parameter, Preiss.integrable_exp_abs_inner_mul_gaussian_parameter}
\uses{def:polynomial-growth-bound, lem:gaussian-polynomial-integrable}

If $\operatorname{HasPolynomialGrowthBound}(\mu,C,p)$, then for every $w\in\mathbb R^n$ and $t\in\mathbb R$, the function $e^{-\|z\|^2}e^{t\langle z,w\rangle}$ is $\mu$-integrable. Equivalently, $e^{t\langle z,w\rangle}$ is integrable for the Gaussian weighted measure.
\end{lemma}

\begin{lemma}[Existence of a plane in every admissible dimension]
\label{lem:plane-exists}
\lean{Preiss.nonempty_plane_of_le}
\uses{def:plane}

If $m\leq n$, the type of $m$-dimensional linear subspaces of $\mathbb R^n$ is nonempty. Choose $m$ linearly independent vectors and take their span.
\end{lemma}

\begin{lemma}[Continuity of projection energy]
\label{lem:projection-energy-continuity}
\lean{Preiss.flatnessEnergy_eq_integral_projection, Preiss.continuous_integral_projectionEnergy}
\uses{def:flatness}

The distance energy of a plane is the integral of $\varphi(z)|z-P_Vz|^2$. For any compactly supported continuous $\varphi$ and any measure finite on compact sets, the same integral is a continuous function of an arbitrary continuous linear operator $P$.
\end{lemma}

\begin{lemma}[Minimizing flatness under finiteness on compact sets]
\label{lem:flatness-minimizer-finite-compacts}
\lean{Preiss.flatness_exists_minimizer_of_finite_on_compacts}
\uses{def:flatness, lem:plane-exists, lem:projection-energy-continuity, lem:grassmannian}

For $m\leq n$, every compactly supported continuous real cutoff and every measure finite on compact sets admit a plane attaining the infimum of distance energy. The cutoff need not be nonnegative. Minimize the continuous operator energy over the compact set of rank-$m$ orthogonal projections.
\end{lemma}

\begin{lemma}[Integrability of compactly weighted distance energy]
\label{lem:flatness-energy-integrable}
\lean{Preiss.integrable_flatnessEnergy_integrand}
\uses{def:flatness}

For every compactly supported continuous real cutoff, every plane, and every measure finite on compact sets, the cutoff times squared distance to the plane is integrable.
\end{lemma}

\begin{lemma}[Continuous finite ball profiles have null spheres]
\label{lem:continuous-ball-profile-boundary}
\lean{Preiss.ball_ae_eq_closedBall_of_continuous_profile, Preiss.measure_sphere_eq_zero_of_continuous_profile, Preiss.integral_ball_eq_closedBall_of_continuous_profile}

Let $\mu$ be a Borel measure on a pseudometric space, $x$ a point, and $r>0$. Suppose a real function $f$ is continuous at $r$ and $\mu(\overline B(x,s))=\operatorname{ofReal}(f(s))$ for every $s>0$. Then $B(x,r)=\overline B(x,r)$ almost everywhere, the sphere of radius $r$ has measure zero, and every normed real vector-valued Bochner integrand has equal integrals on the open and closed balls.
\end{lemma}

\begin{lemma}[Open and closed uniform balls agree almost everywhere]
\label{lem:uniform-ball-boundary}
\lean{Preiss.unitUniform_ball_ae_eq_closedBall}
\uses{def:unit-uniform, lem:continuous-ball-profile-boundary}

For every unit $m$-uniform measure $\mu$ on $\mathbb R^n$, support center $x$, and positive radius $r$, $B(x,r)=\overline B(x,r)$ almost everywhere for $\mu$. This includes $m=0$.
\end{lemma}

\begin{lemma}[Derivative of a product with the identity at zero]
\label{lem:iterated-derivative-id-product}
\lean{Preiss.iteratedDeriv_id_mul_zero}

For every positive natural $n$ and every real function $f$ that is $C^n$ at zero, $D^n(t\mapsto t f(t))(0)=nD^{n-1}f(0)$.
\end{lemma}

\begin{lemma}[All even derivatives of a quadratic exponential]
\label{lem:quadratic-exponential-even-derivatives}
\lean{Preiss.iteratedDeriv_exp_quadratic_even}
\uses{lem:iterated-derivative-id-product}

For every real $a$ and natural $j$, $D^{2j}(t\mapsto e^{at^2/2})(0)=\frac{(2j)!}{2^j j!}a^j$. No sign restriction on $a$ is required.
\end{lemma}

\begin{lemma}[All even moments of a centered real Gaussian]
\label{lem:real-gaussian-even-moments}
\lean{Preiss.integral_gaussianReal_zero_pow_even}
\uses{lem:quadratic-exponential-even-derivatives}

For every nonnegative real variance $v$, including zero, and every natural $j$, $\int x^{2j}\,d\mathcal N(0,v)=\frac{(2j)!}{2^j j!}v^j$.
\end{lemma}

\begin{lemma}[Integrability of finite products of Gaussian projections]
\label{lem:gaussian-projection-integrability}
\lean{Preiss.integrable_gaussian_inner_pow, Preiss.integrable_gaussian_inner_prod, Preiss.integrable_gaussian_four_inner}

For a standard Gaussian on any finite dimensional real inner product space $E$, every natural power of a real linear projection is integrable. For every finite index type $I$ and vectors $(u_i)_{i\in I}$, the function $x\mapsto\prod_i\langle u_i,x\rangle$ is integrable, including the empty product. A four-projection companion is also provided.
\end{lemma}

\begin{lemma}[Detecting convergence by further subsequences]
\label{lem:subsequence-limit-criterion}
\lean{Preiss.tendsto_of_subsequence_subsequence}

A sequence in a pseudometric space converges to $x$ if every strictly increasing subsequence has a further strictly increasing subsequence converging to $x$.
\end{lemma}

\begin{lemma}[Moving projection energies under vague convergence]
\label{lem:vague-projection-energy}
\lean{Preiss.vague_integral_projectionEnergy_moving}
\uses{lem:vague-moving-tests}

Let $\mu_i\rightharpoonup\nu$ be regular Euclidean measures, $P_i\to P$ in operator norm, and $\varphi$ a compactly supported continuous real function. Then $\int\varphi(z)|z-P_i z|^2\,d\mu_i\to\int\varphi(z)|z-Pz|^2\,d\nu$. Operators need not be projections and the cutoff need not be nonnegative.
\end{lemma}

\begin{lemma}[Vague continuity of minimum energy for general cutoffs]
\label{lem:flatness-vague-general}
\lean{Preiss.flatness_tendsto_of_vague}
\uses{def:flatness, lem:subsequence-limit-criterion, lem:vague-projection-energy, lem:flatness-minimizer-finite-compacts, lem:grassmannian}

For $m\leq n$, every compactly supported continuous real cutoff has vaguely continuous minimum squared-distance energy on regular measures. Extract convergent minimizing projections along every subsequence, use the moving-integral lemma, and compare with a fixed minimizing plane for the limit measure. No nonnegativity or special cutoff normalization is needed.
\end{lemma}

\begin{lemma}[Measure bounds from eventual membership in moving sets]
\label{lem:moving-set-measure-bounds}
\lean{Preiss.measure_le_of_pointwise_eventually_mem, Preiss.le_measure_of_pointwise_eventually_notMem}

Let $S_i,T\subseteq X$ and $c\in[0,\infty]$. If eventually $\mu(S_i)\le c$ and every $x\in T$ eventually belongs to $S_i$, then $\mu(T)\le c$. If $\mu$ is finite, each $S_i$ is $\mu$-null-measurable, eventually $c\le\mu(S_i)$, and every $x\notin T$ is eventually outside $S_i$, then $c\le\mu(T)$.
\end{lemma}

\begin{lemma}[Distance to a subspace through orthogonal projection]
\label{lem:projection-distance}
\lean{Preiss.infDist_submodule_eq_norm_sub_starProjection}

In a real or complex inner-product space, let $V$ be a linear subspace admitting an orthogonal projection $P_V$. For every $x$, $\operatorname{dist}(x,V)=\lVert x-P_Vx\rVert$.
\end{lemma}

\begin{lemma}[Convergent subsequences of fixed-rank plane projections]
\label{lem:plane-projection-subsequence}
\lean{Preiss.exists_tendsto_planeProjection_subseq}
\uses{def:plane, lem:grassmannian}

For every sequence of $m$-planes $V_i\subseteq\mathbb R^n$, there exist an $m$-plane $V$ and a strictly increasing $\phi:\mathbb N\to\mathbb N$ such that $P_{V_{\phi(i)}}\to P_V$ in operator norm.
\end{lemma}

\begin{lemma}[Compact approximation inside an increasing measurable exhaustion]
\label{lem:compact-measurable-exhaustion}
\lean{Preiss.exists_compact_subset_of_ae_mem_iUnion, Preiss.exists_pos_compact_subset_of_ae_mem_iUnion}

Let $\mu$ be a finite regular Borel measure on a Hausdorff space, let $E$ and $S_i$ be measurable, and assume $S_i\subseteq S_{i+1}$ and $\mu(E\setminus\bigcup_i S_i)=0$. For every $\varepsilon>0$ there exist $N$ and a compact set $C\subseteq E\cap S_N$ such that $\mu(E\setminus C)<\varepsilon$.
\end{lemma}

\begin{lemma}[Projection onto a finite orthonormal span]
\label{lem:orthonormal-span-projection}
\lean{Preiss.starProjection_span_orthonormal, Preiss.starProjection_span_orthonormal_apply}

For every finite orthonormal family $(u_i)$ in a finite dimensional real inner product space, the orthogonal projection onto its span is the sum of the rank-one maps $x\mapsto\langle u_i,x\rangle u_i$. The family may be empty.
\end{lemma}

\begin{definition}[Orthogonal grading and its weighted complement]
\label{def:orthogonal-grading}
\uses{}
\lean{Preiss.IsOrthogonalGrading, Preiss.weightedGrading, Preiss.weightedComplement}

An orthogonal grading consists of finitely many self-adjoint idempotents $P_i$ with $P_iP_j=0$ for distinct indices and $\sum_iP_i=\operatorname{id}$. For every real $s$, define $A(s)=\sum_i s^{i+1}P_i$ and $F_s=A(s)^{-1}(V^\perp)$ for a subspace $V$. These are data and total definitions, with no analytic-extension assertion.
\end{definition}

\begin{lemma}[A metric from countably many separating real coordinates]
\label{lem:countable-coordinate-metric}
\lean{Preiss.countableCoordinateMap, Preiss.countableCoordinateMap_injective, Preiss.countableCoordinateMetric, Preiss.countableCoordinateMetric_dist, Preiss.countableCoordinateMetric_tendsto_iff}

For an arbitrary type $A$, let $f_j:A\to\mathbb R$, $j\in\mathbb N$, separate points. The explicit shifted weighted coordinate map induces a metric $d(a,b)=\sum_{j\geq0}2^{-(j+1)}\min(1,|f_j(a)-f_j(b)|)$. For every type $B$, filter $\mathcal F$ on $B$, function $g:B\to A$, and point $a\in A$, convergence of $g$ to $a$ in this metric is equivalent to convergence of every $f_j\circ g$ to $f_j(a)$.
\end{lemma}

\begin{lemma}[Norm powers as explicit feature functionals]
\label{lem:feature-norm-functionals}
\lean{Preiss.norm_pow_eq_sum_coordinate_monomials, Preiss.exists_normPolynomial_feature_functionals}
\uses{def:polynomial-feature}

For every natural $d$, $\|x\|^{2d}=\sum_{f:\operatorname{Fin}d\to\operatorname{Fin}n}(\prod_jx_{f(j)})^2$. Consequently, for all natural $n,q$, there exist continuous linear functionals $L_j$ on the degree-$2q$ feature space, indexed by $j:\operatorname{Fin}q$, with $L_j(\Phi_{2q}(x))=\|x\|^{2(j+1)}$ for every $x$.
\end{lemma}

\begin{lemma}[General tangents, scalar multiplication and dilation]
\label{lem:dimensional-tangent-operations}
\lean{Preiss.IsDimensionalTangent.isTangentMeasure, Preiss.IsDimensionalTangent.smul, Preiss.IsDimensionalTangent.dilate}
\uses{def:dimensional-tangent, lem:blowup-composition, lem:dimensional-blowup-support, lem:vague-moving-blowups, def:tangent-api}

Every dimensional tangent is a general tangent measure. If $c$ is positive and finite, $\nu\in\operatorname{Tan}_m(\mu,x)$ implies $c\nu\in\operatorname{Tan}_m(c\mu,x)$. For locally finite $\mu$ and $s>0$, origin dilation also preserves $\operatorname{Tan}_m(\mu,x)$. The first two assertions need no local-finiteness hypothesis.
\end{lemma}

\begin{lemma}[Vanishing upper density at a fixed positive ball scale]
\label{lem:vanishing-density-scaled-balls}
\lean{Preiss.tendsto_scaled_ball_mass_div_pow_zero_of_upperDensity_eq_zero}
\uses{def:density-api}

Let $\nu$ be an outer measure on a pseudometric space and $m\in\mathbb N$. If $\Theta^{*m}(\nu,x)=0$, then for every fixed $c>0$, $\nu(B_{cr}(x))/r^m\to0$ as $r\downarrow0$.
\end{lemma}

\begin{lemma}[Vanishing density of a deleted restriction gives negligible deletion]
\label{lem:vanishing-density-negligible-deletion}
\lean{Preiss.negligibleDeletionAt_of_upperDensity_restrict_compl_eq_zero}
\uses{def:density-api, def:weak-approximation, lem:vanishing-density-scaled-balls}

For a measure $\mu$ on $\mathbb R^n$ and a measurable set $F$, if $\Theta^{*m}(\mu|_{F^c},x)=0$, then $\mu(B_{2r}(x)\setminus F)/r^m\to0$ as $r\downarrow0$.
\end{lemma}

\begin{definition}[Dense compact tests with common supports and cutoffs]
\label{def:dense-compact-tests}
\lean{Preiss.DenseCompactTestFamily}

For a topological space $X$, DenseCompactTestFamily consists of a sequence $f_j\in C_c(X,\mathbb R)$ containing nonnegative cutoffs majorizing $1$ on every compact set. For every $f\in C_c(X,\mathbb R)$ there is one compact $K$ containing its support such that, for every $\varepsilon>0$, some $f_j$ is supported in $K$ and satisfies $|f_j-f|<\varepsilon$ everywhere. The compact set is fixed before choosing the accuracy.
\end{definition}

\begin{lemma}[Dense continuous maps approximate uniformly on compact sets]
\label{lem:dense-continuous-approximation}
\lean{Preiss.dense_continuousMap_uniform_approximation}

For an arbitrary topological space $X$, a sequence with dense range in $C(X,\mathbb R)$ approximates each continuous function to every positive accuracy uniformly on every fixed compact set.
\end{lemma}

\begin{lemma}[Existence of a dense compact-test family]
\label{lem:dense-compact-tests-exist}
\lean{Preiss.exists_denseCompactTestFamily}
\uses{def:dense-compact-tests, lem:dense-continuous-approximation}

Every locally compact, second-countable Hausdorff space admits DenseCompactTestFamily data.
\end{lemma}

\begin{lemma}[Supplied compact tests determine vague convergence and measures]
\label{lem:vague-determining-tests}
\lean{Preiss.weaklyConverges_iff_denseCompactTests, Preiss.measure_eq_of_denseCompactTests}
\uses{def:dense-compact-tests, gmt:vague, mathlib:regular-measure-ext}

For a supplied DenseCompactTestFamily on $\mathbb R^n$ and regular measures $\mu_i,\nu$, vague convergence $\mu_i\rightharpoonup\nu$ is equivalent to convergence of the integrals of every enumerated test. Two regular measures with equal integrals on every enumerated test are equal.
\end{lemma}

\begin{definition}[Constructing a vague metric from supplied compact tests]
\label{def:vague-metrization-from-tests}
\lean{Preiss.vagueMetrizationOfDenseCompactTests}
\uses{def:dense-compact-tests, def:vague-metrization, lem:vague-determining-tests, lem:countable-coordinate-metric}

A supplied DenseCompactTestFamily on $\mathbb R^n$ produces explicit VagueMetrization data by applying the countable-coordinate metric to test integrals of regular measures. Point separation and the exact convergence specification follow from the supplied-test determination theorem.
\end{definition}

\begin{lemma}[An operator-valued integral preserves a common kernel vector]
\label{lem:integral-operator-annihilation}
\lean{Preiss.integral_apply_eq_zero_of_ae}

Let $A:\Omega\to\mathcal L(H,F)$ for arbitrary real normed spaces $H,F$ and a measure $\mu$. If $A(z)u=0$ for almost every $z$, then $(\int A\,d\mu)u=0$. No integrability, completeness, or measurability assumption on $A$ is required beyond the almost-everywhere assertion.
\end{lemma}

\begin{definition}[The support feature span and weighted feature moment]
\label{def:feature-moment}
\lean{Preiss.supportFeatureSpan, Preiss.featureMoment}
\uses{def:polynomial-feature, def:orthogonal-grading, def:moments}

For any measure $\mu$, let $V=\operatorname{span}(\Phi_N(\operatorname{supp}\mu))$. Define the continuous linear form $L_s=I(s)^{-1}\int e^{-s\|z\|^2}\langle\Phi_N(z),A(s)(\cdot)\rangle\,d\mu(z)$ by the total operator-valued Bochner integral.
\end{definition}

\begin{lemma}[Weighted feature moments vanish on the weighted complement]
\label{lem:feature-moment-vanishing}
\lean{Preiss.featureMoment_eq_zero_of_mem_weightedComplement, Preiss.featureMoment_vanishes_weightedComplement}
\uses{def:feature-moment, def:orthogonal-grading, lem:integral-operator-annihilation}

For every measure $\mu$ on $\mathbb R^n$, every natural $N$, every real $s$, and every $u$ such that $A(s)u\in\operatorname{span}(\Phi_N(\operatorname{supp}\mu))^\perp$, the feature moment satisfies $L_s(u)=0$. The original unit-uniform, positive-parameter statement is retained as a wrapper.
\end{lemma}

\begin{lemma}[Nonzero vague limits of scalar rescalings are proportional]
\label{lem:vague-scalar-limits}
\lean{Preiss.exists_pos_smul_of_vague_and_scaled_vague}
\uses{gmt:vague, mathlib:regular-measure-ext}

On any locally compact Hausdorff space with its Borel measurable structure, suppose $\mu_i\rightharpoonup\nu$ and $c_i\mu_i\rightharpoonup\xi$, where $\nu,\xi$ are nonzero regular measures and $c_i\in[0,\infty]$. Then $\xi=a\nu$ for some positive finite $a$. No regularity assumption on the sequence is needed for the compact-test integral formulation.
\end{lemma}

\begin{lemma}[Positive finite scalar multiples preserve general tangents]
\label{lem:general-tangent-scalar}
\lean{Preiss.isTangentMeasure_smul}
\uses{def:tangent-api}

If $\nu\in\operatorname{Tan}(\mu,x)$ and $0<a<\infty$, then $a\nu\in\operatorname{Tan}(\mu,x)$. No local-finiteness assumption on the source measure is required.
\end{lemma}

\begin{lemma}[Dimensional subsequence extraction from two-sided ball bounds]
\label{lem:dimensional-extraction-ball-bounds}
\lean{Preiss.exists_dimensional_tangent_subseq_of_ball_bounds}
\uses{def:dimensional-tangent, lem:dimensional-tangent-operations, gmt:blowup-balls, gmt:portmanteau, thm:compactness}

Let $\mu$ be locally finite and suppose $p r^m\leq\mu(\overline B(x,r))\leq q r^m$ for all $0<r<R_0$, where $p,q,R_0>0$. Every positive sequence $r_i\to0$ admits a strictly increasing subsequence for which the dimensionally normalized blow-ups converge vaguely to a nonzero regular measure, which is also a general tangent. All natural dimensions, including zero, are allowed.
\end{lemma}

\begin{lemma}[Tangent-ray correspondence under density bounds]
\label{lem:tangent-ray-density-bounds}
\lean{Preiss.tangent_iff_dimensional_ray_of_ball_bounds, Preiss.tangent_iff_dimensional_ray_of_density_bounds}
\uses{lem:dimensional-extraction-ball-bounds, lem:vague-scalar-limits, lem:general-tangent-scalar, gmt:density-bounds}

For a locally finite measure with positive lower and finite upper $m$-density at $x$, $\nu\in\operatorname{Tan}(\mu,x)$ if and only if $\nu=c\xi$ for some $0<c<\infty$ and $\xi\in\operatorname{Tan}_m(\mu,x)$. A preceding version assumes only explicit two-sided local ball bounds. No existence of a density limit is used.
\end{lemma}

\begin{lemma}[Prescribed upper-density thresholds give quantitative ball bounds]
\label{lem:quantitative-upper-density-bound}
\lean{Preiss.exists_closedBall_power_bound_of_upperDensity_lt}
\uses{def:density-api}

For an outer measure $\nu$ on a pseudometric space, $m\in\mathbb N$ and $b>0$, if $\Theta^{*m}(\nu,x)<b$, then there is $R>0$ such that $\nu(\overline B_r(x))\le b\,2^m r^m$ for every $0<r<R$.
\end{lemma}

\begin{lemma}[Closed sets defined by uniform upper ball bounds]
\label{lem:closed-upper-ball-bound-sets}
\lean{Preiss.ballPowerUpperBoundSet, Preiss.isClosed_ballPowerUpperBoundSet, Preiss.ballPowerUpperBoundSet_mono}
\uses{lem:moving-set-measure-bounds}

For a measure $\mu$ on a pseudometric space, define $U_{m,b,R}=\{x: \mu(B_r(x))\le\operatorname{ofReal}(b r^m)\text{ for every }0<r<R\}$. This set is closed. If $b\le b^{\prime}$ and $R^{\prime}\le R$, then $U_{m,b,R}\subseteq U_{m,b^{\prime},R^{\prime}}$.
\end{lemma}

\begin{lemma}[Compactness of dimensional tangent families under ball bounds]
\label{lem:dimensional-family-compact-ball-bounds}
\lean{Preiss.dimensionalTangent_nonempty_compact_of_ball_bounds}
\uses{lem:dimensional-extraction-ball-bounds, def:vague-metrization, lem:vague-countable-metric}

Under positive two-sided local $m$-power closed-ball bounds at $x$, the family $\operatorname{Tan}_m(\mu,x)$ is nonempty and every sequence in it has a strictly increasing subsequence converging vaguely to another member. The source measure is locally finite; no density-limit existence is required.
\end{lemma}

\begin{lemma}[Radius-power limits from diameter-normalized densities]
\label{lem:radius-power-density}
\lean{Preiss.tendsto_radiusPowerRatio_of_densityRatio, Preiss.HasPositiveFiniteDensityAt.exists_radiusPowerDensity}
\uses{def:positive-density, def:density-api, gmt:density-bounds}

For outer measures on pseudometric spaces, convergence of the diameter-normalized $m$-density ratio to $\theta$ implies convergence of $\mu(\overline B(x,r))/r^m$ to $2^m\theta$. In particular, a positive finite density point of a measure on any metric space has a positive real radius-power limit $q$. No local-finiteness or Borel-space assumption is needed.
\end{lemma}

\begin{lemma}[Gaussian integrability of polynomial feature kernels]
\label{lem:gaussian-feature-integrability}
\lean{Preiss.integrable_gaussian_polynomialFeature, Preiss.integrable_featureMomentKernel_of_growth}
\uses{def:polynomial-feature, lem:gaussian-polynomial-integrable}

For a measure with a polynomial ball-growth bound, every finite polynomial feature vector and its weighted linear-functional kernel are integrable against $e^{-s\lVert z\rVert^2}$ for $s>0$. Uniformity and regularity are not needed.
\end{lemma}

\begin{lemma}[Gaussian integral from a centered ball law at arbitrary scale]
\label{lem:gaussian-ball-power-parameter}
\lean{Preiss.integral_gaussian_parameter_of_closedBall_power, Preiss.integral_gaussian_center_parameter_of_closedBall_power}
\uses{lem:gaussian-ball-power}

Let $\mu$ be an s-finite Borel measure on a real normed space with $\mu(\overline B(0,r))=Cr^d$ for $r>0$, where $C\geq0$ and $d>0$. For every $s>0$, the function $e^{-s\lVert z\rVert^2}$ is integrable and its integral is $C\Gamma(d/2+1)s^{-d/2}$.
\end{lemma}

\begin{lemma}[Support under continuous closed maps]
\label{lem:closed-map-support}
\lean{Preiss.measure_support_map_of_continuous_of_isClosedMap}

For a regular Borel measure $\mu$ and a continuous closed map $f$, the support of $f_*\mu$ is $f(\operatorname{supp}\mu)$.
\end{lemma}

\begin{lemma}[Outer regularity under continuous closed maps]
\label{lem:closed-map-outer-regularity}
\lean{Preiss.outerRegular_map_of_continuous_of_isClosedMap}

A continuous closed map pushes an outer regular Borel measure forward to an outer regular Borel measure.
\end{lemma}

\begin{lemma}[Regularity under proper maps]
\label{lem:proper-map-regularity}
\lean{Preiss.regular_map_of_isProperMap}
\uses{lem:closed-map-outer-regularity}

A proper map into a Hausdorff Borel space pushes a regular measure forward to a regular measure.
\end{lemma}

\begin{lemma}[Vague convergence under cocompact and proper maps]
\label{lem:vague-cocompact-pushforward}
\lean{Preiss.weaklyConverges_map_cocompact, Preiss.weaklyConverges_map_of_isProperMap}
\uses{gmt:vague}

Vague convergence is preserved by pushforward along a continuous cocompact map, and hence along a proper map. No regularity hypothesis is required for this test-function identity.
\end{lemma}

\begin{lemma}[Flat-measure regularity and positive plane-ball mass]
\label{lem:flat-measure-regularity}
\lean{Preiss.regular_flatOn, Preiss.flatOn_ball_pos}
\uses{lem:plane-hausdorff}

For every $m$-plane $V$, including $m=0$, $c\mathcal H^m\!\restriction V$ is regular when $c<\infty$. When $c>0$, every positive-radius ball centered in $V$ has positive measure.
\end{lemma}

\begin{lemma}[Vanishing mass on compact null sets]
\label{lem:vague-compact-null-mass}
\lean{Preiss.tendsto_mass_compact_null_of_vague}
\uses{gmt:portmanteau}

If regular Euclidean measures $\nu_i$ converge vaguely to a regular $\mu$ and $K$ is compact with $\mu(K)=0$, then $\nu_i(K)\to0$.
\end{lemma}

\begin{lemma}[Uniform lower ball mass on compact sets]
\label{lem:vague-uniform-ball-lower}
\lean{Preiss.eventually_uniform_ball_mass_lower_bound_of_vague}
\uses{gmt:portmanteau}

Let regular Euclidean measures $\nu_i$ converge vaguely to a regular $\mu$. If $K$ is compact, $\eta>0$, and $\mu(B(z,\eta/2))>0$ for every $z\in K$, then there exists $a>0$ such that eventually $\nu_i(B(z,\eta))\ge a$ simultaneously for every $z\in K$.
\end{lemma}

\begin{definition}[A common-plane two-sided approximation predicate]
\label{def:two-sided-approximation}
\lean{Preiss.TwoSidedApproximation}
\uses{def:cone}

The set $F$ is nonempty. At each $a\in G$ and scale $0<r<R$, a single $m$-plane gives tube containment of $F\cap B(a,2r)$, distance less than $\eta r$ from its points in $B(a,r)$ to $F$, and a $\gamma(\eta r)^m$ lower bound for the corresponding ball masses.
\end{definition}

\begin{lemma}[Compact two-sided approximation from uniform lower mass]
\label{lem:compact-two-sided-general}
\lean{Preiss.exists_compact_twoSided_approximation_of_uniform_lower_bound}
\uses{def:two-sided-approximation, lem:closed-weak-approximation-good-set, lem:weak-approximation-good-exhaustion, lem:closed-upper-ball-bound-sets, lem:compact-measurable-exhaustion, lem:tube-support-upgrade}

For a finite regular measure and a positive-measure compact $F$ with uniform lower ball bounds, assume weak linear approximation for $\mu\restriction F$ almost everywhere and arbitrarily small local upper ball bounds for $\mu\restriction F^c$ almost everywhere on $F$. For every $0<\eta<1$, there is a positive-measure compact $G\subseteq F$ with common-plane two-sided approximation at uniform positive scales and constants.
\end{lemma}

\begin{lemma}[A scalar Gaussian envelope for polynomial growth]
\label{lem:gaussian-power-envelope}
\lean{Preiss.pow_mul_gaussian_half_le}

For $r\geq0$, $s>0$ and integer $k>0$, $r^k e^{-sr^2/2}\leq(k/s)^{k/2}$.
\end{lemma}

\begin{lemma}[Integrated exponential series under a Gaussian weight]
\label{lem:integrated-gaussian-exponential-series}
\lean{Preiss.hasSum_integral_gaussian_exp_inner}
\uses{lem:gaussian-weight-exponential}

For a Euclidean measure with polynomial ball growth, every $s>0$, $t\in\mathbb R$ and vector $w$, the series of Gaussian-weighted moments $(t^k/k!)\int\langle z,w\rangle^k e^{-s\lVert z\rVert^2}\,d\mu$ converges absolutely to $\int e^{t\langle z,w\rangle}e^{-s\lVert z\rVert^2}\,d\mu$.
\end{lemma}

\begin{lemma}[Rescaled distances and approximate nearest-point selection]
\label{lem:rescaled-distance-selection}
\lean{Preiss.dist_affine_rescale, Preiss.exists_points_rescaled_tendsto_of_infDist}

In a real normed space, $\operatorname{dist}(a+ry,z)=r\operatorname{dist}(y,r^{-1}(z-a))$ for $r>0$. If $D$ is nonempty and $\operatorname{dist}(a+r_i y,D)/r_i\to0$ for positive $r_i$, one can choose $z_i\in D$ with $r_i^{-1}(z_i-a)\to y$. No closedness, compactness, or convergence of $r_i$ is required.
\end{lemma}

\begin{lemma}[Uniform small-scale bounds from subsequence properties]
\label{lem:sequential-uniform-scale}
\lean{Preiss.exists_uniform_scale_of_subsequence_property}

Let $P(a,r)$ be monotone downward in its positive parameter $a$. If every positive sequence $r_i\to0$ has a strictly increasing subsequence on which $P(a,r_i)$ eventually holds for some $a>0$, then there are $a,R>0$ such that $P(a,r)$ holds for every $0<r<R$.
\end{lemma}

\begin{lemma}[Homogeneity of distance to a linear plane]
\label{lem:plane-distance-homogeneous}
\lean{Preiss.infDist_smul_plane}
\uses{lem:projection-distance}

For a linear plane $V$, $\operatorname{dist}(cz,V)=|c|\operatorname{dist}(z,V)$ for every real $c$, including zero.
\end{lemma}

\begin{lemma}[Measure inequalities under dimensional normalization]
\label{lem:dimensional-blowup-set-estimates}
\lean{Preiss.dimensionalBlowup_apply_lt_iff, Preiss.le_dimensionalBlowup_apply_iff}
\uses{def:dimensional-blowup}

For every Borel set $S$, $r>0$ and $a\geq0$, upper and lower comparisons of $D^m_{x,r}\mu(S)$ with $a$ are equivalent to the corresponding comparisons of $\mu(T_{x,r}^{-1}S)$ with $ar^m$. No regularity or local finiteness is needed.
\end{lemma}

\begin{lemma}[Transfer of normalized plane approximation to a physical scale]
\label{lem:weak-approximation-scale-transfer}
\lean{Preiss.weakApproximation_scale_of_normalized_bounds}
\uses{lem:plane-distance-homogeneous, lem:dimensional-blowup-set-estimates, gmt:blowup-balls}

At a positive scale $r$, suppose $D^m_{x,r}\mu$ has mass less than $\eta$ outside the $\eta$-tube in the closed unit ball, and gives every $\eta$-ball centered on $V$ in that ball mass at least $a\geq0$. Then $\mu$ has the required weak plane-approximation estimates at $x,r$, with bounds $\eta r^m$ and $ar^m$.
\end{lemma}

\begin{lemma}[Second countability of countable-coordinate metrics]
\label{lem:countable-coordinate-second-countability}
\lean{Preiss.countableCoordinateMetric_secondCountable}
\uses{lem:countable-coordinate-metric}

The metric induced by countably many separating real coordinates is second countable, because its topology is induced from the countable product of real lines.
\end{lemma}

\begin{lemma}[Summable coefficient majorants and a global Gaussian-series remainder]
\label{lem:scalar-gaussian-series-remainder}
\lean{Preiss.summable_pow_mul_halfPower_div_factorial, Preiss.abs_power_series_tail_le, Preiss.abs_sum_positive_powers_le, Preiss.power_series_gaussian_taylor_remainder}

For $a\geq0$, the sequence $a^k k^{k/2}/k!$ is summable. Nonnegative summable coefficient majorants control power-series tails for $0\leq r\leq1$ and finite polynomials for $r\geq1$. If $\sum f_k=e^{r^2}$, $f_0=1$, and $|f_k|\leq c_k r^k$, then the discrepancy between the positive-degree truncations of $f$ and the Gaussian exponential is at most $(\sum c_k+e)r^{2q+1}$ for every $r\geq0$.
\end{lemma}

\begin{lemma}[A second-countable metrization of vague convergence]
\label{lem:vague-second-countable}
\lean{Preiss.exists_secondCountable_vagueMetrization}
\uses{def:vague-metrization-from-tests, lem:countable-coordinate-second-countability}

The regular measures on Euclidean space admit a second-countable metric whose sequential convergence is vague convergence.
\end{lemma}

\begin{lemma}[Metric approximation and separation of dimensional tangents]
\label{lem:dimensional-tangent-metric-separation}
\lean{Preiss.dimensionalTangent_of_metric_approximation, Preiss.exists_metric_separation_of_not_dimensionalTangent}
\uses{def:dimensional-tangent, def:vague-metrization, lem:dimensional-blowup-support}

A nonzero regular measure is a dimensional tangent if every metric neighborhood contains a dimensional blow-up at arbitrarily small positive scales. Consequently, a nontangent has a metric neighborhood separated from all sufficiently small blow-ups.
\end{lemma}

\begin{lemma}[Moving blow-up centers at the same dimensional scale]
\label{lem:dimensional-moving-centers}
\lean{Preiss.vague_dimensionalBlowup_moving_center}
\uses{lem:vague-moving-blowups, lem:blowup-composition, lem:dimensional-blowup-support}

If $D^m_{a,r_i}\mu\to\nu$ vaguely and $r_i^{-1}(a_i-a)\to y$, then $D^m_{a_i,r_i}\mu\to D^m_{y,1}\nu$. The original dimensional normalization is unchanged.
\end{lemma}

\begin{lemma}[Nonvanishing after recentering at a support point]
\label{lem:recenter-nonzero}
\lean{Preiss.dimensionalBlowup_one_ne_zero_of_mem_support}
\uses{lem:dimensional-blowup-support}

For every $y\in\operatorname{supp}\nu$, the unit-scale dimensional blow-up $D^m_{y,1}\nu$ is nonzero.
\end{lemma}

\begin{lemma}[Nullity of uniformly separated recentering witnesses]
\label{lem:recenter-separated-set-null}
\lean{Preiss.measure_eq_zero_of_recenter_separation}
\uses{gmt:density-points, lem:hull-support-approximation, lem:rescaled-distance-selection, lem:dimensional-moving-centers}

An arbitrary finite-outer-measure set of centers is null if each center has a recentered tangent witness within $\rho$ of one fixed regular measure while all its sufficiently small dimensional blow-ups are at distance at least $4\rho$ from that witness. No measurability of the center set is assumed.
\end{lemma}

\begin{lemma}[Operator bounds from pointwise bounds in finite dimensions]
\label{lem:finite-dimensional-operator-bigO}
\lean{Preiss.isBigO_continuousLinearMap_iff_forall_apply, Preiss.isBigO_restrict_span_of_forall_mem}

For a family of continuous linear maps on a finite-dimensional real normed space, an operator-norm big-$O$ bound is equivalent to the bound at each vector. Bounds on a generating set give an operator-norm bound after restriction to its span.
\end{lemma}

\begin{lemma}[Extending vague convergence using uniform integral tails]
\label{lem:vague-uniform-norm-tails}
\lean{Preiss.norm_integral_sub_cutoff_le, Preiss.vague_integral_of_uniform_norm_tails}
\uses{gmt:vague}

For a continuous integrable real function $f$, vague convergence implies convergence of its integrals whenever the norm integrals outside suitable compact sets are uniformly small for the sequence and limit measure. A cutoff between zero and one, equal to one on $K$, changes the integral by at most $\int_{K^c}|f|$.
\end{lemma}

\begin{lemma}[Polynomial norm growth outside the unit ball]
\label{lem:polynomial-norm-growth}
\lean{Preiss.exists_polynomial_norm_bound_outside_unit_ball}

Every real polynomial on Euclidean space admits constants $C\ge0$ and $d\in\mathbb N$ with $|P(z)|\le C\|z\|^d$ whenever $\|z\|\ge1$.
\end{lemma}

\begin{lemma}[Uniform tails for polynomial Gaussian tests]
\label{lem:polynomial-gaussian-norm-tails}
\lean{Preiss.polynomial_gaussian_uniform_norm_tail}
\uses{lem:polynomial-norm-growth, lem:gaussian-uniform-tails, lem:polynomial-gaussian-integrability}

For fixed polynomial $P$ and $s>0$, the integral of $|P(z)e^{-s\|z\|^2}|$ outside large balls tends to zero uniformly over regular measures with a common polynomial ball-growth bound.
\end{lemma}

\begin{lemma}[Centered ball laws from a positive radius-power density]
\label{lem:dimensional-tangent-origin-ball-law}
\lean{Preiss.dimensionalBlowup_ball_mass_tendsto, Preiss.dimensionalTangent_ball_law_at_zero}
\uses{def:dimensional-tangent, gmt:blowup-balls, gmt:portmanteau}

If $\mu(\overline B(x,r))/r^m\to q>0$, every dimensional tangent $\nu$ satisfies $\nu(\overline B(0,R))=qR^m$ for $R>0$ and has $0$ in its support. The ball-mass convergence for the blow-ups themselves also holds for $q=0$. Dimension zero is allowed.
\end{lemma}

\begin{lemma}[Flatness under positive finite scalar multiplication]
\label{lem:flat-positive-scalar}
\lean{Preiss.IsFlat.smul, Preiss.isFlat_smul_iff}
\uses{def:flat}

For $0<c<\infty$, the measure $c\nu$ is $m$-flat if and only if $\nu$ is $m$-flat.
\end{lemma}

\begin{lemma}[A Gaussian fourth-moment contraction for bilinear forms]
\label{lem:gaussian-bilinear-contraction}
\lean{Preiss.integrable_gaussian_bilinear_diagonal, Preiss.integrable_gaussian_inner_sq_smul_bilinear, Preiss.gaussian_integral_inner_sq_mul_bilinear}
\uses{lem:gaussian-wick-fourth, lem:gaussian-quadratic-trace-bound}

For a continuous real bilinear form $A$ on a finite-dimensional real inner-product space and standard Gaussian $\gamma$, $\int\langle y,v\rangle^2 A(v,v)\,d\gamma=\|y\|^2\operatorname{tr}A+2A(y,y)$. Symmetry of $A$ is unnecessary. The diagonal and weighted-diagonal integrability lemmas allow normed real vector-valued forms.
\end{lemma}

\begin{lemma}[A lower packing mass exceeding an upper bound is impossible]
\label{lem:column-mass-contradiction}
\lean{Preiss.false_of_column_mass_exceeds_upper_bound}

If nonnegative retained radii have total $m$-power at least $b$, and a measure has mass at least $Ab$ but at most a strictly smaller finite bound, the two estimates contradict one another.
\end{lemma}

\begin{lemma}[Positivity and complementary subspaces for weighted gradings]
\label{lem:weighted-grading-positive-complement}
\lean{Preiss.weightedGrading_inner_self, Preiss.weightedGrading_inner_pos, Preiss.weightedGrading_injective, Preiss.weightedGrading_inner_left_right, Preiss.isCompl_weightedComplement}
\uses{def:orthogonal-grading}

For an orthogonal grading $(P_i)$ and positive $s$, the weighted operator $A(s)=\sum_i s^{i+1}P_i$ is self-adjoint and positive definite, with quadratic form $\sum_i s^{i+1}\|P_i x\|^2$. It is injective, and every subspace $V$ is complementary to $A(s)^{-1}(V^\perp)$.
\end{lemma}

\begin{lemma}[Isometric parametrization of a linear plane]
\label{lem:plane-isometric-parametrization}
\lean{Preiss.exists_linearIsometry_range_plane}
\uses{def:plane}

Every $m$-dimensional plane in $\mathbb R^n$ is the range of a linear isometric embedding of $\mathbb R^m$, including $m=0$.
\end{lemma}

\begin{lemma}[Support and regularity of normalized flat measures]
\label{lem:unit-flat-support-regularity}
\lean{Preiss.flatOn_support, Preiss.regular_unitFlat, Preiss.unitFlat_support, Preiss.isFlat_unitFlat, Preiss.unitFlat_ne_zero}
\uses{lem:flat-measure-regularity, lem:measure-support-transport, lem:plane-hausdorff, def:flat}

A positive multiple of plane Hausdorff measure has support exactly that plane. The normalized plane measure is regular and has the same support, in every dimension.
\end{lemma}

\begin{lemma}[Normalized flat measures are unit uniform]
\label{lem:unit-flat-unit-uniform}
\lean{Preiss.isUnitUniform_unitFlat}
\uses{lem:unit-flat-support-regularity, lem:plane-hausdorff, def:unit-uniform}

The normalized measure on every linear $m$-plane is unit $m$-uniform, with mass $r^m$ on every positive-radius closed ball centered in that plane.
\end{lemma}

\begin{lemma}[The packing contradiction from compact two-sided approximations]
\label{lem:two-sided-packing-core}
\lean{Preiss.false_of_compact_twoSided_approximation}
\uses{def:two-sided-approximation, lem:projection-hole-cover, lem:discarded-cylinder-sum, lem:column-packing, lem:column-mass-contradiction}

A measure cannot simultaneously have a uniform lower $m$-power bound on $F$, arbitrarily small-aperture compact two-sided approximations inside $F$, base points with a common upper mass constant and arbitrarily small deleted mass, and null $m$-dimensional projections of $F$ onto every $m$-plane. The proof does not require the measure itself to be Hausdorff measure.
\end{lemma}

\begin{lemma}[Orthogonal layers of a descending subspace flag]
\label{lem:nested-orthogonal-layers}
\lean{Preiss.starProjection_comp_of_ge, Preiss.nestedProjectionLayer, Preiss.nestedProjectionLayer_symmetric, Preiss.nestedProjectionLayer_mem, Preiss.nestedProjectionLayer_eq_zero_of_mem, Preiss.sub_nestedProjectionLayer_mem, Preiss.nestedProjectionLayer_idempotent, Preiss.nestedProjectionLayer_comp_eq_zero, Preiss.sum_nestedProjectionLayer}

For a descending finite flag $(U_i)$ in a finite-dimensional inner-product space, differences of successive orthogonal projections are self-adjoint idempotents with pairwise orthogonal ranges. Their sum telescopes to the difference of the endpoint projections. Residual and annihilation identities describe each layer.
\end{lemma}

\begin{lemma}[The kernel flag induced on a subspace by an orthogonal grading]
\label{lem:kernel-grading-flag}
\lean{Preiss.kernelGradingFlag, Preiss.mem_kernelGradingFlag, Preiss.kernelGradingFlag_antitone, Preiss.kernelGradingFlag_zero, Preiss.kernelGradingFlag_end}
\uses{def:orthogonal-grading}

Given an orthogonal grading $(P_i)$ and subspace $V$, the subspaces $K_i=\{v\in V:P_jv=0\text{ for }j<i\}$ form a descending flag from $V$ to zero.
\end{lemma}

\begin{lemma}[Orthogonal grading by kernel-flag layers]
\label{lem:kernel-flag-layer-grading}
\lean{Preiss.flagProjectionLayer, Preiss.flagProjectionLayer_isOrthogonalGrading}
\uses{lem:kernel-grading-flag, lem:nested-orthogonal-layers}

The differences of successive orthogonal projections onto the kernel flag form an orthogonal grading on $V$.
\end{lemma}

\begin{lemma}[Injectivity of the constant divided grading operator]
\label{lem:divided-grading-constant-injective}
\lean{Preiss.compressedGradingPiece, Preiss.flagProjectionLayer_compressed_inner, Preiss.dividedGradingConstant, Preiss.flagProjectionLayer_dividedGradingConstant, Preiss.dividedGradingConstant_injective}
\uses{lem:kernel-flag-layer-grading, lem:kernel-grading-flag}

Let $D_i$ be the kernel-flag layers and $R_i$ the compression of $P_i$ to $V$. The operator $\sum_i D_iR_i$ is injective. The proof uses $\langle D_iR_ix,x\rangle=\|P_ix\|^2$ on $K_i$, layer orthogonality, and induction through the flag.
\end{lemma}

\begin{lemma}[Uniqueness of full-dimensional uniform measures with fixed ball constant]
\label{lem:uniform-full-dimension-unique}
\lean{Preiss.uniformWithConstant_eq_of_fullDimension}
\uses{gmt:uniform-full-dimension, lem:uniform-scaling}

In positive dimension $m$, two $m$-uniform measures on $\mathbb R^m$ with the same positive ball constant are equal. This reuses FormalizingGMT\textquotesingle s full-dimensional classification and identifies the scalar by the unit-ball mass.
\end{lemma}

\begin{lemma}[Taylor coefficient vanishing from a growth exponent]
\label{lem:asymptotic-taylor-vanishing}
\lean{Preiss.isLittleO_rpow_nhdsGT_zero, Preiss.eq_zero_of_isLittleO_pow_smul, Preiss.taylorCoefficient_eq_zero_of_lt_growth}

For positive scales, $s^a=o(s^b)$ when $b<a$. A constant vector whose multiple $s^qv$ is $o(s^q)$ must vanish. Consequently, a normed-space-valued function bounded by $O(s^a)$ and admitting Taylor expansions has zero coefficients in all positive integer degrees below $a$.
\end{lemma}

\begin{lemma}[Operator-norm growth of Gaussian moments at zero]
\label{lem:moment-operator-growth}
\lean{Preiss.gaussianMoment_isBigO_rpow}
\uses{lem:moment-norm-bound}

For a positive-dimensional unit uniform measure and $k>0$, its $k$-linear Gaussian moment is $O(s^{k/2})$ in operator norm as $s\downarrow0$.
\end{lemma}

\begin{lemma}[Polynomial division along a triangular grading]
\label{lem:graded-polynomial-division}
\lean{Preiss.dividedGradingMap, Preiss.dividedGradingMap_row, Preiss.weightedGrading_comp_dividedGradingMap, Preiss.dividedGradingMap_zero, Preiss.analyticAt_dividedGradingMap}
\uses{def:orthogonal-grading}

For an orthogonal grading $(D_i)$ and operators $R_j$ with $D_iR_j=0$ when $j<i$, the polynomial $H(s)=\sum_{i\le j}s^{j-i}D_iR_j$ satisfies $(\sum_i s^{i+1}D_i)H(s)=\sum_j s^{j+1}R_j$. Its rows and value at zero are explicit, and it is analytic at every real parameter.
\end{lemma}

\begin{lemma}[Analytic divided operators for the kernel flag]
\label{lem:kernel-flag-divided-operators}
\lean{Preiss.projectedGradingPiece, Preiss.flagProjectionLayer_projectedGradingPiece_eq_zero, Preiss.dividedFlagMap, Preiss.dividedFlagOperator, Preiss.dividedFlagOperator_zero, Preiss.weightedGrading_comp_dividedFlagMap, Preiss.analyticAt_dividedFlagMap, Preiss.analyticAt_dividedFlagOperator, Preiss.isUnit_dividedFlagOperator_zero}
\uses{lem:graded-polynomial-division, lem:kernel-flag-layer-grading, lem:divided-grading-constant-injective}

Compressing the original grading to a subspace gives the triangular relations needed for polynomial division. The resulting divided map and its restriction to the subspace are analytic; their weighted factorization is exact, and the restricted operator is invertible at zero.
\end{lemma}

\begin{lemma}[Invertibility of the divided operator at positive scales]
\label{lem:divided-operator-positive-inverse}
\lean{Preiss.isUnit_dividedFlagOperator_pos}
\uses{lem:kernel-flag-divided-operators, lem:weighted-grading-positive-complement}

The restricted divided kernel-flag operator is invertible at every positive parameter, by its weighted factorization and positivity of the original weighted grading.
\end{lemma}

\begin{lemma}[Uniform Hausdorff bounds for slabs in linear planes]
\label{lem:plane-slab-measure}
\lean{Preiss.hausdorffMeasure_plane_slab_le}
\uses{def:plane}

For $m>0$, a universal constant bounds the $m$-dimensional Hausdorff measure of $\{z\in V:\|z-c\|\le R,\ |\langle v,z-c\rangle|\le dR\}$ by $CdR^m$, for every $m$-plane, center in that plane, unit normal in that plane, and nonnegative $R,d$.
\end{lemma}

\begin{lemma}[A projected tube lies in a thin plane slab]
\label{lem:projected-tube-contained-slab}
\lean{Preiss.projected_tube_subset_slab}
\uses{def:cone, lem:projection-distance}

If a unit vector $v\in V$ has $\|P_Wv\|\le\varepsilon$, the projection onto $V$ of the radius-$R$ portion of an $\eta R$-tube about $a+W$ lies in the radius-$R$ slab centered at $P_Va$ of relative thickness $\varepsilon+\eta$.
\end{lemma}

\begin{lemma}[Image-measure bounds from a buffered Vitali cover]
\label{lem:fivefold-image-covering}
\lean{Preiss.measure_image_le_mul_of_ball_image_bounds}

In a separable pseudometric space with measurable open sets, assume every point of $C$ admits a positive radius bounded by a common upper bound and $\nu(f(C\cap B(x,5r)))\le K\mu(B(x,r))$. A disjoint Vitali subfamily gives $\nu(f(C))\le K\mu(X)$. No measurability of $f$ or $C$ is needed.
\end{lemma}

\begin{lemma}[Unit uniform dimensional tangents at every support point]
\label{lem:unit-uniform-tangent-existence}
\lean{Preiss.exists_unitUniform_dimensionalTangent}
\uses{lem:uniform-compactness, lem:uniform-scaling, def:dimensional-tangent}

A unit $m$-uniform measure admits a unit $m$-uniform dimensional tangent at every support point, including $m=0$.
\end{lemma}

\begin{lemma}[Finite Taylor jets of normed-space-valued analytic functions]
\label{lem:analytic-finite-jets}
\lean{Preiss.exists_taylor_jets_of_analyticAt_zero, Preiss.exists_positive_taylor_jets_of_analyticAt_zero}

Every real-parameter analytic function into a real normed space has one coefficient sequence whose degree-$q$ truncation differs from the function by $o(s^q)$ near zero, for every $q$. Its constant coefficient is the value at zero. If that value is zero, the truncations start in degree one.
\end{lemma}

\begin{lemma}[Extending analytic approximations from a subspace]
\label{lem:analytic-projection-approximation}
\lean{Preiss.functional_comp_projection_eq_self, Preiss.exists_analyticApproximation_of_weightedComplement}
\uses{lem:graded-analytic-projection}

Let $f(s)$ be continuous linear maps into a normed space annihilating the weighted complement of $V$ for positive $s$. If an analytic family $g(s)$, vanishing at zero, approximates $f(s)$ on $V$ with operator error $O(r(s))$, then an analytic family vanishing at zero approximates $f(s)$ on the whole space with the same error order.
\end{lemma}

\begin{lemma}[Dimensional tangents lie in the bounded-radius scaling orbit closure]
\label{lem:dimensional-tangent-bounded-orbit}
\lean{Preiss.IsDimensionalTangent.isUniformPseudotangent}
\uses{def:dimensional-tangent, lem:pseudotangent-transitive}

Every dimensional tangent at a support point lies in the support-centered scaling orbit closure with radii at most one. The proof discards finitely many large radii; no uniformity of the source measure is needed.
\end{lemma}

\begin{lemma}[Small projected directions from plane proximity and cone escape]
\label{lem:cone-projection-direction}
\lean{Preiss.exists_unit_small_projection_of_near_plane, Preiss.norm_projection_lt_of_outside_cone}
\uses{def:cone, lem:projection-distance}

A nonzero displacement close to a plane and with small projection onto another plane yields a unit direction in the first plane with quantitatively small projection. Escaping a cone supplies the required nonzero displacement and projection bound.
\end{lemma}

\begin{lemma}[Quantitative projection bounds from uniform tubes and cone escape]
\label{lem:uniform-tube-projection-cover}
\lean{Preiss.measure_projection_le_of_uniform_tubes_and_cone_escape}
\uses{lem:fivefold-image-covering, lem:projected-tube-slab-bound, lem:cone-projection-direction}

If $D\subseteq C$ has uniform lower ball mass $\delta r^m$, the carrier $C$ lies in sufficiently thin approximating tubes around points of $D$, and every such point admits arbitrarily small cone-escaping displacements, then $\mathcal H^m(P_VD)\le (K\varepsilon/\delta)\mu(X)$. The fivefold-cover estimate uses the tube information at a larger scale, preserving the entire image mass.
\end{lemma}

\begin{lemma}[Coherent Taylor coefficients from compatible finite jets]
\label{lem:coherent-finite-jets}
\lean{Preiss.exists_coherent_taylor_coefficients_of_finite_jets}
\uses{lem:asymptotic-taylor-vanishing}

Let $f$ take values in a real normed space. Suppose every positive order $q$ admits coefficients giving $o(s^j)$ remainders for all $1\leq j\leq q$, as $s\downarrow0$. Then one coefficient sequence gives these remainders simultaneously at every positive order. Uniqueness of the first differing coefficient yields compatibility.
\end{lemma}

\begin{lemma}[Analytic approximation of the complete feature moment functional]
\label{lem:feature-moment-analytic-approximation}
\lean{Preiss.exists_featureMoment_analytic_approximation}
\uses{lem:analytic-projection-approximation, lem:feature-norm-functionals, lem:moment-feature-remainder, lem:feature-moment-vanishing}

For every unit-uniform measure of positive dimension and every positive integer $q$, the degree-$2q$ feature moment functional admits an analytic approximation vanishing at zero, with operator error $O(s^{q+1/2})$ as $s\downarrow0$.
\end{lemma}

\begin{lemma}[Compact tube uniformization preserving a fixed lower density bound]
\label{lem:compact-uniform-tubes}
\lean{Preiss.exists_compact_uniform_tubes_subset}
\uses{lem:uniform-compact-selection, lem:tube-support-upgrade}

For a finite regular measure, a measurable carrier with a fixed positive lower ball-mass bound and weak linear approximation almost everywhere contains a compact subset losing arbitrarily little measure, with uniform small-scale tube containment at any prescribed positive aperture at most one.
\end{lemma}

\begin{lemma}[A flat dimensional tangent from uniformity and simultaneous recentering]
\label{lem:flat-tangent-from-uniform-recentering}
\lean{Preiss.exists_flat_dimensionalTangent_of_uniform_recenter}
\uses{thm:flat-pseudotangent, lem:uniform-scaling, lem:blowup-composition, lem:dimensional-tangent-closed, lem:unit-flat-support-regularity}

For $0<m\leq n$, suppose a $q$-uniform measure has zero in its support and every support-centered dimensional rescaling is a dimensional tangent of a locally finite measure $\mu$ at $x$. There is an $m$-plane $V$ such that $q$ times the normalized plane measure on $V$ is a dimensional tangent of $\mu$ at $x$. The coefficient $q$ is retained throughout the limiting argument.
\end{lemma}

\begin{definition}[Symmetric multilinear forms and permutation averaging]
\label{def:symmetric-multilinear}
\lean{Preiss.symmetricContinuousMultilinearSubmodule, Preiss.symmetrizeContinuousMultilinearMap}

Symmetric continuous multilinear maps form a concrete real submodule. Symmetrization is the average over all permutations of the finite index set.
\end{definition}

\begin{lemma}[Symmetrization identities]
\label{lem:multilinear-symmetrization}
\lean{Preiss.symmetrizeContinuousMultilinearMap_apply_perm, Preiss.symmetrizeContinuousMultilinearMap_eq_self, Preiss.symmetrizeContinuousMultilinearMap_diagonal, Preiss.comp_symmetrizeContinuousMultilinearMap}
\uses{def:symmetric-multilinear}

Permutation averaging produces a symmetric map, fixes every symmetric map, preserves diagonal evaluation, and commutes with continuous linear maps on the target.
\end{lemma}

\begin{lemma}[Multilinear coordinate tensors for polynomial feature blocks]
\label{lem:feature-multilinear-block}
\lean{Preiss.polynomialFeatureMultilinearBlock, Preiss.polynomialFeatureMultilinearBlock_apply, Preiss.inner_polynomialFeature_multilinearBlock, Preiss.polynomialFeatureMultilinearBlock_mem, Preiss.weightedGrading_polynomialFeatureMultilinearBlock}
\uses{def:polynomial-feature, lem:polynomial-feature-identities}

The degree-$k$ coordinate tensor takes values in its polynomial feature block. Its inner product with a feature vector is the normalized product of $k$ inner products, and weighted grading multiplies it by $s^k$.
\end{lemma}

\begin{definition}[Bounded recovery of symmetric forms from feature functionals]
\label{def:symmetric-feature-recovery}
\lean{Preiss.symmetricPolynomialFeatureMultilinearBlock, Preiss.polynomialFeatureMomentRecovery, Preiss.polynomialFeatureMomentRecovery_symmetric, Preiss.symmetricPolynomialFeatureMomentRecovery}
\uses{lem:feature-multilinear-block, lem:multilinear-symmetrization}

Composition with the symmetrized coordinate tensor is a bounded linear map from feature functionals to symmetric multilinear forms. Its range is corestricted to the concrete symmetric submodule.
\end{definition}

\begin{lemma}[Exact Gaussian moment recovery from polynomial features]
\label{lem:feature-moment-exact-recovery}
\lean{Preiss.featureMoment_comp_multilinearBlock_of_growth, Preiss.polynomialFeatureMomentRecovery_featureMoment_of_growth}
\uses{def:symmetric-feature-recovery, lem:gaussian-feature-integrability, lem:moment-form-symmetry}

For a measure with polynomial ball growth and positive Gaussian parameter, applying the degree-$k$ recovery map to its feature moment functional recovers its Gaussian multilinear moment exactly.
\end{lemma}

\begin{definition}[Gaussian moments as elements of the symmetric-form submodule]
\label{def:symmetric-moment-form}
\lean{Preiss.SymmetricMomentForm, Preiss.symmetricGaussianMomentForm}
\uses{def:symmetric-multilinear, def:moments, lem:moment-form-symmetry}

The symmetric moment space is the concrete submodule of permutation-invariant continuous multilinear forms. Gaussian moment forms lie in it by the proved permutation identity.
\end{definition}

\begin{lemma}[Finite Gaussian moment jets in the symmetric-form space]
\label{lem:finite-moment-taylor-jets}
\lean{Preiss.exists_gaussianMoment_finite_jets}
\uses{def:symmetric-moment-form, lem:feature-moment-analytic-approximation, def:symmetric-feature-recovery, lem:feature-moment-exact-recovery, lem:analytic-finite-jets, lem:asymptotic-taylor-vanishing}

For every positive moment degree and positive truncation order $q$, a unit-uniform measure of positive dimension admits symmetric multilinear coefficients giving $o(s^j)$ remainders simultaneously for every $1\leq j\leq q$. Recovery from sufficiently many polynomial features preserves symmetry and operator estimates.
\end{lemma}

\begin{lemma}[Reflection invariance of flat measures]
\label{lem:flat-measure-reflection}
\lean{Preiss.flatOn_map_neg, Preiss.unitFlat_map_neg}
\uses{lem:plane-hausdorff}

Every scalar flat measure and every normalized plane measure is invariant under reflection $z\mapsto-z$, including dimension zero.
\end{lemma}

\begin{lemma}[Odd Gaussian moments of reflection-invariant measures]
\label{lem:reflection-odd-gaussian-moments}
\lean{Preiss.gaussianMomentForm_eq_zero_of_odd_of_map_neg}
\uses{def:moments, lem:moment-kernel-calculus}

Every odd Gaussian multilinear moment of a reflection-invariant measure vanishes. No integrability, uniformity, or positivity assumption on the Gaussian parameter is required.
\end{lemma}

\begin{lemma}[Odd Gaussian moments of flat measures vanish]
\label{lem:flat-gaussian-odd-moments}
\lean{Preiss.flat_gaussianMoment_odd}
\uses{lem:flat-measure-reflection, lem:reflection-odd-gaussian-moments}

All odd Gaussian multilinear moments of a normalized plane measure vanish, including for zero-dimensional planes.
\end{lemma}

\begin{lemma}[Projection nullness with fixed lower ball mass]
\label{lem:fixed-lower-density-projection-null}
\lean{Preiss.measure_projection_eq_zero_of_uniformLower_and_weakApproximation}
\uses{lem:compact-uniform-tubes, lem:uniform-tube-projection-cover, lem:pure-cone-escape}

A measurable purely unrectifiable subset of a finite Hausdorff-measure carrier, with a uniform positive lower ball-mass bound and weak linear approximation almost everywhere, has zero projection measure onto every plane of the same dimension.
\end{lemma}

\begin{lemma}[Vanishing weighted projection energy controls the vague-limit support]
\label{lem:weighted-energy-limit-support}
\lean{Preiss.support_subset_fixedPoints_of_vague_of_weighted_projectionEnergy_tendsto_zero}
\uses{lem:vague-projection-energy}

Let regular Euclidean measures $\mu_i$ converge vaguely to $\nu$, and continuous linear maps $P_i$ converge to $Q$. For a strictly positive continuous weight $w$, suppose $w(z)\lVert z-P_i z\rVert^2$ is integrable for each $\mu_i$ and its integral tends to zero. Then the support of $\nu$ lies in $\{z:Qz=z\}$. No global integrability assumption on $\nu$ is required.
\end{lemma}

\begin{lemma}[Gaussian plane energy is integrable for unit-uniform measures]
\label{lem:gaussian-plane-energy-integrable}
\lean{Preiss.integrable_gaussian_planeEnergy_of_unitUniform}
\uses{lem:moment-norm-bound, lem:gaussian-polynomial-integrable}

For every unit-uniform measure $\mu$ and $m$-plane $V$, the function $e^{-\lVert z\rVert^2}\operatorname{dist}(z,V)^2$ is $\mu$-integrable. This includes $m=0$.
\end{lemma}

\begin{lemma}[Density bounds inherited under absolute continuity]
\label{lem:absolute-continuity-density-bounds}
\lean{Preiss.ae_hasPositiveFiniteDensitiesAt_of_absolutelyContinuous}
\uses{def:positive-density, mathlib:besicovitch-differentiation}

If locally finite Euclidean measures satisfy $\nu\ll\mu$ and $\mu$ has positive finite lower and upper $m$-densities almost everywhere, then $\nu$ has the same density-bound property at $\nu$-almost every point. The proof uses positive finite Radon--Nikodym derivatives and differentiation along closed balls.
\end{lemma}

\begin{lemma}[Tangent cones agree under finite absolute continuity]
\label{lem:absolute-continuity-tangent-cones}
\lean{Preiss.ae_isTangentMeasure_iff_of_absolutelyContinuous}
\uses{thm:matching}

For locally finite Euclidean $\mu$ and finite $\nu\ll\mu$, their full tangent cones agree at $\nu$-almost every point. Apply the pinned FormalizingGMT density-matching theorem to the real Radon--Nikodym derivative.
\end{lemma}

\begin{lemma}[Orthogonal representatives and normal-coordinate identities]
\label{lem:normal-coordinate-identities}
\lean{Preiss.exists_orthogonal_representative, Preiss.normal_coordinates_of_quadratic_identity}

A continuous linear functional vanishing on a closed subspace has a representing vector in its orthogonal complement. If $2\langle b,x\rangle+\lVert P_Vx\rVert^2=\lVert x\rVert^2$ with $b\perp V$, then $\lVert x-P_Vx\rVert^2=2\langle b,x\rangle$, $\lVert x-P_Vx-b\rVert=\lVert b\rVert$, and $\lVert x-P_Vx\rVert\leq2\lVert b\rVert$.
\end{lemma}

\begin{lemma}[Even diagonal Gaussian moments are nonnegative]
\label{lem:even-moment-positivity}
\lean{Preiss.gaussianMomentForm_even_nonneg}
\uses{def:moments, lem:moment-form-evaluation}

Every even diagonal Gaussian moment is nonnegative, for any measure and real Gaussian parameter. The statement also covers the nonintegrable case in the Bochner integral definition.
\end{lemma}

\begin{lemma}[Exact moment homogeneity for conical measures]
\label{lem:conical-moment-scaling}
\lean{Preiss.gaussianMomentForm_homogeneous_of_isConical}
\uses{def:conical, lem:moment-covariance}

For every conical measure, its degree-$k$ Gaussian multilinear moment at positive parameter $s$ equals $s^{k/2}$ times its moment at parameter one. No uniformity assumption is required.
\end{lemma}

\begin{lemma}[Conical covariance equals squared norm on the support]
\label{lem:conical-support-covariance}
\lean{Preiss.conical_covariance_eq_norm_sq_on_support}
\uses{lem:conical-moment-scaling, lem:scaled-moment-limits, lem:taylor-two, lem:moment-support-polynomial}

For a positive-dimensional unit-uniform conical measure, its second Gaussian moment at parameter one equals squared norm at every support point.
\end{lemma}

\begin{lemma}[The second-moment trace as a bounded linear functional]
\label{lem:moment-trace-operator}
\lean{Preiss.momentTraceCLM, Preiss.momentTraceCLM_apply, Preiss.momentTraceCLM_momentKernel}
\uses{lem:bilinear-trace-basis, def:bilinear-trace, def:moments}

The finite sum of diagonal evaluations defines a continuous linear functional on second moment forms. It equals the associated bilinear trace and maps the moment kernel at $z$ to $\lVert z\rVert^2$.
\end{lemma}

\begin{lemma}[Radial second Gaussian moment of a unit-uniform measure]
\label{lem:gaussian-radial-second-moment}
\lean{Preiss.integral_norm_sq_gaussian_unitUniform}
\uses{lem:moment-norm-bound, lem:gaussian-parameter-derivative}

For a positive-dimensional unit-uniform measure and $s>0$, the integral of $\lVert z\rVert^2e^{-s\lVert z\rVert^2}$ is $m/(2s)$ times its Gaussian mass. Differentiate the explicit Gaussian mass formula under the integral.
\end{lemma}

\begin{lemma}[Finite ordered trace bounds and extremizers]
\label{lem:ordered-trace-bounds}
\lean{Preiss.ordered_nonnegative_trace_cutoff_bounds, Preiss.ordered_nonnegative_trace_eq_initial_indicator, Preiss.antitone_weighted_sum_le_initial_sum}

An antitone nonnegative tuple of total $m$ has its $m$-th entry at most one; equality forces the indicator of its first $m$ entries. For any antitone real tuple and weights in $[0,1]$ summing to $m$, the weighted sum is at most the sum of its first $m$ entries.
\end{lemma}

\begin{lemma}[Quadratic rigidity near a critical basis vector]
\label{lem:quadratic-near-basis-rigidity}
\lean{Preiss.ordered_quadratic_identity_rigid_of_near_basis_vector}
\uses{lem:ordered-trace-bounds}

In a real inner-product space with a finite orthonormal basis, an ordered nonnegative spectrum of total $m$ is the rank-$m$ indicator spectrum if its quadratic identity holds at a point within distance $\delta<1/m$ of the $m$-th basis vector.
\end{lemma}

\begin{lemma}[Taylor coefficients of an exact monomial]
\label{lem:exact-monomial-coefficients}
\lean{Preiss.taylorCoefficient_eq_of_exact_monomial}
\uses{lem:asymptotic-taylor-vanishing}

If a normed-space-valued function equals $s^d a$ for positive $s$ and admits compatible Taylor remainders of every order, its only nonzero positive-order Taylor coefficient is $a$ in degree $d>0$.
\end{lemma}

\begin{lemma}[A normal vector from first-moment plane vanishing]
\label{lem:first-moment-normal-from-vanishing}
\lean{Preiss.exists_firstMomentNormalVector_of_restrict_eq_zero}
\uses{lem:normal-coordinate-identities, def:rigidity-normal-vector}

If the first coefficient form vanishes on a plane, half its Riesz representative lies in the orthogonal complement and represents that form as $x\mapsto2\langle b,x\rangle$.
\end{lemma}

\begin{lemma}[Inverse-square scales approach zero through positive values]
\label{lem:inverse-square-positive-limit}
\lean{Preiss.tendsto_inv_sq_atTop_nhdsGT_zero}

As $r\to\infty$, the scale $r^{-2}$ tends to zero through positive values.
\end{lemma}

\begin{lemma}[Gaussian product integrals and moment forms on unit-uniform measures]
\label{lem:gaussian-product-moment-conversion}
\lean{Preiss.gaussian_integral_product_eq_moment, Preiss.integrable_polynomial_gaussian_unitUniform}
\uses{lem:moment-norm-bound, lem:polynomial-gaussian-integrability}

For a unit-uniform measure, Gaussian integrals of products of linear functionals equal explicit scalar multiples of their moment forms. Polynomial Gaussian tests are integrable.
\end{lemma}

\begin{lemma}[Coordinate monomials are finite products of linear inner products]
\label{lem:coordinate-product-representation}
\lean{Preiss.exists_coordinate_product_representation}

Every coordinate monomial can be represented as a finite product of inner products with coordinate basis vectors.
\end{lemma}

\begin{lemma}[Gaussian linear-product tests have blowdown limits]
\label{lem:blowdown-product-limits}
\lean{Preiss.exists_blowdown_product_gaussian_limit}
\uses{lem:blowdown-form-limits, lem:gaussian-product-moment-conversion, lem:uniform-scaling}

Gaussian integrals of every finite product of linear functionals have a limit along dimensional blowdowns of a positive-dimensional unit-uniform measure.
\end{lemma}

\begin{lemma}[Every positive-degree Gaussian moment has a blowdown limit]
\label{lem:blowdown-form-limits}
\lean{Preiss.exists_blowdown_gaussianMoment_limit}
\uses{lem:inverse-square-positive-limit, lem:taylor-two, lem:scaled-moment-limits, lem:moment-covariance}

Every positive-degree Gaussian multilinear moment at parameter one has a limit along the dimensional blowdowns of a positive-dimensional unit-uniform measure. The even limits come from leading Taylor coefficients; the odd limits vanish.
\end{lemma}

\begin{lemma}[Orthonormal sums are bounded by leading eigenvalues]
\label{lem:orthonormal-spectral-sum-bound}
\lean{Preiss.orthonormal_diagonal_sum_le_initial_sum}
\uses{lem:ordered-trace-bounds}

For an ordered real spectrum, the sum of its quadratic form over any orthonormal family of size $m$ is bounded by the sum of its first $m$ eigenvalues.
\end{lemma}

\begin{lemma}[Ordered spectrum of a positive symmetric bilinear form]
\label{lem:ordered-bilinear-spectrum}
\lean{Preiss.exists_ordered_spectral_representation}
\uses{def:bilinear-trace}

A positive symmetric bilinear form on a finite-dimensional real inner-product space admits an orthonormal eigenbasis, with nonnegative decreasing eigenvalues summing to its trace.
\end{lemma}

\begin{lemma}[Distance to a plane in orthonormal coordinates]
\label{lem:orthonormal-plane-distance}
\lean{Preiss.infDist_sq_eq_norm_sq_sub_sum_orthonormalBasis, Preiss.exists_plane_of_orthonormalBasis_finset}

Squared distance to a subspace equals squared norm minus the sum of squared coordinates in an orthonormal basis. Selecting $m$ vectors of an ambient orthonormal basis constructs an $m$-plane with this distance formula.
\end{lemma}

\begin{lemma}[Explicit even moment limits of blowdowns]
\label{lem:blowdown-even-moment-limit}
\lean{Preiss.blowdown_gaussianMoment_even_limit}
\uses{lem:scaled-moment-limits, lem:inverse-square-positive-limit, lem:moment-covariance}

The degree $2j$ Gaussian moment of dimensional blowdowns converges to $B_{2j}^{j}/j!$ for every positive $j$.
\end{lemma}

\begin{lemma}[Gaussian moment continuity under vague convergence]
\label{lem:vague-gaussian-moments}
\lean{Preiss.vague_gaussianMomentForm_apply}
\uses{lem:vague-gaussian-polynomials, lem:moment-norm-bound}

Along vaguely convergent unit-uniform measures, every Gaussian multilinear moment evaluated on fixed vectors converges at every positive Gaussian parameter.
\end{lemma}

\begin{lemma}[Taylor coefficients of an exactly homogeneous moment]
\label{lem:exact-homogeneous-moment-coefficients}
\lean{Preiss.gaussianMomentCoefficient_eq_of_exact_homogeneity}
\uses{lem:exact-monomial-coefficients}

If a positive-degree Gaussian moment is exactly $s^d A$, its Taylor coefficient family has only one nonzero positive-order coefficient, equal to $d!A$ in degree $d$.
\end{lemma}

\begin{lemma}[Sparse moment coefficients of conical measures]
\label{lem:conical-moment-coefficients}
\lean{Preiss.gaussianMomentCoefficient_conical}
\uses{lem:exact-homogeneous-moment-coefficients, lem:conical-moments}

For a positive-dimensional conical unit-uniform measure, $B_k^j$ vanishes unless $k=2j$; in that degree it equals $j!$ times the moment at parameter one.
\end{lemma}

\begin{lemma}[Even Gaussian moments at support points]
\label{lem:conical-even-support-moments}
\lean{Preiss.conical_even_moment_on_support}
\uses{lem:conical-moment-coefficients, lem:moment-support-polynomial, lem:taylor-two}

For a positive-dimensional conical unit-uniform measure and $x$ in its support, the degree $2j$ diagonal moment at parameter one equals $\lVert x\rVert^{2j}/j!$ for every $j>0$.
\end{lemma}

\begin{lemma}[A symmetric multilinear form is determined by its diagonal]
\label{lem:symmetric-multilinear-diagonal-injective}
\lean{Preiss.symmetricMultilinear_eq_zero_of_diagonal}

A continuous symmetric multilinear map over real normed spaces is zero if all its diagonal evaluations vanish. This includes degree zero and requires neither finite dimension nor completeness.
\end{lemma}

\begin{lemma}[Weighted plane energy in orthonormal coordinates]
\label{lem:weighted-plane-energy-decomposition}
\lean{Preiss.integral_weighted_planeEnergy_eq_sub_sum}
\uses{lem:orthonormal-plane-distance}

An integrable weighted squared distance to a subspace equals the weighted squared norm integral minus the sum of directional second moments.
\end{lemma}

\begin{lemma}[Leading covariance eigenspaces minimize plane energy]
\label{lem:covariance-minimizing-plane}
\lean{Preiss.exists_minimizing_plane_of_ordered_covariance}
\uses{lem:weighted-plane-energy-decomposition, lem:orthonormal-spectral-sum-bound}

Given an ordered spectral representation of the weighted covariance, the plane spanned by its first $m$ eigenvectors minimizes weighted squared distance among all $m$-planes.
\end{lemma}

\begin{lemma}[Gaussian bilinear covariance identities]
\label{lem:gaussian-bilinear-covariance}
\lean{Preiss.momentBilinear_diag, Preiss.momentBilinear_gaussianMomentForm_symmetric, Preiss.integrable_gaussian_inner_sq_of_integrable_momentKernel, Preiss.integral_gaussian_inner_sq_eq_momentBilinear}
\uses{def:bilinear-trace, def:moments}

The degree-two Gaussian moment defines a symmetric bilinear form. Its diagonal equals twice the Gaussian directional second moment divided by Gaussian mass, with directional integrability following from moment-kernel integrability.
\end{lemma}

\begin{lemma}[Normalized flat measures are conical]
\label{lem:unit-flat-conical}
\lean{Preiss.isConical_unitFlat}
\uses{lem:flat-plane-blowup}

The normalized Hausdorff measure on a linear plane is invariant under dimensional dilations.
\end{lemma}

\begin{lemma}[Zero-dimensional flat measures and their moments]
\label{lem:zero-dimensional-flat-moments}
\lean{Preiss.unitFlat_eq_dirac_of_dim_zero, Preiss.gaussianMomentForm_dirac_zero_diagonal, Preiss.flat_gaussianMoment_diagonal_dim_zero}
\uses{lem:plane-hausdorff}

A normalized zero-dimensional flat measure is the Dirac mass at the origin. Its zeroth moment is one and every positive-degree moment vanishes, including the required projected even-moment formula.
\end{lemma}

\begin{lemma}[Flat moments depend only on projected vectors]
\label{lem:flat-moment-projection}
\lean{Preiss.gaussianMomentForm_unitFlat_diagonal_project, Preiss.gaussianMomentForm_unitFlat_zero_diagonal}
\uses{lem:moment-form-evaluation, lem:gaussian-polynomial-integrable}

Diagonal Gaussian moments of a normalized flat measure are unchanged by orthogonal projection of their argument to the plane. The zeroth moment equals one at every positive Gaussian parameter.
\end{lemma}

\begin{lemma}[Unit-uniform measures are nonzero]
\label{lem:unit-uniform-nonzero}
\lean{Preiss.IsUnitUniform.ne_zero}

A unit-uniform measure is nonzero because its support contains the origin.
\end{lemma}

\begin{lemma}[Existence and uniformity of tangents at infinity]
\label{lem:unit-uniform-infinity-existence}
\lean{Preiss.isUnitUniform_of_tangentAtInfinity, Preiss.exists_unitUniform_tangentAtInfinity}
\uses{lem:unit-uniform-nonzero, lem:uniform-vague-closed, lem:uniform-compactness, lem:uniform-scaling}

A positive-dimensional unit-uniform measure admits a tangent at infinity, and every such tangent remains unit-uniform.
\end{lemma}

\begin{lemma}[Uniqueness of tangents at infinity]
\label{lem:unit-uniform-infinity-unique}
\lean{Preiss.tangentAtInfinity_eq_of_unitUniform}
\uses{lem:unit-uniform-infinity-existence, lem:blowdown-moment-limits, lem:vague-gaussian-polynomials, lem:moment-determination}

Any two tangents at infinity of a positive-dimensional unit-uniform measure agree, by equality of their Gaussian polynomial moments.
\end{lemma}

\begin{lemma}[Leading homogeneous terms along escaping sequences]
\label{lem:leading-homogeneous-term}
\lean{Preiss.tendsto_scaled_homogeneous_sum, Preiss.leading_diagonal_nonneg_of_rescaled_points}

For a finite sum of continuous multilinear homogeneous polynomials, dividing its values at $r_i y_i$ by the highest power of $r_i$ yields the leading diagonal value when $r_i\to\infty$ and $y_i\to v$. Nonnegativity on those points implies nonnegativity of the leading value.
\end{lemma}

\begin{lemma}[Conical Gaussian energy and unit-ball energy]
\label{lem:conical-gaussian-ball-energy}
\lean{Preiss.conical_gaussian_planeEnergy_eq_ball}
\uses{lem:homogeneous-energy-radial}

For a conical unit-uniform measure, Gaussian squared-distance energy equals unit-ball squared-distance energy times the corresponding Gamma factor.
\end{lemma}

\begin{lemma}[Gaussian norm moments bounded by half-scale mass]
\label{lem:gaussian-norm-half-mass}
\lean{Preiss.integral_norm_pow_mul_gaussian_le_halfMass}

An integrable Gaussian norm moment of positive degree $k$ is bounded by $(k/s)^{k/2}$ times the Gaussian mass at parameter $s/2$, for any measure on a real normed additive group.
\end{lemma}

\begin{lemma}[Normalized Gaussian norm bounds for uniform measures]
\label{lem:uniform-normalized-gaussian-norm}
\lean{Preiss.gaussianMass_half_unitUniform, Preiss.integral_norm_pow_mul_gaussian_le_unitUniform, Preiss.normalized_gaussian_norm_pow_le_unitUniform}
\uses{lem:gaussian-norm-half-mass, lem:moment-norm-bound}

For positive-dimensional unit-uniform measures, halving the Gaussian parameter scales its mass by $2^{m/2}$. Dividing Gaussian norm moments by Gaussian mass gives an explicit bound of order $s^{-k/2}$.
\end{lemma}

\begin{lemma}[Coefficient vanishing from growth of linear images]
\label{lem:moment-coefficient-growth-images}
\lean{Preiss.gaussianMomentCoefficient_map_eq_zero_of_lt_growth, Preiss.gaussianMomentCoefficient_apply_eq_zero_of_lt_growth}
\uses{lem:asymptotic-taylor-vanishing}

If a continuous linear image of a positive-degree Gaussian moment is $O(s^a)$, every positive Taylor coefficient of degree $j<a$ has zero image. Fixed-vector evaluation is a specialization.
\end{lemma}

\begin{lemma}[Continuous trace of moments on a subspace]
\label{lem:subspace-moment-trace}
\lean{Preiss.subspaceMomentTraceCLM, Preiss.subspaceMomentTraceCLM_apply, Preiss.subspaceMomentTraceCLM_momentKernel}
\uses{def:bilinear-trace, def:moments}

Tracing a degree-two moment on a subspace is a continuous linear functional. On the moment kernel it equals the squared norm of orthogonal projection.
\end{lemma}

\begin{lemma}[Subspace Gaussian trace under a cylinder identity]
\label{lem:subspace-gaussian-trace-cylinder}
\lean{Preiss.subspaceTrace_gaussianMoment_eq_integral, Preiss.subspaceTrace_gaussianMoment_of_cylinder}
\uses{lem:subspace-moment-trace}

The subspace trace of the Gaussian second moment is an explicit integral of squared projection. If that squared projection equals $2\langle b,z\rangle$ almost everywhere, the trace equals $2s$ times the first moment evaluated at $b$.
\end{lemma}

\begin{lemma}[Taylor coefficients after multiplication by scale]
\label{lem:taylor-multiply-scale}
\lean{Preiss.first_second_taylor_coefficients_of_mul_scale}
\uses{lem:asymptotic-taylor-vanishing}

If $f(s)=sa+o(s)$, $g(s)=sb+s^2c/2+o(s^2)$, and $g(s)=dsf(s)$ near zero through positive values, then $b=0$ and $c=2da$.
\end{lemma}

\begin{lemma}[Bounded holes in a plane near a uniform support]
\label{lem:bounded-tube-plane-to-support}
\lean{Preiss.exists_plane_to_support_bound_of_bounded_normal_coordinates}
\uses{lem:uniform-compactness, lem:uniform-vague-closed, lem:vague-support-approximation}

If the support of an $m$-dimensional unit-uniform measure lies in a bounded tube around an $m$-plane, every point of that plane lies within a uniform finite distance of the support.
\end{lemma}

\begin{lemma}[Taylor remainders under scalar linear maps]
\label{lem:moment-taylor-linear-image}
\lean{Preiss.gaussianMoment_taylor_remainder_map}
\uses{def:moment-coefficients}

Applying a continuous real linear functional to a Gaussian moment preserves its Taylor remainder estimate and applies the same functional to every coefficient.
\end{lemma}

\begin{lemma}[Second coefficient trace from a cylinder identity]
\label{lem:cylinder-second-coefficient-trace}
\lean{Preiss.subspaceTrace_secondCoefficient_of_cylinder}
\uses{lem:moment-taylor-linear-image, lem:subspace-gaussian-trace-cylinder, lem:taylor-multiply-scale}

For a positive-dimensional unit-uniform measure satisfying the quadratic cylinder identity almost everywhere, the subspace trace of $B_2^2$ equals $4B_1^1(b)$.
\end{lemma}

\begin{lemma}[Trace splits over orthogonal complementary subspaces]
\label{lem:bilinear-trace-orthogonal-split}
\lean{Preiss.bilinearTrace_eq_restrict_add_orthogonal}
\uses{lem:bilinear-trace-basis}

The trace of any continuous real bilinear form on a finite-dimensional inner-product space equals the sum of its traces restricted to a subspace and its orthogonal complement.
\end{lemma}

\begin{lemma}[Scalar Gaussian moments from product growth]
\label{lem:moment-product-growth}
\lean{Preiss.gaussianMomentForm_apply_isBigO_of_product_bound}
\uses{lem:uniform-normalized-gaussian-norm}

If the product of the degree-$k$ evaluation functionals is bounded almost everywhere by $C\lVert z\rVert^d$, its Gaussian moment is $O(s^{k-d/2})$. Zero degrees are included.
\end{lemma}

\begin{lemma}[Cubic coefficients with two normal inputs vanish]
\label{lem:bounded-tube-cubic-vanishing}
\lean{Preiss.thirdMomentCoefficient_two_normal_eq_zero_of_bounded_tube}
\uses{lem:moment-product-growth, lem:moment-coefficient-growth-images}

For a unit-uniform measure whose support lies in a bounded tube around a subspace, $B_3^2(v,w,w)=0$ whenever $w$ is normal to that subspace and $v$ is arbitrary.
\end{lemma}

\begin{lemma}[Binomial formula for a symmetric trilinear form]
\label{lem:symmetric-cubic-binomial}
\lean{Preiss.symmetricTrilinear_diagonal_add}

The diagonal of a continuous symmetric real trilinear form at $x+y$ expands with coefficients $1,3,3,1$ into its pure and mixed evaluations.
\end{lemma}

\begin{lemma}[Leading coefficient from a little-o remainder]
\label{lem:leading-coefficient-limit}
\lean{Preiss.tendsto_scaled_of_isLittleO_sub}

A normed-space-valued function equal to $s^j a+o(s^j)$ converges to $a$ after division by $s^j$ as $s$ tends to zero through positive values.
\end{lemma}

\begin{lemma}[Vanishing of scaled normalized bounded integrals]
\label{lem:scale-normalized-bounded-integral}
\lean{Preiss.tendsto_scale_normalized_integral_of_ae_bound}

For any nonnegative integrable weights with positive mass, the weighted average of an essentially bounded function tends to zero after multiplication by a positive scale tending to zero.
\end{lemma}

\begin{lemma}[A convergent function has no nonzero scaled derivative limit]
\label{lem:scale-derivative-limit}
\lean{Preiss.scale_derivative_limit_eq_zero}

If $f(s)$ converges to a finite real limit and $s f\prime(s)$ has a finite limit as $s\downarrow0$, that scaled derivative limit is zero. The proof uses the quotient $(sf(s))/s$.
\end{lemma}

\begin{lemma}[Rigidity forms from symmetric coefficient tensors]
\label{lem:rigidity-forms-from-symmetry}
\lean{Preiss.exists_rigidity_forms_of_symmetry}

Symmetric second and third coefficient forms determine the required normal-valued bilinear map and adjusted scalar bilinear form through Riesz representation. No measure or tangent assumption is required.
\end{lemma}

\begin{lemma}[Quadratic bound from bounded offsets]
\label{lem:quadratic-bound-bounded-offsets}
\lean{Preiss.quadratic_square_bound_of_bounded_offsets}

If $y_i\to v$, $\lVert q_i\rVert\le B$ eventually, and $A(y_i,y_i)+\langle\omega(y_i,y_i),q_i\rangle\to0$, then $A(v,v)^2\le B^2\lVert\omega(v,v)\rVert^2$. The spaces and continuous bilinear maps are arbitrary.
\end{lemma}

\begin{lemma}[Gaussian integrals of functions with linear growth]
\label{lem:gaussian-linear-growth-integral}
\lean{Preiss.gaussian_integral_bound_of_linear_growth}

A measurable function bounded almost everywhere by $C(1+\lVert z\rVert)$ is Gaussian integrable whenever Gaussian mass and first norm moment are integrable, with the corresponding explicit integral bound.
\end{lemma}

\begin{lemma}[Vanishing of scaled Gaussian linear-growth averages]
\label{lem:uniform-gaussian-linear-growth-limit}
\lean{Preiss.normalized_gaussian_integral_tendsto_zero_of_linear_growth}
\uses{lem:gaussian-linear-growth-integral, lem:uniform-normalized-gaussian-norm}

For a positive-dimensional unit-uniform measure, the Gaussian integral of a linearly bounded measurable function, multiplied by $s$ and divided by Gaussian mass, tends to zero as $s\downarrow0$.
\end{lemma}

\begin{lemma}[Scalar lower bound from a quadratic inequality]
\label{lem:quadratic-square-upper-lower}
\lean{Preiss.quadratic_lower_bound_of_square_upper}

If $d>0$, $q\ge0$, and $(1+2/d)t^2\le-8qt$, then $t\ge-8dq/(d+2)$.
\end{lemma}

\begin{lemma}[Gaussian trace lower bound from a squared-form upper bound]
\label{lem:gaussian-square-upper-trace-lower}
\lean{Preiss.gaussian_bilinear_integral_lower_bound_of_square_upper}
\uses{lem:quadratic-square-upper-lower, lem:gaussian-quadratic-trace-bound}

For a symmetric bilinear form on a positive-dimensional real inner-product space, a Gaussian squared-form bound by $-8q$ times its mean forces the mean to be at least $-8dq/(d+2)$.
\end{lemma}

\begin{lemma}[Difference of bilinear diagonal values]
\label{lem:bilinear-diagonal-difference}
\lean{Preiss.norm_bilinear_diagonal_sub_le}

For a continuous bilinear map, $\lVert Q(z,z)-Q(p,p)\rVert\le\lVert Q\rVert\lVert z-p\rVert(\lVert z\rVert+\lVert p\rVert)$.
\end{lemma}

\begin{lemma}[Uniqueness of regular vague limits]
\label{lem:vague-limit-unique}
\lean{Preiss.vague_limit_unique}

Two regular measures that are vague limits of the same sequence agree.
\end{lemma}

\begin{lemma}[Flatness transfers between tangent normalizations]
\label{lem:flat-tangent-normalization}
\lean{Preiss.exists_flat_dimensionalTangent_of_flat_tangent, Preiss.all_tangents_flat_of_dimensionalTangents, Preiss.all_dimensionalTangents_flat_of_smul}
\uses{lem:tangent-ray-density-bounds}

Under positive finite density bounds, a flat tangent yields a flat dimensional tangent, and flatness of all dimensional tangents yields flatness of all tangents. A common positive finite scalar normalization can be undone.
\end{lemma}

\begin{lemma}[Positive finite density under scalar multiplication]
\label{lem:density-positive-scalar}
\lean{Preiss.HasPositiveFiniteDensityAt.smul}
\uses{def:positive-density}

Multiplication of a measure by a positive finite scalar preserves positive finite existing dimensional density.
\end{lemma}

\begin{lemma}[Integrability of squared norms of Gaussian bilinear diagonals]
\label{lem:gaussian-bilinear-norm-integrable}
\lean{Preiss.integrable_gaussian_norm_sq_bilinear}

The squared norm of a continuous Hilbert-valued bilinear form evaluated diagonally is integrable under standard Gaussian measure on a finite-dimensional real inner-product space.
\end{lemma}

\begin{lemma}[Vector-valued Gaussian fourth moment with zero mean]
\label{lem:gaussian-vector-bilinear-wick}
\lean{Preiss.gaussian_integral_norm_sq_bilinear_of_mean_zero}
\uses{lem:gaussian-bilinear-norm-integrable, lem:gaussian-quadratic-trace-bound}

For a symmetric Hilbert-valued bilinear form with zero Gaussian mean, its Gaussian squared norm equals twice the sum of squared norms of its matrix entries in an orthonormal basis.
\end{lemma}

\begin{lemma}[Gaussian means and contractions of vector bilinear forms]
\label{lem:gaussian-vector-trace-contractions}
\lean{Preiss.gaussian_integral_bilinear_eq_basis_sum, Preiss.gaussian_integral_inner_bilinear_eq_zero_of_mean_zero, Preiss.gaussian_integral_inner_sq_mul_bilinear_of_mean_zero}
\uses{lem:gaussian-bilinear-contraction, lem:gaussian-quadratic-trace-bound}

The Gaussian mean of a vector-valued bilinear diagonal is its basis trace. Zero mean implies zero scalar component means and the corresponding factor-two weighted Gaussian contraction.
\end{lemma}

\begin{lemma}[Polarization of weighted integral limits]
\label{lem:weighted-integral-polarization}
\lean{Preiss.tendsto_weighted_integral_mul_of_square_limits, Preiss.integrable_weighted_mul_of_squares, Preiss.tendsto_weighted_bilinear_mixed_of_diagonal_limits}

Integrability and convergence of weighted diagonal square integrals imply their cross-product counterparts by polarization. This applies to diagonal limits of Hilbert-valued bilinear forms.
\end{lemma}

\begin{lemma}[Finite-basis bilinear expansion and integral limits]
\label{lem:finite-bilinear-integral-limits}
\lean{Preiss.bilinear_diagonal_eq_sum_orthonormalBasis, Preiss.tendsto_weighted_integral_finite_linear_combination}

A bilinear diagonal expands in a finite orthonormal basis. Eventual integrability and limits of finitely many weighted integrals pass through their fixed linear combinations.
\end{lemma}

\begin{lemma}[First and second scaled moment limits]
\label{lem:low-order-moment-limits}
\lean{Preiss.gaussianMoment_first_scaled_limit, Preiss.gaussianMoment_second_scaled_limit}
\uses{lem:leading-coefficient-limit, def:moment-coefficients}

For every positive moment degree, division by $s$ yields the first coefficient in the limit. If that coefficient vanishes, division by $s^2$ yields half the second coefficient.
\end{lemma}

\begin{lemma}[Centered mixed Gaussian limit from Taylor coefficients]
\label{lem:mixed-centered-gaussian-limit}
\lean{Preiss.mixed_gaussian_limit_of_moment_expansion}
\uses{lem:low-order-moment-limits, lem:moment-coefficient-vanishing}

The normalized centered cubic Gaussian integral converges to $3B_3^2(v,v,w)-4B_2^1(v,v)\langle b,w\rangle$.
\end{lemma}

\begin{lemma}[Centered Gaussian mean from its first coefficient]
\label{lem:centered-gaussian-mean-limit}
\lean{Preiss.normalized_centered_gaussian_limit_of_first_coefficient}
\uses{lem:low-order-moment-limits}

If $B_1^1(w)=2\langle b,w\rangle$, the normalized Gaussian average of $\langle z-b,w\rangle$ tends to zero.
\end{lemma}

\begin{definition}[The quadratic rigidity error]
\label{def:rigidity-error}
\lean{Preiss.rigidityError}
\uses{def:rigidity-omega, def:rigidity-adjusted-form}

The error is $A(Pz,Pz)+\langle\omega(Pz,Pz),z-Pz-b\rangle$, where $P$ is the orthogonal projection onto the limiting plane.
\end{definition}

\begin{lemma}[Linear growth of the rigidity error]
\label{lem:rigidity-error-linear-growth}
\lean{Preiss.rigidityError_le_linear_growth}
\uses{lem:rigidity-error-remainder, lem:rigidity-fourth-support-identity, lem:uniform-cylinder-identity, lem:bilinear-diagonal-difference}

On the support of the unit-uniform measure, the rigidity error is bounded in absolute value by $C(1+\lVert z\rVert)$ for some nonnegative constant $C$.
\end{lemma}

\begin{lemma}[Coefficient remainder identity for the rigidity error]
\label{lem:rigidity-error-remainder}
\lean{Preiss.rigidityError_eq_coefficient_remainder}
\uses{def:rigidity-error, def:rigidity-omega, def:rigidity-adjusted-form}

Assuming the cylinder and fourth-order coefficient identities at $z$, the rigidity error equals $\lVert z-Pz\rVert^4-B_1^2(z)-(B_2^2(z,z)-B_2^2(Pz,Pz))$.
\end{lemma}

\begin{lemma}[Pointwise quadratic bound from linear error growth]
\label{lem:rigidity-pointwise-from-growth}
\lean{Preiss.adjustedForm_sq_le_of_rigidityError_linear_growth}
\uses{def:rigidity-error, lem:quadratic-bound-bounded-offsets, lem:uniform-cylinder-identity, lem:vague-support-approximation}

For flat-at-infinity moment data, a linear growth bound on the rigidity error implies $A(v,v)^2\le\lVert b\rVert^2\lVert\omega(v,v)\rVert^2$. The proof uses support points approximating the plane along blowdowns.
\end{lemma}

\begin{lemma}[Gaussian error limit from linear growth]
\label{lem:rigidity-error-limit-from-growth}
\lean{Preiss.continuous_rigidityError, Preiss.gaussian_rigidityError_limit_of_linear_growth}
\uses{def:rigidity-error, lem:uniform-gaussian-linear-growth-limit}

The rigidity error is continuous. If it has linear growth on a unit-uniform support, its Gaussian integral multiplied by $s$ and divided by Gaussian mass tends to zero.
\end{lemma}

\begin{lemma}[Gaussian traces are unchanged by a zero-mean correction]
\label{lem:gaussian-mean-zero-correction}
\lean{Preiss.gaussian_integral_bilinear_eq_trace_of_mean_zero_correction}
\uses{lem:gaussian-vector-trace-contractions}

Adding a fixed scalar component of a vector-valued bilinear form with zero Gaussian mean leaves the Gaussian integral of a scalar bilinear diagonal unchanged.
\end{lemma}

\begin{lemma}[First-crossing sequences from a large-flatness tangent]
\label{lem:first-crossing-from-large-flatness}
\lean{Preiss.exists_first_crossing_sequences_of_flatness_gt}
\uses{lem:flatness-radius-continuous, lem:greatest-crossing, lem:dimensional-compactness}

A flat dimensional tangent and a dimensional tangent whose flatness exceeds the threshold produce the full first-crossing scale data, including every intermediate-scale bound.
\end{lemma}

\begin{lemma}[Gaussian quadratic bounds from a mean-square identity]
\label{lem:gaussian-rigidity-integral-bounds}
\lean{Preiss.gaussian_quadratic_bounds_of_mean_square_identity}
\uses{lem:gaussian-square-upper-trace-lower, lem:gaussian-bilinear-norm-integrable}

If $A(v,v)^2\le q\lVert\omega(v,v)\rVert^2$ and the Gaussian mean square of $\omega$ equals $-8$ times the mean of $A$, then the Gaussian square of $A$ is at most $-8q$ times its mean, and that mean is at least $-8dq/(d+2)$.
\end{lemma}

\begin{lemma}[Derivatives of normalized Gaussian expectations]
\label{lem:normalized-gaussian-derivative}
\lean{Preiss.hasDerivAt_normalized_polynomial_gaussian_unitUniform, Preiss.hasDerivAt_normalized_centered_gaussian_unitUniform}
\uses{lem:gaussian-parameter-derivative, lem:gaussian-radial-second-moment, lem:moment-norm-bound}

The derivative of a normalized polynomial Gaussian expectation is $m/(2s)$ times that expectation minus its normalized squared-norm-weighted expectation. The centered linear polynomial is a specialization.
\end{lemma}

\begin{lemma}[Existence of the centered full Gaussian contraction limit]
\label{lem:centered-gaussian-full-contraction-exists}
\lean{Preiss.exists_normalized_norm_sq_centered_gaussian_limit}
\uses{lem:mixed-centered-gaussian-limit}

Summing the centered mixed limits over an ambient orthonormal basis gives a finite limit for the normalized centered squared-norm Gaussian contraction.
\end{lemma}

\begin{lemma}[Vanishing of the centered full Gaussian contraction]
\label{lem:centered-gaussian-full-contraction-zero}
\lean{Preiss.normalized_norm_sq_centered_gaussian_limit_eq_zero}
\uses{lem:centered-gaussian-full-contraction-exists, lem:centered-gaussian-mean-limit, lem:normalized-gaussian-derivative, lem:scale-derivative-limit}

When the first moment coefficient represents the center, the normalized centered squared-norm Gaussian contraction tends to zero. This follows by differentiating the normalized centered mean and applying the scaled derivative limit lemma.
\end{lemma}

\begin{lemma}[Vanishing projected contraction for a bounded tube]
\label{lem:bounded-tube-projected-contraction}
\lean{Preiss.contracted_gaussian_limit_eq_zero_of_bounded_tube}
\uses{lem:centered-gaussian-full-contraction-zero, lem:scale-normalized-bounded-integral}

For a uniform support in a bounded tube around a subspace, the corresponding projected squared-norm Gaussian contraction tends to zero against every normal vector satisfying the centered first-coefficient identity.
\end{lemma}

\begin{lemma}[Projected Gaussian contraction as a finite basis sum]
\label{lem:projected-mixed-gaussian-sum}
\lean{Preiss.sum_mixed_gaussian_eq_projected}
\uses{lem:moment-norm-bound}

The projected squared-norm centered Gaussian contraction is the sum of its directional mixed contractions over an orthonormal basis of the plane.
\end{lemma}

\begin{lemma}[Gaussian contraction from weighted moment limits]
\label{lem:gaussian-weighted-contraction}
\lean{Preiss.gaussian_bilinear_contraction_of_weighted_moment_limits}
\uses{lem:finite-bilinear-integral-limits, lem:gaussian-vector-bilinear-wick}

Finite families of second and mixed weighted moment limits, together with a vanishing quadratic error and a zero Gaussian mean, imply that the Gaussian mean square of the bilinear vector form is $-8$ times the mean of the scalar quadratic form.
\end{lemma}

\begin{lemma}[Gaussian inner-product integrability and limits]
\label{lem:gaussian-inner-product-limits}
\lean{Preiss.integrable_gaussian_inner_sq_unitUniform, Preiss.integrable_gaussian_inner_sq_centered_unitUniform, Preiss.normalized_gaussian_inner_sq_limit}
\uses{lem:low-order-moment-limits, lem:moment-product-growth}

Inner-product squares and their centered linear multiples are Gaussian-integrable for a uniform measure. The normalized square integral converges to half of the first quadratic moment coefficient.
\end{lemma}

\begin{lemma}[Polarized Gaussian moment limits on the limiting plane]
\label{lem:gaussian-cross-product-limits}
\lean{Preiss.normalized_gaussian_inner_mul_limit_on_plane, Preiss.normalized_gaussian_mixed_limit_rigidityOmega}
\uses{lem:weighted-integral-polarization, lem:gaussian-inner-product-limits, lem:rigidity-mixed-gaussian-limit, lem:flat-infinity-moments}

Polarization gives the normalized second and centered mixed Gaussian limits for two arbitrary vectors in the limiting plane.
\end{lemma}

\begin{lemma}[Mean-square contraction of the rigidity form]
\label{lem:rigidity-omega-contraction}
\lean{Preiss.rigidityOmega_contraction_of_mean_zero_and_error_limit}
\uses{lem:gaussian-weighted-contraction, lem:gaussian-cross-product-limits, def:rigidity-error}

Under the zero-mean and vanishing-error identities, the Gaussian mean of $\lVert\omega(v,v)\rVert^2$ equals $-8$ times the Gaussian mean of $A(v,v)$.
\end{lemma}
